% concatenated from burnsurv.tex
\documentclass[12pt,reqno]{amsart}
\usepackage{amssymb}
%\usepackage[notref,notcite]{showkeys}
\usepackage{graphicx}
\usepackage{xcolor}
\usepackage{bbm}


\usepackage[all]{xy}
%we only want one symbol from mathb, we don't need all of mathabx
%\usepackage{mathabx}
\DeclareFontFamily{U}{mathb}{\hyphenchar\font45}
\DeclareFontShape{U}{mathb}{m}{n}{
      <5> <6> <7> <8> <9> <10> gen * mathb
      <10.95> mathb10 <12> <14.4> <17.28> <20.74> <24.88> mathb12
      }{}
\DeclareSymbolFont{mathb}{U}{mathb}{m}{n}
\DeclareMathSymbol{\righttoleftarrow}{3}{mathb}{"FD}

\oddsidemargin 15mm
\evensidemargin 15mm
\textwidth 130mm


\theoremstyle{plain}
\newtheorem{prop}{Proposition}[section]
\newtheorem{theo}[prop]{Theorem}
\newtheorem{coro}[prop]{Corollary}
\newtheorem{lemm}[prop]{Lemma}
\theoremstyle{remark}
\newtheorem{rema}[prop]{Remark}
\newtheorem{assu}{Assumption}
\theoremstyle{definition}
\newtheorem{defi}[prop]{Definition}
\newtheorem{nota}[prop]{Notation}
\newtheorem{exam}[prop]{Example}
\numberwithin{equation}{section}

\newcommand{\A}{{\mathbb A}}
\newcommand{\C}{{\mathbb C}}
\newcommand{\PP}{{\mathbb P}}
\newcommand{\bP}{{\mathbb P}}
\newcommand{\Q}{{\mathbb Q}}
\newcommand{\F}{{\mathbb F}}
\newcommand{\G}{{\mathbb G}}
\newcommand{\N}{{\mathbb N}}
\newcommand{\R}{{\mathbb R}}
\newcommand{\Z}{{\mathbb Z}}
\newcommand{\cA}{{\mathcal A}}
\newcommand{\cB}{{\mathcal B}}
\newcommand{\cC}{{\mathcal C}}
\newcommand{\cD}{{\mathcal D}}
\newcommand{\cE}{{\mathcal E}}
\newcommand{\cF}{{\mathcal F}}
\newcommand{\cG}{{\mathcal G}}
\newcommand{\cO}{{\mathcal O}}
\newcommand{\cI}{{\mathcal I}}
\newcommand{\cJ}{{\mathcal J}}
\newcommand{\cK}{{\mathcal K}}
\newcommand{\cL}{{\mathcal L}}
\newcommand{\cM}{{\mathcal M}}
\newcommand{\cN}{{\mathcal N}}
\newcommand{\cS}{{\mathcal S}}
\newcommand{\cT}{{\mathcal T}}
\newcommand{\cX}{{\mathcal X}}
\newcommand{\cY}{{\mathcal Y}}
\newcommand{\cZ}{{\mathcal Z}}
\newcommand{\cU}{{\mathcal U}}
\newcommand{\cV}{{\mathcal V}}
\newcommand{\cW}{{\mathcal W}}

\newcommand{\fA}{{\mathfrak A}}
\newcommand{\sD}{{\mathfrak D}}
\newcommand{\fK}{{\mathfrak K}}

\newcommand{\fs}{{\mathfrak s}}
\newcommand{\BC}{{\mathcal B}{\mathcal C}}

\newcommand{\rB}{{\mathrm B}}
\newcommand{\rH}{{\mathrm H}}
\newcommand{\rK}{{\mathrm K}}
\newcommand{\rN}{{\mathrm N}}

\newcommand{\iJ}{{\mathrm{IJ}}}
\newcommand{\GL}{{\mathrm{GL}}}
\newcommand{\PGL}{{\mathrm{PGL}}}


\newcommand{\ra}{\rightarrow}
\newcommand{\lra}{\longrightarrow}

\newcommand{\bA}{{\mathbb A}}
\newcommand{\bC}{{\mathbb C}}
\newcommand{\bF}{{\mathbb F}}
\newcommand{\bG}{{\mathbb G}}
\newcommand{\bL}{{\mathbb L}}
\newcommand{\bN}{{\mathbb N}}
\newcommand{\bQ}{{\mathbb Q}}
\newcommand{\bR}{{\mathbb R}}
\newcommand{\bZ}{{\mathbb Z}}
\newcommand{\bfQ}{{\mathbf Q}}
\newcommand{\bfZ}{{\mathbf Z}}


\newcommand{\Am}{\mathrm{Am}}


\newcommand{\fS}{{\mathfrak S}}
\newcommand{\mo}{{\mathfrak o}}

\newcommand{\upperhalfplane}{{\mathbb H}}
\newcommand{\eqto}{\stackrel{\lower1.5pt\hbox{$\scriptstyle\sim\,$}}\to}
\newcommand{\eqdashto}{\stackrel{\lower1.5pt\hbox{$\scriptstyle\sim\,$}}\dashrightarrow}
\newcommand{\actsfromleft}{\mathrel{\reflectbox{$\righttoleftarrow$}}}
\newcommand{\actsfromright}{\righttoleftarrow}
\DeclareMathOperator{\Alg}{Alg}
\DeclareMathOperator{\Gal}{Gal}
\DeclareMathOperator{\inv}{inv}
\DeclareMathOperator{\disc}{disc}
\DeclareMathOperator{\dv}{div}
\DeclareMathOperator{\Pic}{Pic}
\DeclareMathOperator{\Spec}{Spec}
\DeclareMathOperator{\Div}{Div}
\DeclareMathOperator{\Tor}{Tor}
\DeclareMathOperator{\Hom}{Hom}
\DeclareMathOperator{\End}{End}
\DeclareMathOperator{\Br}{Br}
\DeclareMathOperator{\Aut}{Aut}
\DeclareMathOperator{\Ker}{Ker}
\DeclareMathOperator{\Jac}{Jac}
\DeclareMathOperator{\Var}{Var}
\DeclareMathOperator{\Burn}{Burn}
\DeclareMathOperator{\Bir}{Bir}
\DeclareMathOperator{\Cr}{Cr}
\DeclareMathOperator{\Isom}{Isom}
\DeclareMathOperator{\Proj}{Proj}
\DeclareMathOperator{\Sym}{Sym}
\DeclareMathOperator{\Ind}{Ind}
\DeclareMathOperator{\im}{im}
\DeclareMathOperator{\ord}{ord}
\DeclareMathOperator{\codim}{codim}
\DeclareMathOperator{\cores}{cores}
\DeclareMathOperator{\chara}{char}
\newcommand{\ocB}{{\overline\cB}}
\newcommand{\oBurn}{{\overline\Burn}}


\begin{document}
\title[Equivariant birational geometry]{Invariants in equivariant birational geometry}



\author[Andrew Kresch]{Andrew Kresch}
\address{
  Institut f\"ur Mathematik,
  Universit\"at Z\"urich,
  Winterthurerstrasse 190,
  CH-8057 Z\"urich, Switzerland
}
\email{andrew.kresch@math.uzh.ch}


\author[Yuri Tschinkel]{Yuri Tschinkel}
\address{
  Courant Institute,
  251 Mercer Street,
  New York, NY 10012, USA
}
\email{tschinkel@cims.nyu.edu}

\address{Simons Foundation\\
160 Fifth Avenue\\
New York, NY 10010\\
USA}


\date{February 20, 2026}

\begin{abstract}
We discuss invariants in equivariant birational geometry.
\end{abstract}

\maketitle

\section{Introduction}
\label{sec.intro}
In this survey, we consider problems in classical birational geometry and their extensions to the setting of equivariant geometry. The focus is on new invariants, introduced in \cite{BnG}, building on \cite{KT}, \cite{KPT}, \cite{Bbar}, and 
\cite{HKTsmall}. 
%on their range of applicability and their limitations. 
In our applications, we work over an algebraically closed ground field $k$ of characteristic zero. 

The term {\em birational geometry} refers to the characterization of isomorphism classes of function fields of algebraic varieties over $k$, i.e., fields $K$ which are finitely generated over $k$. 
Of particular interest is the {\em  (stable) rationality problem}, i.e., 
the characterization of {\em (stably) rational} function fields. 
Recall that an algebraic variety $X$ over $k$ is called {\em rational}, respectively, {\em stably rational}, if its function field $K=k(X)$, respectively, the field $K(x_1,\ldots, x_n)$ for suitable $n$, 
is isomorphic to a purely transcendental extension of $k$. 
An interesting related notion is that of {\em unirationality}, 
i.e., the property that $K\subseteq k(x_1,\ldots, x_n)$.
The (stable) rationality problem is settled in dimension 2, but remains elusive in dimensions $\ge 3$. 

One can ask about birationality of varieties with group actions. Concretely, let $G$ be a finite group, acting on $K=k(X)$, and trivially on $k$. If the action is generically free, then $G$ is identified with a subgroup of 
$$
\Bir\Aut(X),
$$
the group of birational automorphisms of $X$ over $k$. The main question in this context is: 

\medskip

\noindent
{\em When are two such subgroups conjugate in $\Bir\Aut(X)$?}

\medskip

There is always a smooth projective model for which the action is \emph{regular} (see Section \ref{sect:basic}).
Taking $X$ to be such a model, we can view $G$ as a subgroup of 
$$
\Aut(X),
$$
the group of biregular automorphisms of $X$. 
For example, when $X=\bP^2$, the projective plane, we have $\Aut(X)\simeq \PGL_3$, and it is known classically how to distinguish finite subgroups up to conjugation in $\PGL_3$.
But the problem of distinguishing finite subgroups of $\PGL_3$ up to conjugation in the \emph{Cremona group} 
$$
\mathrm{Cr}_2:=\Bir\Aut(\bP^2),
$$ 
the group of birational automorphism of $\bP^2$, is a recent advance \cite{CTZ-p2}.


In this paper, we discuss new 
invariants in $G$-equivariant birational 
geometry, and in particular, the 
{\em equivariant Burnside group}
$$
\Burn_n(G),
$$
introduced in \cite{BnG}. 
%These invariants allow to distinguish actions that were previously indistinguishable. 
Needless to say, there 
is an enormous literature on equivariant birational geometry; we will not be able to give a complete account of all important ideas and techniques. We will point to sources that should allow an interested reader to learn more about this fascinating subject. 

One of the motivations to study $G$-birational geometry lies in strongly suggestive analogies with birational geometry over nonclosed fields, a booming subject. 
These analogies have been emphasized in \cite{manin} and \cite{manin2}:
the role of $G$ is played by Galois symmetries, which act on special loci, geometric invariants, etc. 
There are also appreciable differences between these theories. We will not pursue this theme here, but refer to \cite{HTtors}, \cite{HT-quad}, \cite{KT-dp} for recent developments inspired by this dual point of view. 

\medskip

Here is a roadmap of this paper.
In Sections \ref{sec.key} and \ref{sect:basic} we discuss the basics of equivariant birational geometry and the evolution of ideas and constructions that led us to the definition of $\Burn_n(G)$:
\begin{itemize}
\item {\em analytic tools}: motivic integration and (specialization of) rationality,
\item {\em geometric tools}: equivariant weak factorization, birational rigidity, stacks. 
\end{itemize}
We give a short summary of known obstructions in equivariant birational geometry in  Section~\ref{sect:obstr}. In Section~\ref{sect:amitsur}, we discuss the Amitsur group and introduce a higher version of the classical Amitsur invariant, based on the Leray spectral sequence for  motivic complexes of stacks.
In Section~\ref{sect:model}, we explain how to choose 
suitable birational models for $G$-actions. In Section~\ref{sect:defi}, we recall the definition of equivariant Burnside groups from \cite{BnG} and explain how to compute the class of the action. 
We comment on first structural properties of these new invariants in Section~\ref{sect:first}: 
\begin{itemize}
\item {\em filtrations},
\item {\em incompressible symbols},
\item {\em specialization}.
\end{itemize}
We present several applications in Section~\ref{sect:apply} and discuss limitations of these invariants towards the end of the paper, in Section~\ref{sect:scifi}.




\bigskip

\noindent
\textbf{Acknowledgments:}
The authors thank Joseph Ayoub for extensive discussions and essential input for the treatment of motivic cohomology.
The authors are very grateful to Brendan Hassett for his interest and collaboration on related projects. The second author was partially supported by NSF grant 2301983.



\section{Key ideas}
\label{sec.key}
Throughout, we work over a field $k$ of characteristic zero. 

\subsection*{Motivic integration}
Motivic integration was introduced by Kontsevich in his generalization of Batyrev's theorem 
\cite{batyrev}: birational Calabi-Yau varieties have equal {\em Hodge numbers}. 
%As a first approximation, one can think of motivic integration as a way of manipulating formulas \eqref{eqn:tamagawa} or \eqref{eqn:den} without reference to points over finite fields. %The theory is of course much deeper; 
It is now a major research direction, with many applications, e.g., to the study of singularities, McKay correspondence, etc.; see \cite{CLNS}. 


Let 
$$
\Var_{n,k}
$$
be the set of isomorphism classes of algebraic varieties over $k$ of dimension $n$, and put
$$
\Var_k = \bigsqcup_{n\ge 0} \Var_{n,k},
\qquad
\Var_k^{\le d} = \bigsqcup_{0\le n\le d} \Var_{n,k}.
$$
We will write $[X]$ for the class of an algebraic variety (i.e., a reduced, separated, finite-type $k$-scheme) in $\Var_k$. 
The {\em Grothendieck ring}
$$
\mathrm K_0(\Var_k)
$$
is defined as the quotient of the $\bZ$-module 
$$
\bZ[\Var_k]
$$
by the relations:
\begin{itemize}
\item {\em Excision}:  $[X]=[Z]+[X\setminus Z]$, for closed $Z\subset X$,
\item {\em Product}: $[X\times Y] = [X]\times [Y]$, for all $X$, $Y$.
\end{itemize}
%These are of course obvious from the perspective of counting points on varieties over finite fields, if we had $[X]=\#X(\bF_q)$. The advantage here, is that $k$ need not be finite. 
The ring $\mathrm K_0(\Var_k)$ carries a natural filtration, by the images of $\bZ[\Var_k^{\le d}]$. 

A {\em motivic measure} with values in a commutative ring $R$ is a ring homomorphism 
$$
\mu\colon \mathrm K_0(\Var_k) \to R.
$$
For example, for $k=\C$ we take $R=\bZ[u,v]$ and 
define, for smooth projective $X$,  
$$
\mu([X]) := \sum_{p,q} (-1)^{p+q}\,  h^{p,q}(X)u^pv^q,
$$
the {\em Hodge-Deligne polynomial}. 
%Essentially the same argument as in Section~\ref{sect:p-adic} shows that 
Kontsevich's refinement of Batyrev's theorem shows that 
birational Calabi-Yau varieties have the same class in $\mathrm K_0(\Var_k)$, thus the same Hodge-Deligne polynomials, and thus the same Hodge numbers. 


\subsection*{Stable rationality and $\mathrm K_0(\Var_k)$}
The beautiful paper \cite{larsenlunts} revealed an unexpected connection with birational geometry: smooth projective varieties $X$ and $Y$ are {\em stably birational} if and only if 
\[
[X] \equiv [Y] \pmod{[\bA^1]}, 
\]
We have a surjective homomorphism
\begin{equation}
\label{eqn.toSBir}
\mathrm K_0(\Var_k) \to \bZ[\mathrm{SBir}_k]
\end{equation}
to the free abelian group on {\em stable birationality} classes of algebraic varieties over $k$, with kernel the ideal generated by $[\bA^1]$.
In particular, every motivic measure $\mu$ which is trivial on $\bA^1$ factors through the homomorphism \eqref{eqn.toSBir}.

The proof is based on 
\begin{itemize}
\item {\em Weak factorization}: every birational map can be factored into a sequence of blow-ups and blow-downs with smooth centers,
\item {\em Bittner's presentation}: the defining 
{\em Excision} relation in  $\mathrm K_0(\Var_k)$
can be replaced by the {\em Blow-up} relation: If $\widetilde{X}$ is the blow-up of a smooth projective variety $X$ in smooth $Z\subsetneq X$ and $E$ is the exceptional divisor then  
$$
[X] - [Z]= [\widetilde{X}]- [E].
$$
\end{itemize}

The paper \cite{KT} introduced the set 
$$
\Bir_{n,k}
$$
of isomorphism classes of function fields of algebraic varieties of dimension $n$ over $k$, i.e., the set of {\em birationality} classes of $n$-dimensional algebraic varieties. The free abelian group on $\Bir_{n,k}$ is
$$
\Burn_{n,k}:=\bZ[\Bir_{n,k}].
$$
The group
$$
\Burn_k:=\bigoplus_{n\ge 0} \Burn_{n,k}
$$
has a natural ring structure, induced by the product operation on algebraic varieties. There is a surjection 
$$
\Burn_k \to \mathrm{gr}(\mathrm K_0(\Var_k)),
$$
with nontrivial kernel 
\cite[Thm.\ 2.13]{borisov}.




\subsection*{Specialization of (stable) rationality}
Applications of {\em specialization} to rationality problems go back (at least) to \cite{Beau}, which established failure of rationality of certain Fano threefolds by degenerating them to conic bundles; for  these it is easier to compute the intermediate Jacobian, and if it fails to be the Jacobian of a curve, then the same holds for the generic fiber. 

The next conceptual leap was due to Voisin \cite{voisin} and Colliot-Th\'el\`ene--Pirutka \cite{CTP} who realized that specialization holds for cycle-theoretic (stable) birational invariants such as {\em integral decomposition of the diagonal}, respectively, {\em universal $\mathrm{CH}_0$-triviality}. Again, the computation of these invariants can be easier in the special fiber. 
Immediately, this led to a wealth of new results on failure of stable rationality, see, e.g., \cite{voisin}, \cite{CTP}, 
\cite{P-cyclic}, \cite{totaro-hyper}, \cite{HKTconic}, \cite{HT-Fano}, \cite{HPT-quadric}, 
\cite{HPT-quartic}, \cite{HPT-3quad}, \cite{KT-SRBS}, \cite{Schr-2}, \cite{ABP}, \cite{Schr-small-slopes}, \cite{NO-tropical}, 
as well as the surveys \cite{peyre} and \cite{CT-schrei}.    


It turned out that both $\bZ[\mathrm{SBir}_k]$ and $\Burn_k$ admit specialization homomorphisms as well \cite{nicaiseshinder}, \cite{KT}. The key construction was the {\em motivic volume} homomorphism
$$
\mathrm K_0(\Var_K) \to 
\mathrm K_0(\Var_k),  \quad K=k((t)),
$$
respectively, the {\em Burnside volume} homomorphism
$$
\Burn_K\to \Burn_k. 
$$
A common refinement
$$
\mathrm K_0(\Var_K^{\dim}) \to 
\mathrm K_0(\Var_k^{\dim})
$$
of these two volume homomorphisms, with the
\emph{Grothendieck ring graded by dimension} $\mathrm K_0(\Var_k^{\dim})=\bigoplus_{d\ge 0}\mathrm K_0(\Var_k^{\le d})$,
was given in \cite{NO-refinement}.









\section{Equivariant birational geometry}
\label{sect:basic}
Given the developments outlined in Section \ref{sec.key}, it is natural to seek {\em equivariant} and other analogs of the main constructions. 
Here, we do this in the context of actions of a \emph{finite} group $G$; a version for varieties with logarithmic volume forms has been proposed in \cite{CKT}.

\subsection{Background}
\label{sect:basic1}
We follow the conventions of \cite[Sect.\ 2]{KT-vector}:
a $G$-variety is a reduced, separated, finite-type scheme over $k$, with regular action of $G$ that is transitive on the set of irreducible components.
The $G$-action is generically free if it is free on a nonempty invariant open subvariety.

Let $X$ and $Y$ be $G$-varieties.
We call a $G$-equivariant rational map $X\dashrightarrow Y$ a \emph{$G$-rational map}; if birational, then we call it \emph{$G$-birational}.
A $G$-rational map $X\dashrightarrow Y$ is \emph{proper} if there exist a $G$-variety $Z$ and proper $G$-equivariant morphisms $Z\to X$ and $Z\to Y$, with $Z\to X$ birational (cf.\ \cite[Sect.\ 2.12]{iitaka}, \cite[App.\ A]{hassetthyeon}).
Fundamental results in this setting are:
\begin{itemize}
\item \textbf{Equivariant desingularization},
which holds, e.g., by \cite{bierstonemilman}, and lets us replace any model by a smooth model.
\item \textbf{Equivariant weak factorization}
\cite{abramovichtemkin}: Any proper $G$-birational map $X\dashrightarrow Y$, where $X$ and $Y$ are smooth, can be factored as a composition finitely many $G$-birational maps given by blow-ups of invariant smooth $G$-subvarieties, and their inverses.
\item \textbf{Regular models}.
Any rational $G$-action is regular on some projective model, which by equivariant desingularization we may take to be nonsingular; see \cite{brionmodels}.
\end{itemize}

By convention, we work with regular $G$-actions.
We say that $X$ and $Y$ are \emph{$G$-birational}, written,
$$
X\sim_G Y
$$
when there exists a proper $G$-birational map $X\dashrightarrow Y$.
We say that two $G$-varieties $X$ and $Y$ are {\em stably $G$-birational} if 
$$
X\times \bP^n \sim_G Y\times \bP^m, \quad \text{some}\ n,m \in \bN,
$$
with trivial action on the second factors.
In particular, we do not assume that $X$ and $Y$ have the same dimension. 


A projective $G$-variety $X$ is called {\em linearizable}, respectively {\em stably linearizable}, if there exists a $G$-representation $V$ such that 
$$
X\sim_G \bP(V), \quad \text{resp.} \quad X\times \bP^n \sim_G \bP(V)
$$
for some $n$.
If $X$ is only quasi-projective, then we apply the same terminology for the respective condition, applied to an equivariant projective model.
A $G$-action on $X$ is called {\em unirational} if there exists a $G$-equivariant dominant rational map $$
\bP(V)\dashrightarrow X.
$$


By convention, we let $G$ act on $X$ on the right, and we write $X\actsfromright G$.
Correspondingly, the $G$-action on $K=k(X)$ is on the left, $G\actsfromleft K$, where we adopt the further convention that $k(X)$ is the product of the function fields of the irreducible components of $X$.


\subsection{Cremona group}
\label{sect:cremona}

Of particular interest are varieties $X$ with {\em large} birational automorphism groups
$$
\Bir\Aut(X).
$$
There are currently no general methods to determine this group, given a variety $X$. Of course, this is easy for curves, but is a formidable problem in dimensions $\ge 2$. 




There are important classes of varieties, where $\Bir\Aut(X)$ is accessible, in principle:
\begin{itemize}
\item varieties of {\em general type}, via canonical models
supplied by the Minimal Model Program \cite{HM}, 
\item K3 surfaces: every birational automorphism is a regular automorphism, and those are classified, see, e.g., \cite[Chapter 15]{Huy} or \cite{Brandhorst},
\item {\em birationally rigid varieties}, e.g., smooth quartic threefolds \cite{MI-lu}; see \cite{pukh}.
\end{itemize}


However, one of the most important examples,
$$
\Cr_n:=\Bir\Aut(\bP^n), 
$$
the {\em Cremona group} in dimension $n\ge 2$, remains mysterious.
Taking $k$ to be algebraically closed,
the classification of finite subgroups of 
$\mathrm{Cr}_2$ was a culmination of decades of efforts \cite{DI}. 
There is a classification of finite  nonabelian {\em simple} and {\em quasi-simple} subgroups of $\mathrm{Cr}_3$ \cite{ProSimple}, \cite{blancfinite}, and there are fascinating results about the overall structure of $\mathrm{Cr}_n$ \cite{pro-jordan}, \cite{blanc-quo}, but there are also many concrete unanswered questions, that fall under the general, overarching problems:
\begin{itemize}
\item When are two (necessarily linear) $G$-actions on $\bP^n$ conjugate in $\Cr_n$?
\item Is a given $G$-action on a rational variety (stably) linearizable?
\end{itemize}


\subsection{Equivariant birational rigidity}

The study of $G$-equivariant birational rigidity is a thriving area, with many beautiful constructions, see, e.g., \cite{Pro-ECM}, \cite{CS}.  
In principle, one can approach the problem of (equivariant) birationality via classification of all possible blow-ups and blow-downs. In practice, this is feasible only in small dimensions, and under the assumption that $G$ is large. 
It entails a detailed analysis of singularities, explicit invariant theory, geometric inequalities, etc. 
In dimension 2, the main technical tool is the theory of equivariant {\em Sarkisov links}; it allowed to classify finite subgroups of $\Cr_2$ in \cite{DI}. 
In dimension 3, it is the equivariant {\em Noether-Fano inequality}. Here is a sample application:   

\medskip
\noindent
{\em Example:}
The permutation $\fA_5$-action (and thus also the $\fS_5$-action) on the diagonal quadric 
$$
X=\Bigl\{ \sum_{j=1}^5 x_j^2=0 \Bigl\} \subset \bP^4
$$ 
is not linearizable \cite{CSZ}, \cite{PZ}. By \cite[Thm. 4.1]{CTZ-3}, the $\fS_5$-action is stably linearizable. 


\subsection{Stacks}
\label{sect:stacks}
Let $X$ be a $G$-variety and $[X/G]$ the associated Deligne-Mumford stack.
When $X$ is a point (with trivial $G$-action), this is the classifying stack 
$$
BG=[\Spec(k)/G].
$$
A systematic introduction to the birational geometry of Deligne-Mumford stacks can be found in \cite{KT-stacks}.

While there is loss of information in the passage
$$
X\actsfromright G \quad \rightsquigarrow  \quad [X/G] \quad \rightsquigarrow  \quad  k(X)^G,
$$
there is still a rich source of obstructions to $G$-birationality to be explored in the setting of stacks; we will see such examples below in Section~\ref{sect:amitsur}. 

An equivariant sheaf (of abelian groups) on $X$ (for the $G$-action) determines a sheaf on $[X/G]$, by which we mean a sheaf on the \'etale site of $[X/G]$ \cite[Defn.\ 4.10]{DM}.
As on any site, there are cohomology groups of sheaves of abelian groups.
For $BG$, when the base field $k$ is algebraically closed, sheaf cohomology recovers group cohomology
\begin{equation}
\label{eqn.recovergroupcoho}
\rH^p(BG,\cF)=\rH^p(G,\cF(\Spec(k))),
\end{equation}
by the \v Cech spectral sequence for the covering of $BG$ by $\Spec(k)$.
Still supposing $k$ to be algebraically closed, the Leray spectral sequence,
applied to the morphism of stacks 
$$
[X/G] \to BG,
$$
yields
\begin{equation} 
\label{eqn:spectral}
\rH^p(G,\rH^q(X, \cF)) \Rightarrow \rH^{p+q}([X/G], \cF). 
\end{equation}

One can also consider an arithmetic version, where instead of a $G$-action we consider the action of the absolute Galois group $G_k=\Gal(\bar{k}/k)$. In this case, the spectral sequence takes the form
\begin{equation} 
\label{eqn:spectral-galois}
\rH^p(G_k,\rH^q(X_{\bar{k}}, \cF)) \Rightarrow \rH^{p+q}(X_k, \cF). 
\end{equation}
This spectral sequence has been studied in great detail; see, e.g., \cite[Sect.\ 1.5]{CTSansucDuke}. 





\subsection{No-name lemma}
\label{sect:noname}

This surprisingly useful result in equivariant geometry is applicable to a $G$-vector bundle, whenever the $G$-action on the base is \emph{generically free}.
In one form, it asserts the existence of a $G$-birational map, over the base, from the total space of the vector bundle to a product with an affine space (with trivial $G$-action on the affine space); cf.\ \cite[Lemma 4.4]{CGR}.
Here we state a projective version.




\begin{lemm}[No-name lemma, projective bundle version]
\label{lemm:non-name}
Let $X$ be a $G$-variety, where
the $G$-action on $X$ is generically free.
Let $E\to X$ be a $G$-vector bundle of rank $n\ge 1$.
Then, letting $G$ act trivially on $\bP^{n-1}$ we have
$$
\bP(E) \sim_G X\times \bP^{n-1},
$$
compatibly with the projection maps to $X$.
\end{lemm}

\begin{proof}
By assumption,
the $G$-action is free on some dense invariant open subscheme $W\subset X$.
If $V\subset W$ is an affine dense open subscheme, then the intersection
\[
U=\bigcap_{g\in G} V\cdot g
\]
of the translates of $V$ is an affine dense invariant open subscheme of $X$.
Writing $U=\Spec(R)$, we have quotient $Z=\Spec(R^G)$ such that
$U\to Z$ has a structure of a $G$-torsor. Then, by \cite[Prop.\ 0.9, Ampfl.\ 1.3]{mumford-git} and standard faithfully flat descent, we may identify $G$-vector bundles on $U$
with vector bundles on $Z$.
In particular, when we shrink $U$ to a suitable invariant dense affine open and correspondingly shrink $Z$, we may suppose that the restriction of $E$ to $U$ is $G$-equivariantly trivial.
With this we obtain a proper $G$-birational map $\bP(E)\dashrightarrow X\times \bP^{n-1}$.
\end{proof}


As an immediate consequence, we have

\begin{coro}
\label{coro:no-name}
Let $G\to \GL(V^\vee)$ and $G\to \GL(W^\vee)$ be representations, that determine faithful
projective representations $G\to \PGL(V^\vee)$ and $G\to \PGL(W^\vee)$.
Then $\bP(V)$ and $\bP(W)$ are stably $G$-birational.
\end{coro}

Another useful observation concerns the canonical compactification $\bP(E\oplus 1)$ of a $G$-vector bundle $E\to X$.
The $G$-rational map $E\dashrightarrow \bP(E)$ becomes a morphism after
blowing up the zero-section $0_E$.
This extends to the compactification, and we have
an equivariant isomorphism
\[ B\ell_{0_E}\bP(E\oplus 1)\cong \bP(\cO_{\bP(E)}(-1)\oplus 1) \]
over $\bP(E)$.

\begin{coro}
\label{coro:no-name2}
Let $X$ be a $G$-variety and $E\to X$ a $G$-vector bundle of rank $\ge 1$.
If the $G$-action on $\bP(E)$ is generically free, then
\[ \bP(E\oplus 1)\sim_G \bP(E)\times \bP^1, \]
where $G$ acts trivially on $\bP^1$.
\end{coro}

\section{Obstructions in $G$-equivariant geometry}
\label{sect:obstr}
Here, we assume that $k$ (the base field, of characteristic zero) is algebraically closed.
A regular action of a group $G$ on a smooth projective $n$-dimensional variety $X$ naturally induces an action on its points $X(k)$ and on various geometric invariants associated with $X$. Among those geometric invariants are:
\begin{itemize}
\item spaces of differential forms,
\item Chow and Hilbert schemes,
\item Picard and Brauer groups,
\item \'etale cohomology and Chow groups,
\item Intermediate Jacobians, etc. 
\end{itemize}
Frequently, the information extracted from these data can be turned into a birational invariant of the $G$-action.
In this section, we recall some of these
extensively studied invariants and resulting obstructions to (stable) linearizability, and introduce new invariants.

\subsection{Dimension}
Let $X$ be a smooth projective rational variety with a generically free action of a finite group $G$. 
The simplest obstruction to linearizability 
of the action is the absence of faithful $G$-representations of dimension $\dim(X)+1$. 

\medskip
\noindent
{\em Examples:} Quadrics with actions of large groups. E.g., the smallest faithful representation of the symmetric group $\fS_n$ has dimension $n-1$. It follows that the $(n-3)$-dimensional quadric $X$, defined in $\bP^{n-1}$ by
$$
\sum_{j=1}^n x_j^2 = \sum_{j=1}^n x_j = 0  
$$
with the obvious permutation action of $\fS_n$, is not linearizable.
For $n=6$, the action is stably linearizable, by \cite[Thm.\ 4.1]{CTZ-3}.

Another example is supplied by the Mathieu group $G=\mathfrak M_{11}$. Its smallest faithful representations have dimension $10$. One of these is a self-dual representation $V$ admitting a nontrivial invariant quadratic form, thus a
nonlinearizable action on an $8$-dimensional quadric $X\subset \bP(V)$. 
By \cite[Prop.\ 5.3]{HT-quad}, stable linearizability of the $G$-action on $X$ is implied by stable linearizability of the action of the $2$-Sylow subgroup $G_2\subset G$, which in this case is 
the semi-dihedral group of order $16$. The restriction of $V$ to $G_2$ admits a decomposition $W\oplus 1\oplus 1$, from which we may conclude $X\sim_{G_2}\bP(W\oplus 1)$.
Thus the $G$-action on $X$ is stably linearizable.


Yet another example is the Pfaffian, and thus rational, cubic fourfold $X\subset \bP^5$ considered in \cite[Rmk. 15]{BBT}:
$$
\sum_{j=1}^6 x_j^3 = \sum_{j=1}^6 x_3 = 0,  
$$
with the action of the Frobienius group $\mathrm{AGL}_1(\mathbb F_7)$ of order $42$. Its smallest faithful representation has dimension $6$, and the group cannot act generically freely on $\bP^4$. 
As shown in \cite[Thm.\ 16]{BBT}, $X\times \bP^1$ is linearizable. 



\subsection{Fixed points}
\label{sect:fix}

Actions of cyclic groups on smooth projective rational varieties always have fixed points. However, this may fail for noncyclic groups (or for actions of cyclic groups on nonrational varieties). 

\medskip
\noindent
\emph{Examples:} The faithful action of the Klein four-group $\mathfrak K_4=C_2^2$ on $\bP^1$ has no fixed points.
Also, there are no fixed points for the action of the unique $C_2\oplus C_4\subset \Cr_2$ not fixing a curve of genus $\ge 1$ and not birational to an action on $\bP^2$ or $\bP^1\times \bP^1$, see \cite[Thm.\ 5]{blanc-thesis}.  

\medskip

Clearly, 
$$
X^G\neq \emptyset \quad \Leftrightarrow\quad (X\times \bP^n)^G\neq \emptyset,
$$
with $G$ acting trivially on the second factor. 


\begin{prop}[{\cite{RYessential}}]
\label{prop:point}
Let $A$ be an abelian group and 
let $X$ and $Y$ be smooth projective $A$-birational varieties. 
Then 
$$
X^A\neq \emptyset \quad\Leftrightarrow\quad Y^A\neq \emptyset.
$$
\end{prop}



\begin{exam}
Consider a blow-up of a fixed point $x\in X^A$ on a smooth projective $A$-variety. The exceptional divisor $E=\bP(\cT_{X,x})$ is the projectivization of the tangent bundle $\cT_{X,x}$, a representation of $A$. The projectivization necessarily has fixed points.  

However, for {\em nonabelian} $G$,
the existence of fixed points is {\em not} a $G$-birational invariant of smooth projective $G$-varieties.
Consider $G=\fS_3$, acting on $\bP^2$, canonical compactification of the standard $2$-dimensional representation.
When we blow up the origin, the unique fixed point, we obtain
exceptional divisor $\bP^1$, which has has no fixed points. 
\end{exam}



\medskip
\noindent
{\bf (A)}: Existence of fixed points upon restriction to abelian subgroups 
$$
A\subseteq G.
$$

Condition {\bf (A)} is a $G$-equivariant stable birational invariant of smooth projective varieties.
It holds when the $G$-action is stably linearizable, and more generally, when the $G$-action is unirational.


\subsection{Determinant}
\label{sect:det}
For the action of a finite abelian group $A$, we get a finer invariant by considering the {\em weights} of the $A$-action in the tangent space $\cT_{X,x}$ to a fixed point $x\in X^A$. Let 
$$
\beta(x):=[b_1,\ldots,b_n], \quad b_j\in A^\vee :=\Hom(A,k^\times), 
$$
be the collection of these weights, i.e., characters of $A$. Let
$$
\det(\beta(x)) \in \wedge^n(A^\vee)
$$
be the {\em determinant},
defined as the wedge product $b_1\wedge\dots\wedge b_n$ and well-defined up to sign.



\begin{prop}[{\cite{RYinvariant}}]
\label{prop:ry}
Let 
$$
\pi\colon \widetilde{X}\to X
$$ 
be an equivariant birational morphism of smooth projective $A$-varieties. 
Then for all $x\in X^A$ and
$\tilde x\in \widetilde{X}^A$ with $\pi(\tilde x)=x$ we have
$$
\det(\beta(\tilde{x})) = \pm \det(\beta(x))\in \wedge^n(A^\vee).
$$
\end{prop}





\medskip
\noindent
{\em Example:} The action of the cyclic group $C_5$ on $\bP^1$, canonical compactification of a nontrivial one-dimensional representation, falls into one $C_5$-birational equivalence class for weight $\pm 1$ and another $C_5$-birational equivalence class for weight $\pm 2$ (and we have the same description for $C_5$-isomorphism classes, because of dimension $1$).
Using Proposition \ref{prop:ry} we see that for any $n\ge 1$ and odd prime $p$ the canonical compactifications of $n$-dimensional faithful linear representations of $C_p^n$ fall into precisely $(p-1)/2$ equivariant birational equivalence classes.



\medskip
\noindent
{\bf (Det)}: Existence of $A$-fixed points, for some abelian subgroup $A\subseteq G$, with a given determinant class
$$
[\det(\beta)] \in \wedge^n(A^\vee) / \pm
$$
is a $G$-equivariant birational invariant for smooth projective $G$-varieties. 

\begin{rema}
\label{rem.Detlimitation}
The invariant {\bf (Det)} can only deliver finer information than the mere existence of $A$-fixed points in the case of $G$-varieties of dimension equal to the rank of $A$.
\end{rema}

\subsection{Cohomology}
\label{sect:coho}


One can investigate the $G$-action on the Picard group $\Pic(X)$.
A key observation is that a blow-up of a smooth $G$-stable subvariety adds a permutation module to $\Pic(X)$, and that the multiplication with projective space with trivial $G$-action adds $\Z$ to $\Pic(X)$, a trivial permutation module. It follows that $\Pic(X)$, modulo permutation modules is a stable birational invariant. This yields an invariant and obstruction to stable linearizability:


Let $X$ be a smooth projective rational variety with a generically free $G$-action. 
The similarity class 
$$
[\Pic(X)]
$$
in the set of isomorphism classes of $G$-lattices, modulo the equivalence relation given by direct sums with permutation lattices,
is a stable birational invariant.


\medskip
\noindent
{\bf (SP)}: The property that the $G$-module $\Pic(X)$ is stably permutation:
$$
[\Pic(X)]=[0].
$$

Condition {\bf (SP)} holds for stably linearizable actions.

There exist stably permutation modules that are not permutation modules; some examples are recalled in \cite[Sect.\ 1]{hoshiyamasaki}.
We are not aware of algorithms to check whether a given $G$-module is stably permutation. 
However, the next, related, invariant can be computed effectively.

\medskip
The first group cohomology
$$
\rH^1(G, \Pic(X))
$$
is a stable birational invariant.

\medskip
\noindent
{\bf (H1)}: The vanishing of $\rH^1(H, \Pic(X))$, for all subgroups $H\subseteq G$.

\medskip
If the $G$-action is stably linearizable then Condition {\bf (H1)} holds.
In some cases, the invariant $\rH^1(G,\Pic(X))$ can be determined from the stabilizer stratification of $X$ \cite{BogPro}, \cite{KT-dp}.


\medskip
\noindent
{\em Example:} Let $G=C_p$ act on a smooth projective rational surface $X$. Assume that $X^G$ contains a curve of genus $g \ge 1$; such a curve is necessarily smooth and unique. Then by \cite{BogPro}, 
$$
\rH^1(G,\Pic(X)) \cong (\bZ/p\bZ)^{2g}.
$$

\subsection{Coniveau filtrations}
\label{sect:coniveau}
Let $X$ be a smooth projective $G$-variety of dimension $n$ and $A$ a $G$-module.
As in \cite[Sect.\ 3]{KT-uni} there is the cohomology
$$
\rH^i_G(X,A):=\rH^i([X/G],A).
$$
Let us suppose that $A$ is $m$-torsion, for a positive integer $m$, and consider the subgroups:
\begin{itemize}
\item 
$\rN^1\rH^i_G(X,A)$, the group of all $\alpha \in \rH^i_G(X,A)$, such that $\alpha$ vanishes in $\rH^i_G(U,A)$ for some nonempty $G$-invariant open $U\subset X$, 
\item 
$\widetilde{\rN}^1\rH^i_G(X,A)$, consisting of classes induced via equivariant morphisms $Y\to X$, with $Y$ a smooth projective $G$-variety of dimension $n-1$.
\end{itemize}
In the definition of $\widetilde{\rN}^1\rH^i_G(X,A)$,
classes are induced via the Gysin map
$$
\rH^{i-2}([Y/G],A\otimes \mu_m^{\otimes -1})\to \rH^i([X/G],A).
$$
We have
\begin{equation} 
\label{eqn:nc}
\widetilde{\rN}^1\rH^i_G(X,A) \subseteq \rN^c\rH^i_G(X, A);
\end{equation}
see \cite{BO} for background.
Exactly as in \cite[Prop.\ 2.4]{BO} we obtain:

\begin{prop}
\label{prop:n1}
The quotient
\begin{equation}
    \label{eqn:nonv}
\rN^1\rH^i_G(X,A)/\widetilde{\rN}^1\rH^i_G(X,A) 
\end{equation}
is a stable $G$-birational invariant of $X$. 
\end{prop}

This invariant has been studied in the non-equivariant setting: the quotient \eqref{eqn:nonv} (with $G=1$) can be nontrivial, e.g., with $i=3$, $A=\bZ/2\bZ$ \cite{kameko}.

\begin{rema}
\label{rem.anotherapproach}
Benoist and Ottem \cite[Thm. 4.3]{BO} consider a variant of this construction, where for particular $G$ and $A$ there are smooth projective varieties $X$ which approximate in a certain sense the cohomology of $BG$,
for which the quotient \eqref{eqn:nonv} is nontrivial.
\end{rema}






\subsection{Intermediate Jacobians}
\label{sect:jac}
An important notion in birational geometry of threefolds over closed and nonclosed ground fields is the {\em intermediate Jacobian}. It allowed to show that a smooth cubic threefold over $\bC$ is irrational \cite{CG}, showed irrationality of many conic bundles \cite{Beau}, gave rise to examples of irrational but stably rational threefolds \cite{BCSS}. Recent developments in \cite{BW1}, \cite{HT-cycle}, \cite{HT-rat}, \cite{BW2}, \cite{KuP1}, \cite{KuP2} concerning {\em torsors} under intermediate Jacobians
led to rationality criteria for certain geometrically rational threefolds over nonclosed fields. Equivariant versions of intermediate Jacobians and cycle invariants appeared in \cite{HT-quad}. 

The following gives an analog of a birational invariant introduced in \cite[Thm.\ 3.6]{CKK}; our definition does not depend on the theory of {\em atoms} from that paper.  

\begin{theo}[{\cite[Prop.\ 5.2]{KTT}}]
\label{thm:new}
Let $X$ be a smooth projective rational threefold with a regular action of a cyclic group $G=C_p$ of prime order $p$.  
We have a decomposition of the fixed locus $X^G=\bigsqcup_{\alpha} F_{\alpha}$ into a disjoint union of smooth irreducible components. Let $C$ be a smooth projective curve of genus $g\ge 2$. 
Consider
\begin{equation} 
\label{eqn:I}
I:=-I_1-2I_2+I_3, 
\end{equation}
where 
\begin{itemize}
\item $I_1$ is the number of $F_{\alpha}$ isomorphic to $C$, 
\item $I_2$ is the number of $F_{\alpha}$ birational to $C\times \bP^1$, and 
\item $I_3$ is the number of factors of the intermediate Jacobian $\iJ(X)$ isomorphic to $\mathrm{J}(C)$, with trivial $G$-action. 
\end{itemize}
Then $I$ is an equivariant birational invariant, which vanishes when the $G$-action is linearizable. 
\end{theo}

\begin{proof}
The proof is parallel to the proof in \cite{CKK}. By equivariant weak factorization, it suffices to consider blow-ups of smooth $G$-orbits; only the following blow-ups can change the shape of the invariant: 
\begin{enumerate}
\item blow-up of a $G$-fixed curve isomorphic to $C$ not lying on a $G$-fixed surface, 
\item blow-up of a $G$-fixed curve isomorphic to $C$ on a $G$-fixed surface,
\item blow-up of a $G$-stable but not $G$-fixed curve isomorphic to $C$, 
\item blowing up a $G$-orbit of curves isomorphic to $C$. 
\end{enumerate}
Case (1) admits subcases, depending on whether or not the $G$-weights in the normal bundle of $C$ are equal. 
In the first subcase, 
$$
\Delta(I_1)=2-1,\quad \Delta(I_2)=0, \quad \Delta(I_3)=1,
$$
where $\Delta$ denotes the discrepancy.  
In the second subcase,  
$$
\Delta(I_1)=-1, \quad  \Delta(I_2)=1,\quad \Delta(I_3)=1. 
$$
In Case (2), 
$$
\Delta(I_1)=1; \quad \Delta(I_2)=0, \quad \Delta(I_3)=1.
$$
In Case (3), there are no changes in $I_j$. In Case (4), the intermediate Jacobian changes by a $G$-orbit of $\mathrm{J}(C)$; all $I_j$ remain unchanged. 

Finally, for a linear action on $\bP^3$, all terms in \eqref{eqn:I} vanish. 
\end{proof}




\section{Amitsur groups}
\label{sect:amitsur}


\subsection*{Classical Amitsur invariant}
It is well-known that the canonical line bundle on a nonsingular $G$-variety is automatically $G$-linearized. On the other hand, not every element of 
the invariant Picard group $\Pic(X)^G$ can be represented by a $G$-linearized line bundle.
These observations can be formalized, via an introduction of a cohomological invariant, the Amitsur group; 
see, e.g., \cite[App.\ A]{blancfinite}:
$$
\Am^2(X,G):=\Pic(X)^G/\{[L]\,|\,\text{$G$-linearized line bundle $L$}\}.
$$
Then, for instance,
$\Am^2(X,G)=0$ when $X=\bP(V)$ for a linear representation $V$ of $G$.
If $\Am^2(X,G)=0$, then every element of $\Pic(X)^G$ admits a lift to the group
\[ \Pic(X,G)=\Pic([X/G]) \]
of $G$-linearized line bundle $L$ on $X$.

To see that $\Am^2(X,G)$ is a stable $G$-birational invariant, 
let $X$ be a smooth projective $G$-variety and $Z\subset X$ a smooth $G$-subvariety in codimension $\ge 2$.
When we blow up $Z$ in $X$ to produce $\widetilde{X}$ we obtain $\Pic(\widetilde{X})$, differing from $\Pic(X)$ by a permutation module, generated by 
components of the exceptional divisor. The corresponding $G$-invariant class is automatically $G$-linearized. It follows that 
$$
\Am^2(\widetilde{X},G)=\Am^2(X,G). 
$$
A similar argument takes care of invariance under product with a projective space.

\subsection*{Higher Amitsur invariants}
Additional cohomological invariants
can be extracted from the Leray spectral sequence \eqref{eqn:spectral}:
\begin{equation*} 
\rH^p(G,\rH^q(X, \cF)) \Rightarrow \rH^{p+q}([X/G], \cF),
\end{equation*}
for a $G$-linearized sheaf $\cF$ on a $G$-variety $X$. 

We start by recording some basic results:
\begin{itemize} 
\item 
For $i=1$ and $\cF=\bG_m$, we have 
$$
\rH^1([X/G],\bG_m)=\Pic(X,G); 
$$
see, e.g., \cite[Sect.\ 3]{KT-dp}.
\item 
If $k$ is algebraically closed, then by \eqref{eqn.recovergroupcoho} we have
$$
\rH^i(BG, \bG_m)=\rH^i(G,k^\times),
$$
for all $i$; when $i>0$ this recovers
classical group cohomology
$\rH^i(G,k^\times)\cong \rH^{i+1}(G,\bZ)$, cf.\ \cite[\S 2.1]{KT-dp}.
\end{itemize}

Combining these results, setting $\cF=\bG_m$, and taking $X$ to be a smooth projective variety with regular $G$-action and $k$ algebraically closed, the spectral sequence \eqref{eqn:spectral} 
yields (see, e.g., \cite[Sect.\ 3]{KT-dp}):
\begin{align}
\begin{split}
\label{eqn.BrXmodG}
\qquad 0&\to  \rH^1(G,k^\times)\to \Pic(X,G)\to 
\Pic(X)^G \stackrel{\delta_2}{\lra}  \rH^2(G,k^\times)\\
&\stackrel{}{\lra} \ker(\Br([X/G])\to \Br(X))\to\rH^1(G,\Pic(X))\stackrel{\delta_3}{\lra}  \rH^3(G,k^\times), 
\end{split}
\end{align}
where 
$$
\Br(-) :=\rH^2(-, \bG_m)_{\mathrm{tors}}
$$
is the {\em Brauer group}; since $X$ is smooth, the full group $\rH^2(X,\bG_m)$ is torsion and thus is $\Br(X)$; by \cite[Prop.\ 2.5 (iii)]{ant} the same holds for $[X/G]$.
The obstruction to linearizing line bundles with classes in $\Pic(X)^G$ is cohomological:
$$
\Am^2(X,G)=\mathrm{Im}(\delta_2). 
$$
We can also consider 
$$
\Am^3(X,G):=\mathrm{Im}(\delta_3)\subset \rH^3(G,k^\times); 
$$
this is also a stable birational invariant,
by the stable birational invariance of $\rH^1(G,\Pic(X))$ and
functoriality of the spectral sequence
\cite[Sect.\ 3]{KT-dp}. 
In each of the following cases we have the vanishing of $\Am^2(X,G)$ and $\Am^3(X,G)$:
\begin{itemize}
\item When $X$ has a $G$-fixed point \cite[Sect.\ 3]{KT-dp}.
\item When the $G$-action is stably linearizable, or more generally, when the $G$-action is unirational \cite[Sect.\ 2]{KT-uni}.
\end{itemize}

All possibilities for $\Am^2$ for rational surfaces have been determined in \cite[Prop. 6.7]{blancfinite}; for $\Am^3$ this 
has been done in \cite{TZ-uni}.

\subsection*{Universal torsor obstruction}
One can extend the analysis of the spectral sequence further, using different coefficients.
We suppose that $k$ is algebraically closed.
We can consider
$$
\cF=T_{\mathrm{NS}},
$$
the N\'eron-Severi torus, viewed as a sheaf. 
For $X$ smooth, projective, and rational, 
this is the torus with character group $\mathrm{NS}(X)=\Pic(X).$
Then, the low-degree sequence reads:
$$
0\to  \rH^1(G,T_{\mathrm{NS}}(k))\to \rH^1([X/G],T_{\mathrm{NS}})\to 
\rH^1(X,T_{\mathrm{NS}})^G \stackrel{\delta_2}{\lra}  \rH^2(G,T_{\mathrm{NS}}(k))
$$
We can identify
\begin{equation} \label{eqn:ident}
\rH^1(X,T_{\mathrm{NS}}) = \End(\Pic(X)). 
\end{equation}
This has a canonical, $G$-invariant, element $1_{\Pic(X)}$. 
As explained in \cite{HTtors} and \cite[Sect. 5]{KT-uni}, 
$$
\beta(X,G):=\delta_2(1_{\Pic(X)})
$$
is the obstruction to lifting the $G$-action from $X$ to a {\em universal} torsor 

\centerline{
\xymatrix{
 & \cT_X\ar[d]^{T_{\mathrm{NS}}} \\
G \times X  \ar[r] & X
}
}
%$$
%\mathcal T_X\stackrel{T_{\mathrm{NS}}}%{\xrightarrow{\hspace{1.2cm}}} X.
%$$
\noindent 
This gives rise to: 

\medskip
\noindent 
{\bf Condition (T)}:  $\beta(X,G)=0$.

\medskip
\noindent 
This condition 
is a stable birational invariant of smooth projective irreducible rational $G$-varieties. Its failure is an obstruction to unirationality of the $G$-action on $X$ \cite[Prop.\ 5.1]{KT-uni}.

By \cite[Thm.\ 6.1]{KT-uni}, for toric varieties, where the action is via automorphisms of the torus (as an algebraic variety), the vanishing is also sufficient for the unirationality of the action; indeed, up to strata in codimension $\ge 2$, $\cT_X$ is an affine space, and the lifted action on it is linear. A classification of such liftable actions on 3-dimensional toric varieties can be found in \cite{TZ-toric}. 

The formalism of torsors for varieties over $BG$ 
in \cite{HTtors} was 
inspired by the corresponding theory of torsors over nonclosed fields, developed by Colliot-Th\'el\`ene--Sansuc, and others. Some, but not all, features of the theory carry over. The most important difference is the absence of an analog of Hilbert's Theorem 90, used extensively in \cite{CTSansucDuke}. On the other hand, over $BG$ there are additional geometric applications, including
a wealth of stably linearizable, non-linearizable actions (via birational rigidity or Burnside invariants).

The motivation for the introduction and detailed study of universal torsors is rooted in the expectation that the birational geometry of $\cT_X$ is in some sense ``easier'' than that of the base variety $X$. 
One instance of this is the following analog of 
\cite[Thm.\ 2.1.2]{CTSansucDuke}.

\begin{prop}
\label{prop:vanishbr}
Let $X$ be a smooth projective rational variety with a regular action of a finite group $G$. Assume that $\beta(X,G)=0$, i.e., the $G$-action lifts to a universal torsor $\cT_X\to X$. 
Let $\overline{\cT}_X$ be a smooth projective $G$-equivariant compactification of $\cT_X$. Then
$\Pic(\overline{\cT}_X)$ 
is a permutation module and hence
$$
\rH^1(G,\Pic(\overline{\cT}_X))=0. 
$$
\end{prop}

\begin{proof}
We follow the proof of \cite[Thm.\ 2.1.2]{CTSansucDuke}.
Since, non-equivariantly, $\cT_X$ is the
complement of the zero-sections in a sum of line bundles,
whose classes form a basis of $\Pic(X)$,
we have $k[\cT_X]^\times =k^\times$.
Thus $\Pic(\overline{\cT}_X)$ is freely generated by the components of the boundary $\overline{\cT}_X\setminus \cT_X$.
The $G$-action is a permutation action.
\end{proof}


\begin{exam}
A natural geometric construction of new actions goes as follows: given a regular $G$-action on a smooth projective rational $X$,  we can consider the action of $\widetilde{G}:=G^n\wr \fS_n$  on $\widetilde{X}:=X^n$, combining the given action with the permutation action. Then 
$$
\beta(X,G)=0 \quad \Rightarrow \quad  \beta(\widetilde{X}, \widetilde{G})=0. 
$$
Indeed, we have 
$$
\Pic(\widetilde{X})=\Pic(X)^n; 
$$
Putting $\cT_{\widetilde{X}}:=(\cT_X)^n$, 
where $\cT_X$ be a universal torsor of $X$,
we obtain a universal torsor for $\widetilde{X}$. 
The $G$-action lifts to $\cT_X$, by assumption. Then 
$\widetilde{G}$-action lifts to $\cT_{\widetilde{X}}$. 
\end{exam}





\subsection*{Condition {\bf (T)}, torsors, and Amitsur invariants}
We continue to suppose $k$ algebraically closed and $X$ smooth, projective, and rational.
We connect the $\beta(X,G)$-obstruction with the problem of lifting $G$-actions to general torsors and the Amitsur group $\Am^2(X,G)$.

Consider a $G$-torus $S$, i.e., non-equivariantly $S$ is isomorphic to $\G_m^n$ for some $n\in \bN$, and the $G$-action is given by a representation $G\to \GL(M)$, where $M$ denotes the character lattice of $S$.
We fix a $G$-equivariant homomorphism
\[ \lambda\colon M\to \Pic(X). \]
For this torus as well there is the low-degree sequence of the spectral sequence \eqref{eqn:spectral}:
\[
0\to  \rH^1(G,S(k))\to \rH^1([X/G],S)\to 
\Hom_G(M,\Pic(X)) \to  \rH^2(G,S(k)),
\]
where we use the equivariant identification
$$
\rH^1(X,S)=\Hom(M,\Pic(X)).
$$
We interpret $\rH^1([X/G],S)$ as isomorphism classes of $G$-equivariant $S$-torsors on $X$, and for such, mapping to $\lambda\in \Hom_G(M,\Pic(X))$, we call $\lambda$ the \emph{equivariant type} of the corresponding $G$-equivariant $S$-torsor.

The low-degree sequence is related to  the sequence \eqref{eqn:spectral} by functoriality, with $1_{\Pic(X)}\in \End(\Pic(X))$
mapping to $\lambda$.
Consequently, the image of $\beta(X,G)$ in $\rH^2(G,S(k))$ is the obstruction to realizing $\lambda$ as the $G$-equivariant type of a $G$-equivariant $S$-torsor on $X$.

The case $M=\bZ$, with trivial action, corresponds to a $G$-invariant class $[L]\in \Pic(X)^G$.
Here, $\rH^2(G,S(k))=\rH^2(G,k^\times)$, and the functoriality relation supplies the commutative diagram
\[
\xymatrix{
\End_G(\Pic(X))\ar[d]\ar[r]^{\delta_2} & \rH^2(G,T_{\mathrm{NS}}(k)) \ar[d] \\
\Pic(X)^G \ar[r]^{\delta_2} & \rH^2(G,k^\times) 
}
\]
where the vertical maps are given by evaluation on $[L]$.
Consequently, $$
\delta_2([L])=\mathrm{ev}_{[L]}(\beta(X,G)),
$$
and
\[ \mathrm{Am}^2(X,G)=\{\mathrm{ev}_{[L]}(\beta(X,G))\,\, |\,\, [L]\in \Pic(X)^G\}. \]

\begin{prop}
Let $X$ be a smooth projective rational variety with regular $G$-action.
If $\beta(X,G)=0$, then
every $G$-equivariant homomorphism $\lambda\colon M\to \Pic(X)$ arises as the equivariant type of some $G$-equivariant $S$-torsor on $X$.
In particular,
$$
\beta(X,G)=0\quad \Rightarrow \quad \mathrm{Am}^2(X,G)=0.
$$
\end{prop}

Examples with \emph{nonvanishing} $\beta(X,G)$ and vanishing $\mathrm{Am}^2(X,G)$ appear in \cite{KT-uni} (one with vanishing and one with nonvanishing $\mathrm{Am}^3(X,G)$).



\subsection*{Motivic cohomology}
Let $q$ be an integer.
The $q$th \emph{motivic complex} is the bounded above cochain complex of \'etale sheaves with transfers
$$
\bZ(q):= \begin{cases}
\mathrm{Sing}^{\bA^1}(\bZ_{tr}(\G_m^{\wedge q}))[-q], & \text{if $q\ge 0$},\\
0, & \text{if $q<0$},
\end{cases}
%\end{cases}, \quad q\in \bZ,
$$
%be the motivic complex, 
on smooth separated finite-type $k$-schemes;
%is $\bfZ(q):=C_*\bZ_{tr}(\G_m^{\wedge q})[-q]$ for $q\ge 0$, and by convention is $0$ if $q<0$; 
see, e.g., \cite[Defn.\ 3.1]{weibel}.
Here, 
\begin{itemize} 
\item $\mathrm{Sing}^{\bA^1}$ denotes the Suslin-Voevodsky construction \cite[Defn.\ 2.14]{weibel}, which associates to a given presheaf with transfers the complex, formed by its values on products with algebraic simplices,
\item 
$\bZ_{tr}(X)$ is the presheaf with transfers sending $U$ to the free abelian group on the set of correspondences (finite over $U$) to $X$, and
\item $\G_m$ denotes the pointed scheme $(\A^1\setminus \{0\},1)$, with $q$th iterated smash product $\G_m^{\wedge q}$.
\end{itemize}

For $q=0$, we have $\bZ(0)=\mathrm{Sing}^{\bA^1}(\bZ)$, quasi-isomorphic to $\bZ$.
For $q=1$, we have $\bZ(1)[1]=\mathrm{Sing}^{\bA^1}(\bZ_{tr}(\G_m))$, quasi-isomorphic to the \'etale sheaf $\mathbb G_m$.

We are interested in \'etale motivic cohomology groups
\begin{equation}
\label{eqn.etalemotiviccohomology}
\rH^p(X,\bZ(q)):=\mathbb{H}^p_{\acute et}(X,\bZ(q)|_{X_{\acute et}}).
\end{equation}
For instance, when $k$ is algebraically closed and $X$ is a smooth projective variety, we have
\begin{align*} 
\rH^0(X,\bZ(0))& =\bZ,& \rH^0(X,\bZ(1))& = 0,\\
\rH^1(X,\bZ(0))& = 0,& \rH^1(X,\bZ(1))&  = k^\times,  \\
\rH^2(X,\bZ(0))& = \Hom(\pi_1(X),\bQ/\bZ)), & \rH^2(X,\bZ(1))& = \Pic(X), \\
&&\rH^3(X,\bZ(1))& = \Br(X).
\end{align*}
(Normal schemes have vanishing $\rH^1({-},\bZ)$, with $\rH^i({-},\bQ)$ for all $i\ge 1$ \cite[(2.1)]{deninger}, so we get the assertions in the left-hand column from the short exact sequence $0\to \bZ\to \bQ\to \bQ/\bZ\to 0$, and the right-hand column shows the
\'etale cohomology of $\G_m$ in low degrees.)

In \eqref{eqn.etalemotiviccohomology}, the hypercohomology on the right is $\rH^p(R\Gamma(X,\bZ(q)|_{X_{\acute et}}))$.
Another approach is to work in the category of \'etale motives.
This arises in a two-step construction:
\emph{localization} of the derived category $\mathbf{D}=\mathbf{D}(\mathrm{Sh}_{\acute et}(Cor_k))$ of \'etale sheaves with transfers at $\A^1$-weak equivalences to produce the category of effective \'etale motives $\mathbf{DM}_{\acute et}^{\mathrm{eff}}(k)$ \cite[Defn.\ 9.2]{weibel},
and \emph{stabilization}, which formally inverts tensoring with $\bZ(1)$ to yield the category of \'etale motives $\mathbf{DM}_{\acute et}(k)$.
The natural functor
\[ \mathbf{DM}_{\acute et}^{\mathrm{eff}}(k)\to \mathbf{DM}_{\acute et}(k) \]
is fully faithful, by
\cite[Prop.\ A.3]{huberkahn}
and \cite[Lemme 4.2]{ayoubrealisation}.



\begin{theo}
\label{thm.etalemotiviccohoDM}
Let $k$ be a field of characteristic $0$ and $X$ a smooth separated $k$-scheme of finite type.
Then
\[
\mathbb{H}^p_{\acute et}(X,\bZ(q)|_{X_{\acute et}}) \cong \Hom_{\mathbf{DM}_{\acute et}(k)}(\bZ_{tr}(X),\bZ(q)[p]).
\]
\end{theo}

For the convenience of the reader we record a proof of this standard result; cf.\ \cite[Thm.\ 4.12]{ayoub}.
As described in \cite{spaltenstein}, for unbounded complexes, notions such as $\mathrm{K}$-injective are relevant.
For example, among unbounded complexes the injective objects are those that not only are exact (by factoring $1_{A^\bullet}$
through $\mathrm{cone}(1_{A^\bullet})$) and termwise injective (exact left adjoint to the $n$th term functor), but are also $\mathrm{K}$-injective; cf.\ \cite[Exa.\ 3.2]{hoveyreptheory}.
Over a Grothendieck category
there is the \emph{injective model structure} on cochain complexes
(loc.\ cit., see also \cite{beke}):
fibrations are epimorphisms with termwise injective $\mathrm{K}$-injective kernel,
cofibrations are monomorphisms,
and weak equivalences are quasi-isomorphisms.
In particular, for complexes of sheaves of abelian groups on an essentially small site,
$R\Gamma$ is obtained by applying $\Gamma$ to
a fibrant replacement, i.e., a quasi-isomorphic
termwise injective $\mathrm{K}$-injective complex.

\begin{proof}
We use the identification of \'etale hypercohomology with hyperext:
\begin{equation}
\label{eqn.hyperext}
\mathbb{H}^p_{\acute et}(X,\cL^\bullet|_{X_{\acute et}})\cong \Hom_{\mathbf{D}}(\bZ_{tr}(X),\cL^\bullet[p]),
\end{equation}
for any complex $\cL^\bullet$ of \'etale sheaves with transfers.
This fact \cite[Exer.\ 6.25]{weibel}
is conveniently addressed using
the local projective model structure on complexes of presheaves on $Cor_k$, respectively $Sm_k$ (smooth separated finite-type $k$-schemes), respectively $X_{\acute et}$, described in \cite[\S 5.2]{choudhurygallauer}, and the
corresponding
right-lifted model structures
on respective complexes of sheaves
(cf.\ \cite[Cor.\ 2.7]{GKR}).
In each case, the fibrant objects are those that satisfy \'etale hyperdescent;
exact forgetful functors, that forget the transfers and restrict to $X_{\acute et}$,
send fibrant objects to fibrant objects.
It suffices to show that the object $\bZ(q)$ of $\mathbf{D}$ is $\A^1$-local, since by \cite[Lemma 9.19]{weibel} the $\Hom$ on the right in the statement of the theorem is then identified with
$\Hom_{\mathbf{D}}(\Z_{tr}(X),\bZ(q)[p])$,
and we conclude by \eqref{eqn.hyperext}.

We show, more generally, that for any \'etale sheaf with transfers $\cK$, the Suslin-Voevodsky construction
$\mathrm{Sing}^{\bA^1}(\cK)$ gives the $\A^1$-localization of $\cK$, meaning:
\begin{itemize}
\item the natural map $\cK\to \mathrm{Sing}^{\bA^1}(\cK)$ is an $\A^1$-weak equivalence,
\item $\mathrm{Sing}^{\bA^1}(\cK)$ is an $\A^1$-local complex of \'etale sheaves with transfers.
\end{itemize}
The first assertion holds by \cite[Lemma 9.15]{weibel}.
The second is more subtle and relies on $\mathrm{char}(k)=0$, but quickly reduces to the case that $k$ is algebraically closed.
Indeed, writing $\bar k$ for an algebraic closure and
\[ p\colon \A^1_k\to \Spec(k),\quad \bar p\colon \A^1_{\bar k}\to \Spec(\bar k),\quad
\alpha\colon \Spec(\bar k)\to \Spec(k), \]
we have generally for a complex $\cL^\bullet$ over $\mathrm{Sh}_{\acute et}(Cor_k)$, that
$\cL^\bullet$ is $\A^1$-local if and only if the natural map $\cL^\bullet\to Rp_*p^*\cL^\bullet$ is a quasi-isomorphism.
Being a quasi-isomorphism may be detected after applying $\alpha^*$.
By \cite[Lemme 4.2]{ayoubrealisation} we get an identification of $\alpha^*Rp_*p^*\cL^\bullet$ with $R\bar p_*\bar p^*\alpha^*\cL^\bullet$.
Now the reduction is clear, since essentially by definition, we have an identification
$\alpha^*\mathrm{Sing}^{\bA^1}(\cK)\cong \mathrm{Sing}^{\bA^1}(\alpha^*\cK)$.

For the rest of the proof, we suppose $k$ is algebraically closed.
Then we may apply \cite[Prop.\ 9.30]{weibel}, which tells us that a complex over $\mathrm{Sh}_{\acute et}(Cor_k)$, which has strictly $\A^1$-homotopy invariant cohomology sheaves, is $\A^1$-local.
For $n\in \bZ$,
the cohomology presheaf $\rH^n(\mathrm{Sing}^{\bA^1}(\cK))$ is
$\A^1$-homotopy invariant by \cite[Cor.\ 2.19]{weibel}.
This leads to strict $\A^1$-homotopy invariance of the sheafification $\cF:=a_{\acute et}\rH^n(\mathrm{Sing}^{\bA^1}(\cK))$, i.e., the property that
\[ R\Gamma(U,\cF)\to R\Gamma(U\times \A^1,\cF) \]
is a quasi-isomorphism
for $U$ in $Sm_k$.
Indeed, this holds if and only if the same is true with $\cF$ replaced by $\cF\otimes^{\mathbf{L}}\bZ/\ell\bZ$ for $\ell$ prime and by $\cF\otimes \bQ$, since cohomology commutes with $\otimes^{\mathbf{L}}\bZ/\ell\bZ$, and with $\otimes\bQ$, the latter an instance of commuting with filtered colimits \cite[Rmk.\ III.3.6]{milne}.
The case of $\cF\otimes^{\mathbf{L}}\bZ/\ell\bZ$ is taken care of by the Suslin rigidity theorem \cite[Thm.\ 7.20]{weibel}, and $\cF\otimes \bQ$, by \cite[Cor.\ 14.22]{weibel} and \cite[Thm.\ 5.6, Prop.\ 5.27]{voevodcohomological}.
\end{proof}

\begin{rema}
\label{rem.SuslinVoevodskyofcomplex}
By the same argument, for an arbitrary complex $\cK^\bullet$ of \'etale sheaves with transfers,
$\mathrm{Sing}^{\bA^1}(\cK^\bullet)$ gives the $\A^1$-localization of $\cK^\bullet$ \cite[Cor.\ 4.11]{ayoub}.
\end{rema}

By \cite[Prop.\ 3.3.1]{voevodtriangulated}, we have the exact
sheafification functor 
$$
\mathrm{Sh}_{Nis}(Cor_k)\to \mathrm{Sh}_{\acute et}(Cor_k),
$$
thus a triangulated functor
\[ \mathbf{DM}_{Nis}^{\mathrm{eff}}(k)\to \mathbf{DM}_{\acute et}^{\mathrm{eff}}(k). \]
We have the projective bundle and blow-up formulas in $\mathbf{DM}_{Nis}^{\mathrm{eff}}(k)$ \cite[Thm.\ 15.12, Cor.\ 15.13]{weibel}:
\begin{gather*}
\bZ_{tr}(\bP(E))\cong \bigoplus_{r=0}^{d-1} \bZ_{tr}(X)(r)[2r]\qquad (\mathrm{rk}(E)=d), \\
\bZ_{tr}(B\ell_Z(X))\cong \bZ_{tr}(X)\oplus\bigoplus_{r=1}^{d-1} \bZ_{tr}(Z)(r)[2r]\qquad (\mathrm{codim}(Z)=d).
\end{gather*}
So the same formulas are valid in $\mathbf{DM}_{\acute et}^{\mathrm{eff}}(k)$ and we conclude:

\begin{prop}
\label{prop.blowup}
For a vector bundle $E\to X$ of rank $d$
on a smooth separated finite-type scheme $X$ over $k$,
respectively, a smooth subscheme $Z\subset X$ of codimension $d$,  we have
\begin{gather*}
\rH^p(\bP(E),\bZ(q)) \cong \bigoplus_{r=0}^{d-1}\rH^{p-2r}(\bP(E),\bZ(q-r)),\\
\rH^p(B\ell_Z(X),\bZ(q)) \cong \rH^p(X,\bZ(q))\oplus \bigoplus_{r=1}^{d-1} \rH^{p-2r}(Z,\bZ(q-r)).
\end{gather*}
\end{prop}

\begin{proof}
We apply Theorem \ref{thm.etalemotiviccohoDM} and the projective bundle and blow-up formulas in $\mathbf{DM}_{\acute et}^{\mathrm{eff}}(k)$.
\end{proof}


\subsection*{Applications of motivic cohomology}
We suppose that $k$ is algebraically closed and $X$ is a smooth projective rational variety.
This implies, in particular, that the birational invariant $\Br(X)$ of $X$ vanishes.
Let a regular action of $G$ on $X$ be given.

We use scheme approximations
in order to apply a result such as Proposition \ref{prop.blowup} to $[X/G]$.
This is the construction \cite[Sect. 6.3]{EG},
used to define equivariant (higher) Chow groups.
For each $i>0$, a representation $V$ of $G$ is chosen, with invariant open $U\subset V$ and $V\setminus U$ of codimension $\ge i$, such that a $G$-principal bundle $U\to U/G$ exists in the category of smooth separated schemes of finite type over $k$.
By \cite[Prop.\ 7.1]{mumford-git}, the same holds for
\[ X\times U\to (X\times U)/G. \]
The scheme approximation
$(X\times U)/G$ has the property, that
\[ \rH^p([X/G],\Z(q))\to \rH^p((X\times U)/G,\Z(q)) \]
is an isomorphism, provided $i\ge \min((p+1)/2,q)$,
as we see by
comparing the Leray spectral sequence \eqref{eqn:spectral}
for $X$ and $X\times U$ and applying homotopy invariance and purity
\cite[Prop.\ 3.13]{KN-I}.

In the Leray spectral sequence \eqref{eqn:spectral}, we put $\mathcal F:=\bZ(1)$ and analyze the outcome. In particular, we are interested in the kernel $\ker(\varphi)$ of the natural homomorphism
$$
\varphi\colon\rH^4 ([X/G],\bZ(1)) \to \rH^4(X,\bZ(1))^G. 
$$
Considering deeper terms in the spectral sequence, we obtain an exact sequence
\begin{equation}
\label{eqn.deeper}
\rH^3(G,k^\times)\to \ker(\varphi)\to 
\rH^2(G, \Pic(X)) \stackrel{\delta_4}{\lra} \rH^4(G,k^\times),
\end{equation}
where  $\delta_4:=d_2^{22}$ is the differential of the spectral sequence. 
This allows to define the degree $4$ Amitsur group 
$$
\Am^4(X,G):=\mathrm{Im}(\delta_4) \subseteq \rH^4(G,k^\times). 
$$

\begin{theo}
\label{thm:amit4}
Let $X$ be a smooth projective rational variety, equipped with a regular action of a finite group $G$. Then 
$\Am^4(X,G)$ satisfies the following properties:
\begin{enumerate}
\item It is a stable birational invariant of the $G$-action.
\item It vanishes if $X^G\neq \emptyset$. 
\item Given a $G$-equivariant morphism $Y\to X$ of smooth projective rational
$G$-varieties one has
$$
\Am^4(X,G)\subseteq \Am^4(Y,G). 
$$
If this induces an isomorphism $\Pic(X)\cong\Pic(Y)$ then  
$$
\Am^4(X,G)=\Am^4(Y,G).
$$
\item If $V$ is a linear representation of $G$ then 
$$
\Am^4(\bP(V),G)=0. 
$$
\end{enumerate}
\end{theo}

\begin{proof}
First we address the birational invariance.    By equivariant weak factorization, we only need to consider blow-ups of smooth subvarieties.  

Let $X$ be a smooth projective variety and $Z\subset X$ a smooth subvariety of codimension $d$. Let $\widetilde{X}\to X$ be the blow-up of $X$ in $Z$. 
By Proposition \ref{prop.blowup} and the vanishing of $\bZ(q)$ when $q<0$, we have 
\begin{equation} 
\label{eqn:blow-stack}
\rH^p(\widetilde{X}, \bZ(1))=
\rH^p(X, \bZ(1)) \oplus \rH^{p-2}(Z,\bZ(0)). 
\end{equation}  
For $p=4$, the additional term takes the form
\begin{equation}
\label{eqn:z}
\rH^{2} (Z,\bZ(0))=\Hom(\pi_1(Z),\bQ/\bZ). 
\end{equation}
Same formulas hold for stacks, as can be seen via scheme approximations. 

Assume that $Z$ is a smooth $G$-subvariety of $X$, and let $Z_0$ be a component of $Z$.
We let $H$ be the subgroup of $G$ that stabilizes the component $Z_0$.
Then, $[Z/G]\cong [Z_0/H]$.
We apply the blow-up formula:
\[
\xymatrix@C=18pt{
\rH^4 ([\widetilde{X}/G],\bZ(1)) \ar[r]^{\tilde\varphi}\ar@{=}[d] & \rH^4(\widetilde{X},\bZ(1))^G\ar@{=}[d]\\
\rH^4 ([X/G],\bZ(1))  \oplus \rH^2([Z_0/H],\bZ(0)) \ar[r] & \rH^4(X,\bZ(1))^G  \oplus \rH^2(Z_0,\bZ(0))^H 
}
\]

Using the low-order terms of the Leray spectral sequence for $\cF=\bZ(0)$, we have
$$
0\to \rH^2(H,\bZ)\to \rH^2([Z_0/H],\bZ(0))\to \rH^2(Z_0,\bZ(0))^H,
$$
thus
$$
\ker(\tilde\varphi)\cong \ker(\varphi)\oplus \rH^2(H,\bZ).
$$
One the other hand,
$$
\rH^2(G,\Pic(\widetilde{X}))\cong\rH^2(G,\Pic(X))\oplus \rH^2(H,\bZ).
$$
Comparing the exact sequence \eqref{eqn.deeper} for $X$ and $\widetilde{X}$, we see that
$$
\Am^4(X,G)=\Am^4(\widetilde{X},G).
$$



We now address invariance under upon passage to the product with a projective space with trivial $G$-action. 
This uses the following formula for motivic cohomology, a trivial case of the projective bundle formula:
$$
\rH^p(X\times \bP^n, \bZ(1))=
\rH^p(X, \bZ(1)) \oplus \rH^{p-2} (X,\bZ(0)).
$$
Then we can argue as above and obtain the invariance.

Properties (2) and (3) are immediate from the functoriality of the Leray spectral sequence. Property (4) follows by considering the $G$-rational map
$$
\bP(V\oplus 1)\dashrightarrow \bP(V). 
$$
By Property (2), $\Am^4(\bP(V\oplus 1),G)=0$. 
After blowing up the origin this becomes a morphism
$$
B\ell_0\bP(V\oplus 1)\to \bP(V).
$$
By Property (1), we still have
$\Am^4(B\ell_0\bP(V\oplus 1),G)=0$; now it suffices to apply Property (3). 
\end{proof}

\begin{coro}
\label{cor:amit4}
Let $X$ be a smooth projective rational variety with regular action of a finite group $G$.
If $X$ is $G$-unirational, then $\mathrm{Am}^4(X,G)=0$.
\end{coro}



%Since $k$ is algebraically closed, $k^\times/(k^\times)^{\mathrm{tors}}$ has a structure of $\Q$-vector space and, consequently, vanishing group cohomology in positive degrees.
%So $\rH^i(G,k^\times)$ is independent of $k$ for $i>0$ and can be identified with $\rH^i(G,\Q/\Z)$.




\begin{exam}
\label{exa.K4P1}
One can compute $\Am^4(X,G)$, following \cite[p. 
~84]{KT-quotient}.
We carry this out, taking $G$ to be
the Klein $4$-group, acting on $\bP^1$ by $x\mapsto -x$ and $x\mapsto x^{-1}$.
The result is
\[ \Am^4(\bP^1,G)=\rH^4(G,k^\times)\cong (\bZ/2\bZ)^2. \]
For the computation we use the invariant set of divisors $\{0,\infty\}$, and we get
$$
\rH^2(G,\Pic(\bP^1))=\rH^2(G,\bZ)\cong (\bZ/2\bZ)^2,
$$
mapping isomorphically to
$\rH^3(G,\Z\cdot(0-\infty))\cong (\bZ/2\bZ)^2$, mapping isomorphically to $\rH^4(G,k^\times)$.
\end{exam}


%\noindent
%{\bf (Ami)}:
%The groups 
%$$
%\mathrm{Am}^d(X,G), \quad %d=2,3,4
%$$
%are stable $G$-birational invariants, which vanish when the $G$-action is unirational or admits a fixed point. 


\subsection*{Bogomolov multipliers}
Given that many cohomological invariants vanish upon existence of fixed points, and taking into account Condition {\bf (A)}, which is necessary for equivariant unirationality and (stable) linearizability of the action, it is natural to consider generalized {\em Bogomolov multipliers}
$$
\rB^j(G,M):=\Ker\left( \rH^j(G,M)\to \bigoplus_{A} 
\rH^j(A,M)\right), 
$$
where $M$ is a $G$-module and the sum is over all abelian subgroups $A\subseteq G$, see  \cite[Sect. 2]{TZ-uni}.
In search of interesting examples, we focus on $G$ and $M$ with nontrivial $\rB^j(G,M)$, for some $j$.  

We suppose that $k$ is algebraically closed.
For $j=2$ and $M=k^\times$, with trivial $G$-action, we obtain the classical Bogomolov multiplier
$$
\rB^2(G,k^\times)= \Br_{\mathrm{nr}}(k(V)^G),
$$
the unramified Brauer group of the field of invariants of a faithful representation $V$ of $G$ \cite{B-Brauer} over an algebraically closed field $k$ of characteristic zero.

We can extend \cite[Prop.\ 2]{TZ-uni} to the degree $4$ Amitsur group.

\begin{prop}
\label{prop.Am4Bogo}
Let $X$ be a smooth projective rational variety with regular action of $G$.
If $X$ satisfies Condition \emph{\textbf{(A)}}, then
\[
\mathrm{Am}^4(X,G)\subseteq \rB^4(G).
\]
\end{prop}


\section{Equivariant birational geometry -- choice of a model}
\label{sect:model}
In this section, we fix a finite group $G$ and assume that the base field $k$ contains a primitive $e$th root of unity, where $e$ is the least common multiple of the orders of the elements of $G$.

\subsection{Stabilizer stratification}
\label{sect:stab}
Let $X$ be a smooth $G$-variety, and suppose that the $G$-action is generically free.
One can 
introduce the \emph{stabilizer poset} 
\[
\mathcal{P}(X,G):=
\{ \mathrm{Stab}(\bar x)\,|\,\bar x\in X(\bar k) \}. 
\]
the set of subgroups of $G$ (ordered by inclusion) which occur as stabilizer groups of geometric points of $X$.

We are interested in the fixed locus $X^G$ of the action, i.e., the maximal closed subscheme of $X$ on which $G$ acts trivially, and more generally in the fixed loci $X^H$ for subgroups $H\subseteq G$.
These are smooth (see, e.g., \cite[Prop.\ A.8.10]{conradgabberprasad}), and may have several components, of various dimensions.
They are collected in the \emph{stabilizer stratification} \cite{BG}:
\[
\mathcal{S}(X,G):=\{\text{components of $X^H$}\,|\,H\subseteq G\}.
\]

Any $F\in \mathcal{S}(X,G)$ gives rise to two associated subgroups of $G$:
\begin{itemize}
\item The \emph{generic stabilizer group} \[ \mathrm{I}(F)\in \mathcal{P}(X,G). \]
\item The \emph{component stabilizer group}
\[ \mathrm{D}(F):=\{g\in G\,|\,F\cdot g=F\}. \]
\end{itemize}
The generic stabilizer group is a normal subgroup of the component stabilizer group:
\[ \mathrm{I}(F)\vartriangleleft \mathrm{D}(F). \]

In case $X$ is affine, $X=\Spec(A)$, with quotient variety $\Spec(B)$, $B=A^G$, and $F=\Spec(A/\mathfrak{P})$, with prime ideal $\mathfrak{P}$ over a prime ideal $\mathfrak{p}$ of $B$, these reproduce the classical inertia and decomposition groups:
\begin{itemize}
\item $\mathrm{I}(F)=\mathrm{I}(\mathfrak{P}/\mathfrak{p})$, the inertia group of $\mathfrak{P}$ over $\mathfrak{p}$.
\item $\mathrm{D}(F)=\mathrm{D}(\mathfrak{P}/\mathfrak{p})$, the decomposition group of $\mathfrak{P}$ over $\mathfrak{p}$.
\end{itemize}
The extension of residue fields
$A_{\mathfrak{P}}/\mathfrak{P}A_{\mathfrak{P}}$ of
$B_{\mathfrak{p}}/\mathfrak{p}B_{\mathfrak{p}}$ is Galois, and the Galois group is canonically identified with $\mathrm{D}(F)/\mathrm{I}(F)$.

\begin{exam}
\label{exa.notcentral}
Let the dihedral group $\mathfrak{D}_4$ of order $8$ act on affine space $\A^3=\Spec(k[x,y,z])$, where generators $\rho$ and $\sigma$ act by
\[ (x,y,z)\mapsto (-y,x,z)\qquad\text{and}\qquad (x,y,z)\mapsto (x,-y,-z), \]
respectively.
Then the subvariety $F$, defined by $x=y=0$, belongs to $\mathcal{S}(\A^3,\mathfrak{D}_8)$, with $\mathrm{I}(F)=\langle \rho\rangle$ and $\mathrm{D}(F)=\mathfrak{D}_4$.
\end{exam}





\subsection{Models}
\label{sect:standardf}
One could try to develop a theory of birational invariants of $G$-actions by analyzing the arrangement
$\cS(X,G)$
on a given model $X$. It turns out that the theory admits significant simplifications, under some assumptions on the model, which can be achieved by performing suitable equivariant blow-ups.
We comment on the nature of the simplification in Section~\ref{sect:scifi}.

Our setting is, still, a smooth $G$-variety, where the $G$-action is generically free.
One of the first simplifications is to obtain a model $X$ for which $\mathcal{P}(X,G)$ consists of {\em abelian} subgroups of $G$. Abelianization of stabilizers has been discovered and rediscovered several times;
the earliest reference we found is \cite{bogomolov}. Subsequently, the result appeared in \cite{RYessential}, \cite{Bcan}, and \cite{BG}. 
In a more modern formulation, it appears in the framework of \emph{destackification} in the theory of algebraic stacks \cite{bergh},
\cite{berghrydh}. 
The procedure, described below in Corollary \ref{cor.divisorial}, achieves several conditions:
\begin{itemize}
\item {\bf Abelian stabilizers} --
every $H\in \mathcal{P}(X,G)$ is abelian. 
\item {\bf Divisorial form} --
the action has abelian stabilizers, and
for every $H\in \mathcal{P}(X,G)$ and $F\in \mathcal{S}(X,G)$ with $\mathrm{I}(F)=H$, putting $Y:=\mathrm{D}(F)/H$, the composite
\[
\Pic(X,G)\to \rH^1(\mathrm{D}(F), k(F)^\times )
\to \rH^1(H, k(F)^\times)^Y \to H^\vee
\]
is surjective.
The leftmost map is given by restriction to the generic point of $F$,
the middle map is the restriction map of the inflation-restriction exact sequence, and the rightmost map is
$$
\rH^1(H,k(F)^\times)^Y\subseteq \rH^1(H,k(F)^\times)=\Hom(H,k(F)^\times)\cong H^\vee.
$$
\item {\bf Standard form} -- there exists $G$-invariant open $U\subset X$, with simple normal crossing boundary, such that
\begin{itemize}
\item $G$ acts freely on $U$, 
\item for every irreducible component $Z$ of the boundary $X\setminus U$, every $g\in G$ satisfies
\[
Z\cdot g \cap Z = \emptyset \qquad \text{or} \qquad Z\cdot g=Z.
\]
\end{itemize}
\end{itemize}
In \cite{BnG},
``divisorial form'' was called ``Assumption 2''; a similar divisoriality condition, with surjection from a given group of classes of orbifold line bundles of an algebraic orbifold, appeared in \cite{Bbar}.
The condition ``standard form'' was introduced in \cite{RYessential}.

\begin{exam}
\label{exa.abeliandivisorial}
If $G$ is abelian, then the
action of $G$ on $X$ is automatically in divisorial form.
Indeed, we have $G^\vee\to \Pic(X,G)$,
and the composite with the map from the definition of divisorial form is
$G^\vee\to H^\vee$ is the restriction map, Cartier dual to the inclusion $H\to G$.
This is surjective.
\end{exam}



All three conditions are preserved under equivariant blow-ups.
Among the advantages of
divisorial form are: its formulation does not require a generically free $G$-action,
so for instance we have
Example \ref{exa.abeliandivisorial} also 
without the requirement of a generically free $G$-action;
divisorial form
is stable under general equivariant morphisms,
as we may observe by the following characterization.

\begin{prop}
\label{prop.divisorialform}
Let $X$ be a smooth $G$-variety, where the $G$-action on $X$ is generically free.
Then the action is in divisorial form if and only if the action has abelian stabilizers and at every $\bar x\in X(\bar k)$ the characters of $\mathrm{Stab}(\bar x)$, determined by the elements of $\Pic(X,G)$,
generate
$\mathrm{Stab}(\bar x)^\vee$.
\end{prop}

\begin{proof}
Given $\bar x\in X(\bar k)$, over a closed point $x\in X$ we let $H=\mathrm{Stab}(\bar x)$ and $F\in \mathcal{S}(X,G)$ be the unique component of $X^H$ containing $x$.
The action of $G$ on $X$ restricts to the trivial action of $H$ on $F$, and in particular, on the local ring $\cO_{F,x}$.
The composite map in the definition of divisorial form factors through $\rH^1(H,\cO_{F,x}^\times)$ and agrees with the composite with
\[ \rH^1(H,\cO_{F,x}^\times)\to \rH^1(H,\bar k^\times)\cong H^\vee. \]
So the condition from the definition of divisorial form, for $F$ as above, is equivalent to generation of $\mathrm{Stab}(\bar x)$ by elements coming from $\Pic(X,G)$.
\end{proof}



\begin{rema}
\label{rem.finmanylinebundles}
Given a collection of linearized line bundles $\{L_i\}_{i\in I}$ on $X$, we can replace $\Pic(X,G)$ by the subgroup generated by the classes of the $L_i$ in the definition of divisorial form, to get a
variant, \emph{divisorial form with respect to $\{L_i\}_{i\in I}$}.
This notion appears, e.g., in \cite[Sect.\ 5]{KT-struct},
and Proposition \ref{prop.divisorialform}
remains valid for this notion and subgroup.
If $X$ is in divisorial form, then $X$ is in
divisorial form with respect to some $\{L_i\}_{i\in I}$ with $I$ \emph{finite}.
This is observed in \cite[Remark 3.2]{BnG}, which also points out the following stack-theoretic formulation ($I$ finite):
$X$ is in divisorial form with respect to $\{L_i\}_{i\in I}$
if and only if
the associated morphism of stacks
\begin{equation}
\label{eqn.toproductBGm}
[X/G]\to \prod_{i\in I} B\G_m,
\end{equation}
is \emph{representable} (i.e., induces monomorphisms on geometric stabilizers).
\end{rema}



\begin{exam}
\label{exa.notdivisorial}
Divisorial form has abelian stabilizers, but not every action with abelian stabilizers is in divisorial form:
consider the 
projectivization $X$ of the $3$-dimensional irreducible representation of the alternating group $G:=\fA_4$.
There are points with stabilizer $H$ of order $4$.
We compute $\Pic(X,G)$ using exact sequence
\eqref{eqn.BrXmodG}:
the torsion is $\Hom(\fA_4,k^\times)\cong \Z/3\Z$, the free part is spanned by the class of $\cO_X(1)$ with its natural linearization.
There is no surjective homomorphism from
$\Pic([X/G])\cong \Z\oplus \Z/3\Z$
to
$H^\vee\cong (\Z/2\Z)^2$.
\end{exam}


\begin{exam}
\label{exa.notstandardform}
Standard form is divisorial, but not every divisorial action is in standard form:
consider $X:=\bP^3$ with action of
$G:=(\Z/2\Z)^4$, where respective generators send $(x:y:z:w)$ to
\[ (x:y:-z:-w),\quad
(z:w:x:y),\quad
(x:-y:z:-w),\quad
(y:x:w:z). \]
Every subgroup of $G$ of order $2$ appears as generic stabilizer along curves in $X$.
Since $G$ is abelian, the action is in divisorial form.
With
\[ Q:\ x^2+y^2+z^2+w^2=0, \]
we get $\Pic(X,G)=\langle[Q]\rangle\oplus G^\vee$, by \eqref{eqn.BrXmodG}.
For
a subgroup $H\subseteq G$, $|H|=2$, the homomorphism $\Pic(X,G)\to H^\vee$ (definition of divisorial form) maps $[Q]$ to $1\in H^\vee$ if $C\subset Q$, otherwise to $0$, for a curve $C\in \mathcal{S}(X,G)$ with $\mathrm{I}(C)=H$, and on $G^\vee$ is given by restriction.
Such $C$ can be contained in at most two boundary components of a standard form.
But exhaustive checking reveals: any subset $T\subseteq \Pic(X,G)$, generating $H^\vee$ for every subgroup $H\subseteq G$, $|H|=2$, has to have at least $3$ elements mapping to $1\in H^\vee$ for some subgroup $H\subseteq G$, $|H|=2$.
\end{exam}

\begin{prop}
\label{prop.divisorialcentral}
Let $G$ be a finite group with generically free action on a smooth variety $X$.
If the action is in divisorial form, then
$\mathrm{I}(F)$ lies in the center of $\mathrm{D}(F)$, for every
$F\in \mathcal{S}(X,G)$.
\end{prop}

\begin{proof}
With the notation of the definition of divisorial form, the action of $Y$ on $\rH^1(H,k(F)^\times)$ has to be trivial.
Since the action is given by conjugation,
$\mathrm{I}(F)$ is a central subgroup of $\mathrm{D}(F)$.
\end{proof}

\begin{exam}
\label{exa.morenotdivisorial}
In Example \ref{exa.notcentral},
$\mathrm{I}(F)$ does not lie in the center of $\mathrm{D}(F)$.
With this linear action of $\mathfrak{D}_4$ on $\A^3$, or its projective compactification $\bP^3$, we have additional examples of actions with abelian stabilizers, not in divisorial form.
\end{exam}

\begin{prop}
\label{prop.divisorial}
Let a generically free action of a finite group $G$ on a smooth algebraic variety $X$ be given, together with a simple normal crossing divisor
\[ D=D_1\cup\dots\cup D_\ell, \]
on $X$ such that each $D_i$ is smooth and $G$-invariant.
Given $\bar x\in X(\bar k)$ over $x\in X$, we define
\[ H'\subseteq H:=\mathrm{Stab}(\bar x) \]
to be the intersection of the kernels of the characters determined by the $G$-linearized line bundles $\cO_X(D_i)$ for all $i$ with $x\in D_i$, and let $d(\bar x)$ denote the dimension of the nontrivial part of the representation of $H'$ on the tangent space $\cT_{X,\bar x}$.
Then the function $d$ is identically zero if and only if the action is in standard form with boundary $D$.
If the maximal value $m$ of $d$ is positive, then
\[ \{\bar x\in X(\bar k)\,|\, d(\bar x)=m\} \]
is the set of $\bar k$-points of a smooth $G$-invariant closed subscheme $W$ of pure codimension $m$, meeting $D$ transversally, such that the blow-up $\widetilde{X}$ of $X$ along $W$ has a simple normal crossing divisor formed by the exceptional divisor and the pre-images of the $D_i$, and $d(\bar y)<m$ for all $\bar y\in \widetilde{X}(\bar k)$.
\end{prop}

\begin{proof}
We start by proving the first assertion.
For an action in standard form with boundary $D$, the generation condition of Proposition \ref{prop.divisorialform} is satisfied at $\bar x\in X(\bar k)$
over $x\in X$ with just the $G$-linearized line bundles $\cO_X(D_i)$ with $x\in D_i$,
thus $d(\bar x)=0$.
The converse follows from the observation, that for $x$ in the complement of $D$, the vanishing of $d(\bar x)$ implies the triviality of $H$.

For the remaining assertions, the condition to have geometric stabilizer group contained in $H$ defines $H$-invariant open $X^\circ\subset X$, with $x\in X^\circ$.
At any $\bar k$-point of $X^\circ$ the value of the function $d$ stays the same when we replace $G$ by $H$.
The function $d$ also remains unchanged under passage to an equivariant \'etale cover that induces bijections on geometric stabilizer groups.
So we are reduced to checking the remaining assertions under the assumption that
$G=H$ and
$x$ is the origin in a linear representation $X=V$ for some $G\to \GL(V^\vee)$,
\[ V=\chi_1\oplus\dots\oplus \chi_\ell\oplus V', \]
direct sum of $\ell$ one-dimensional representations of $G$
and a final factor $V'$, with
the $D_i$ given by the first $\ell$ factors.

The subgroup $H'$ is
the intersection of the kernels of $\chi_1$, $\dots$, $\chi_\ell$.
We write
\[ V'=(V')^{H'}\oplus V'' \]
for some subrepresentation $V''$.
Now $m=\dim(V'')$, with
$d(\bar v)\le m$ for all $\bar v\in V(\bar k)$ and equality if and only if $\bar v$ projects to $0\in V''$.
So
\[ W=\chi_1\oplus \dots\oplus \chi_\ell \oplus (V')^{H'}. \]
Blowing up $W$ in $V$ replaces $V''$ by
$B\ell_0(V'')$.
At a $\bar k$-point $\bar y$ of the exceptional divisor,
the nontrivial part of the representation on $\mathcal{T}_{B\ell_WV,\bar y}$ of the intersection of kernels of characters of $\mathrm{Stab}(\bar y)$ is contained in the tangent space to the image of $\bar y$ in $\bP(V'')$.
This has dimension $m-1$.
\end{proof}

\begin{coro}[\cite{berghrydh}]
\label{cor.divisorial}
A smooth algebraic $G$-variety, where the $G$-action is generically free, can be brought into standard form via a sequence of blow-ups at $G$-invariant smooth centers.
\end{coro}

\begin{proof}
Starting with the empty divisor, we repeatedly blow up the locus indicated in Proposition \ref{prop.divisorial}, until the function $d$ is identically zero.
\end{proof}

\begin{rema}
\label{rem.divisorial}
The function $d$ of Proposition \ref{prop.divisorial} depends only on the closed point $x\in X$ and may be extended to arbitrary points of $X$, by replacing the tangent space at $\bar x$ by the dual of the fiber at $x$ of the conormal sheaf of the closure of $\{x\}$.
Then we recover the \emph{divisorial index} of \cite{berghrydh}.
Corollary \ref{cor.divisorial} records the outcome of the \emph{divisorialification} algorithm of \cite[Thm.\ A]{berghrydh}.
\end{rema}

\section{Equivariant Burnside groups -- definitions}
\label{sect:defi}
Consider a smooth projective $G$-variety $X$ of dimension $n$, where the base field $k$ is assumed to contain enough roots of unity; the precise requirement for ``enough'' is stated at the beginning of Section~\ref{sect:model}. 
We assume that the $G$-action is generically free.
Let $\bar x$ be a $\bar k$-point of $X$, with abelian geometric stabilizer group $H\subseteq G$.
The equivalence class of the induced representation of $H$ on the tangent space $\cT_{X,\bar x}$ may be encoded by a sequence
$$
\beta(\bar x):=(b_1,\ldots, b_n) 
$$
of characters of $H$.
If $\bar x$ is an isolated point of $X^H$, then all of the characters will be nontrivial. 
In general, the multiplicity of the zero character will be the
codimension in $X$ of the component of $X^H$ containing $\bar x$.
The characters $b_1$, $\dots$, $b_n$ will generate the character group $H^\vee$.

We formulate, first, the groups of symbols that encode these actions.
At a $\bar k$-point $\bar x\in X$
with abelian geometric stabilizer group $H$, we get a sequence of characters of $H$, that is defined up to order.
At another point in the $G$-orbit of $\bar x$ the geometric stabilizer group $H\subseteq G$ gets replaced by a conjugate subgroup $H':=gHg^{-1}$,
and the characters get replaced accordingly by their $g$-conjugates.
The corresponding symbol groups are subject to order and conjugation relations.
In a next step, blow-up relations are introduced; only after introducing these relations do we obtain groups, where the equivariant birational invariants take their values.

\subsection{Symbols}
\label{ss.symbols}

Let $H$ be a finite abelian group, with character group
\[ A:=H^\vee. \]
For $n\in \N$, we introduce
\[ \cS_n(H), \]
the abelian group generated by \emph{symbols}, which take the form
of an $n$-tuple of characters
\[ \beta=(b_1,\dots,b_n),\qquad b_1,\dots,b_n\in A, \]
such that $\langle b_1,\dots,b_n\rangle=A$.
These are subject to the relation 

\noindent
\textbf{(O)} (order)
$\beta=(b_1,\dots,b_n)$ is equal to $\beta'=(b'_1,\dots,b'_n)$
if there exists a permutation $\sigma\in \fS_n$,
with $b'_i=b_{\sigma(i)}$ for $i=1$, $\dots$, $n$.

Notice that we allow $0$ to appear as a character in $\beta$.

\medskip

Let $G$ be a finite group.
The \emph{combinatorial symbols group}
\[ \mathcal{SC}_n(G) \]
is generated by \emph{symbols}
\[ (H,Y,\beta) \]
where
\begin{itemize}
\item 
$H\subseteq G$ is an abelian subgroup,
\item $Y$ is a subgroup of $Z_G(H)/H$ (with $Z_G(H)$ the centralizer of $H$ in $G$),
and
\item $\beta$ is a sequence of characters in $H^\vee$, generating $H^\vee$, of length at most $n$.
\end{itemize}
These are subject to relations 

\noindent
\textbf{(O)} $(H,Y,\beta)=(H,Y,\beta')$ if $\beta'$ is a reordering of $\beta$, 

\noindent
\textbf{(C)} (conjugation) For $g\in G$ we have $(H,Y,\beta)=(H',Y',\beta')$, where $H'=gHg^{-1}$,
$Y'=gYg^{-1}$, and by conjugation by $g$ we obtain the characters in $\beta'$ from those in $\beta$.

\medskip

Lastly, we introduce the abelian group
\[ \mathrm{Symb}_n(G), \]
where the generators are \emph{symbols}
\[
(H, Y\actsfromleft K, \beta).
\]
Here,
\begin{itemize}
\item $H\subseteq G$ is an abelian subgroup,
\item $Y\subseteq Z_G(H)/H$ is a subgroup,
\item $\beta$ is a sequence of characters in $H^\vee$, generating $H^\vee$,
\item $K$ is a finitely generated extension field of $k$, with faithful action over $k$ by $Y$, such that, if we denote by $Y_0$ the pre-image of $Y$ in $Z_G(H)$ (by the canonical homomorphism), the restriction map
\begin{equation}
\label{eqn:su}
\rH^1(Y_0,K^\times)\to \rH^1(H,K^\times)\cong H^\vee
\end{equation}
is surjective, and
\item denoting by $d$ the transcendence degree of $K$ over $k$,
the length of the sequence of characters $\beta$ is $n-d$.
\end{itemize}
These symbols are subject to relations 

\noindent
\textbf{(O)} $(H, Y\actsfromleft K, \beta)=(H, Y\actsfromleft K, \beta')$ if $\beta'$ is a reordering of $\beta$,

\noindent
\textbf{(C)} For $g\in G$ we have $(H, Y\actsfromleft K, \beta)=(H', Y'\actsfromleft K', \beta')$, where
$$
H'=gHg^{-1}, \quad Y'=gYg^{-1},
$$
$K'$ is isomorphic as a $k$-algebra to $K$ in a manner that is compatible with the
respective actions,
and $\beta'$ obtained from $\beta$ by conjugation by $g$.

\begin{rema}
\label{rem.injective}
The surjectivity condition \eqref{eqn:su} here was called ``Assumption 1'' in \cite{BnG}.
The map \eqref{eqn:su} is extraced from the
inflation-restriction exact sequence of group cohomology.
This has, further to the left, the term
$\rH^1(Y,K^\times)$, which vanishes by Hilbert's Theorem 90.
Consequently, surjectivity of \eqref{eqn:su} is equivalent
to having here an isomorphism.
\end{rema}


\subsection{Relations}
\label{ss.relations}
We define relations, mirroring the impact on the weights when a smooth $G$-variety when $X$ is blown up along a smooth invariant $G$-subvariety.
It suffices to consider the case of a blow-up in codimension $2$ (this will be explained later, in the proof of Theorem \ref{thm.biratinvarclass}); this is reflected in the form of the relations.

\

We consider the quotient 
$$
\cS_n(H) \to \cB_n(H)
$$
by the relation:


\

\textbf{(B)} (blow-up) For $\beta=(b_1,\dots,b_n)$, $n\ge 2$,
\[ 
\beta=\begin{cases}
(0,b_2,\dots,b_n),&
\text{if $b_1=b_2$}, \\
\beta_1+\beta_2,&
\text{if $b_1\ne b_2$},
\end{cases}
\]
where
\[
\beta_1:=(b_1-b_2,b_2,b_3,\dots,b_n),\qquad
\beta_2:=(b_1,b_2-b_1,b_3,\dots,b_n).
\]

\medskip

Similarly, we consider the quotient
$$
\mathcal{SC}_n(G) \to \mathcal{BC}_n(G)
$$
by the relations

\noindent
\textbf{(V)} (vanishing) $(H,Y,\beta)=0$ whenever $b_i=0$ for some $i$, 

\noindent
\textbf{(B)} if the sequence of characters $\beta=(b_1,b_2,\dots)$ has length $\ge 2$, then
$$
(H,Y,\beta)=(H,Y,\beta_1)+(H,Y,\beta_2)+(\overline{H},\overline{Y},\bar\beta),
$$
where $\beta_1:=(b_1-b_2,b_2,\dots)$, $\beta_2:=(b_1,b_2-b_1,\dots)$, $\overline{H}:=\ker(b_1-b_2)$, $\bar\beta:=(\bar b_2,\bar b_3,\dots)$ (restrictions of characters), and $\overline{Y}:=Y_0/\overline{H}$,
with $Y_0$ the pre-image of $Y$ in $Z_G(H)$, as before.

\medskip

Finally, we consider the quotient 
$$
\mathrm{Symb}_n(G) \to \Burn_n(G)
$$
by the relations 

\noindent
\textbf{(V)} $(H,Y\actsfromleft K,\beta)=0$ whenever $b_i=0$ for some $i$, 

\noindent
\textbf{(B)} for a symbol $(H,Y\actsfromleft K,\beta)$ with $n-d\ge 2$
(where $d$ denotes the transcendence degree of $K$ over $k$),
\begin{align}
\label{eqn:v}
(H,Y\actsfromleft K&,\beta)= \notag \\
&(H,Y\actsfromleft K,\beta_1)+(H,Y\actsfromleft K,\beta_2)+(\overline{H},\overline{Y}\actsfromleft \overline{K},\bar\beta),
\end{align}
with notation as above and additionally
\[
\overline{K}:=K(t),
\]
where the following recipe is carried out to produce an action of $\overline{Y}$ on $\overline{K}$.
Consider $b:=b_1-b_2$ as a primitive character of $H/\overline{H}$.
By the surjectivity of \eqref{eqn:su} and Remark \ref{rem.injective}, $b$ admits a unique lift to $\rH^1(Y_0,K^\times)$.
A choice of cocycle representation leads to a $Y_0$-action on $K(t)$,
where $\overline{H}$ acts trivially, and we get an induced action of $\overline{Y}$; see 
\cite[Sect.\ 2]{BnG}.

\begin{rema}
\label{rem.earliersymbols}
The formulation of relations \textbf{(V)} and \textbf{(B)} follows \cite[Sect.\ 2]{KT-struct}.
Earlier papers such as \cite{BnG} and \cite{KT-vector} allowed only symbols with sequences of \emph{nonzero} characters, generating $H^\vee$, with relations \textbf{(B1)}--\textbf{(B2)}, that are however equivalent to the formulation with more general symbols and relations \textbf{(V)} and \textbf{(B)}.
An additional, conventional difference in the present formulation is the requirement for $K$ to be a \emph{field}, rather than a finite product of fields, as in \cite{BnG} (formulation with Galois algebras for the group $N_G(H)/H$), or \cite{KT-struct} (with Galois algebras for subgroups of $N_G(H)/H$);
here, $N_G(H)$ denotes the normalizer of $H$ in $G$.
\end{rema}

\subsection{Class of the action} 
\label{sect:class}
Given a good model for the $G$-action, i.e., one that is in \emph{divisorial form}, as in Section~\ref{sect:standardf}, we define its class in $\cB_n$, $\BC_n$, respectively, $\Burn_n$ as follows: 
\begin{itemize}
\item For $G$ {\em abelian}, we let $\cF(X,G)$ be the set of components of the fixed locus
\[ X^G=\bigsqcup_{F\in \cF(X,G)} F \]
and, for each, record the characters of the $G$-action 
$$
\beta_{F}(X)=(b_1,\ldots, b_n), \quad b_j\in G^\vee,
$$   
in the tangent bundle at a geometric point of $F$. The class 
\[
[X\actsfromright G]:=\sum_{F\in \cF(X,G)} [k':k]\beta_{F}(X)\in \cB_n(G)
\]
with $k'$ the algebraic closure of $k$ in $k(F)$,
yields a well-defined $G$-equivariant birational invariant \cite[Thm.\ 3]{KPT}.
\item 
For arbitrary $G$, we consider the stabilizer stratification $\cS(X,G)$ introduced in Section~\ref{sect:stab}; recall that all stabilizers are abelian.  
For every component $F\in \cS(X,G)$
of dimension $d$ we record a symbol
$$
(H, Y, \beta)\qquad\text{resp.}\qquad
(H, Y\actsfromleft k(F), \beta), 
$$
where $H=\mathrm{I}(F)$ (an abelian subgroup of $G$), and the symbol records
$Y=\mathrm{D}(F)/H\subseteq Z_G(H)/H$ (containment by Proposition \ref{prop.divisorialcentral}),
respectively $Y$ with residual action,
and
$\beta=\beta_F(X)=(b_1,\ldots, b_{n-d})$ is the sequence of characters of $H$, for the action of $H$ 
in the normal bundle to $F$ in $X$ at the generic point of $F$. We record only one such symbol for the $G$-orbit of $F$. In particular, $H$, $Y$, and $\beta$ are defined only up to conjugation (reflecting the passage to a different component in the $G$-orbit of $F$); furthermore, the sequence $\beta$ is only defined up to order.
The class is
$$
[X\actsfromright G] :=\sum_{F\in \cS(X,G)/G} [k':k](H, Y, \beta_F(X))\in\BC_n(G),
$$
with $k'$ as before the algebraic closure of $k$ in $k(F)$, respectively,
$$
[X\actsfromright G] :=\sum_{F\in \cS(X,G)/G}  (H, Y\actsfromleft k(F), \beta_F(X))\in\Burn_n(G),
$$
where the sum runs over $G$-orbit representatives of the stabilizer stratification.
\end{itemize}

\begin{rema}
\label{rem.classofaction}
In parallel to the conventional differences in the definition of symbols, mentioned in Remark \ref{rem.earliersymbols}, there are alternative expressions for $[X\actsfromright G]$.
In \cite{BnG}, $[X\actsfromright G]\in \Burn_n(G)$ is expressed as a sum over conjugacy class representatives $H$ of abelian subgroups of $G$, of sums of symbols $\sum_F \mathfrak{s}_F$, where $F$ in the inner sum runs over $N_G(H)$-orbits of elements of $\mathcal{S}(X,G)$ with generic stabilizer $H$.
In the analogous expression in \cite{KT-vector}, the inner sum is over $Z_G(H)$-orbits, but this comes at the cost of a more complicated outer sum, requiring the notion of $\Pic(X,G)$-pairing on an abelian subgroup of $G$; an advantage of the shift to $Z_G(H)$-orbits was a simpler formulation of the blow-up relations.
The closest to the present formulation is \cite[p.~3027]{KT-toric}, an expression as a sum over points $x_0\in X/G$ (in the scheme sense, i.e., not just closed points).
When we have the generic point of some $F\in \mathcal{S}(X,G)/G$ over $x_0$, the contribution is $(H,Y\actsfromleft k(F),\beta_F(X))$.
Otherwise, the contribution is a symbol with $0$ in the sequence of characters, and this vanishes by relation \textbf{(V)}.
More details are provided in \cite[Sect.\ 2]{KT-struct}.
\end{rema}

\begin{theo}
\label{thm.biratinvarclass}
The recipe, for a projective $n$-dimensional $G$-variety with generically free $G$-action, to choose a smooth projective model $X$ in divisorial form and 
associate the element
\[ [X\actsfromright G], \]
defines a $G$-equivariant birational invariant class in $\BC_n(G)$, respectively in $\Burn_n(G)$.
\end{theo}

We recall, a smooth projective model in divisorial form may be obtained by performing equivariant resolution of singularities, followed by the sequence of blow-ups given by Corollary \ref{cor.divisorial}.

\begin{proof}
It needs to be checked that two different smooth projective models in divisorial form give rise to the same class in $\BC_n(G)$, respectively in $\Burn_n(G)$.
This is established in \cite[Thm.\ 5.1]{BnG} by using
equivariant weak factorization to reduce to the case of a blow-up of a smooth $G$-subvariety of codimension $j\ge 2$, and relations in $\Burn_n(G)$
of the form
\[
(H,Y\actsfromleft K,\beta)=
\sum_{\emptyset\ne I\subset \{1,\dots,j\}} 
(H_I,Y_I\actsfromleft K_I,\beta_I),
\]
provided $\mathrm{trdeg}_k(K)\le n-j$
(or analogous relations in $\BC_n(G)$);
a key point is that these relations are shown, in \cite[Prop.\ 4.7(ii)]{BnG}, to follow from relations \textbf{(B)} and \textbf{(V)}.
Here,
\[ H_I:=\bigcap_{i,i'\in I}\ker(b_i-b_{i'}) \]
and, writing $Y=Y_0/H$ with $Y_0\subseteq Z_G(H)$,
\[ Y_I:=Y_0/H_I. \]
We construct an action of $Y_I$ on the purely transcendental extension
$K_I$ of $K$, $\mathrm{trdeg}_K(K_I)=|I|-1$,
in an analogous fashion to
$\overline{Y}\actsfromleft \overline{K}$ in \eqref{eqn:v}; this is the {\em action construction} from \cite[Sect.\ 2]{BnG}, see also \cite[Sect.\ 2]{KT-vector}. 
All $b_i$ for $i\in I$ have a common class $\bar b\in (H_I)^\vee$; then,
\[ \bar \beta:=(\bar b,\bar b_{i_1},\dots, \bar b_{i_r}), \]
with $\{1,\dots,n-d\}\setminus I=\{i_1,\dots,i_r\}$.
\end{proof}

Algorithms for the computation of the class $[X\actsfromright G]\in \Burn_n(G)$ for actions arising from (projectivizations of) linear representations and actions on smooth projective toric varieties can be found in \cite{KT-vector} and \cite{KT-toric}. 



\subsection{Class of an open subvariety}
\label{ss.opensubvariety}
In applications, we also need to consider a $G$-invariant open $U$ in a smooth projective $G$-variety $X$.
We suppose that the $G$-action on $X$ is generically free, and $X$ is in divisorial form.
There are \emph{two} ways to define the class of $U\actsfromright G$ (in any flavor of Burnside group).
The most naive way to do this is to copy the formula from Section \ref{sect:class}, replacing $X$ by $U$:
\begin{align*}
[U\actsfromright G]^{\mathrm{naive}}&:=\sum_{F\in \cF(U,G)} [k':k]\beta_{F}(U)\in \cB_n(G), \\
[U\actsfromright G]^{\mathrm{naive}}&:=\sum_{F\in \cS(U,G)/G} [k':k](H, Y, \beta_F(U))\in\BC_n(G), \\
[U\actsfromright G]^{\mathrm{naive}}&:=\sum_{F\in \cS(U,G)/G}  (H, Y\actsfromleft k(F), \beta_F(U))\in\Burn_n(G).
\end{align*}
These classes are birational invariants of $U\actsfromright G$, by \cite[Lemma 5.3]{BnG}.

However, motivated by the form of the specialization map, even in its original, nonequivariant form \cite{KT},
we also need a more sophisticated formula with an alternating sum over boundary strata.
For this we make the further assumption, that $X\setminus U$ is a simple normal crossing divisor \[ D=D_1\cup\dots \cup D_\ell, \]
such that each $D_i$ is $G$-invariant.
For $\emptyset\ne I\subseteq \cI:=\{1,\dots,\ell\}$ we have $D_I;=\bigcap_{i\in I}D_i$,
with normal bundle $\mathcal{N}_{D_I/X}$
and \emph{punctured normal bundle} $\mathcal{N}_{D_I/X}^\circ$.
The latter is defined by removing, from $\mathcal{N}_{D_I/X}\cong \bigoplus_{i\in I}\mathcal{N}_{D_i/X}|_{D_I}$,
the fiber over $D_{I\cup\{j\}}$, for all $j\in \cI\setminus I$, as well as the zero-section of $\mathcal{N}_{D_i/X}|_{D_I}$, for all $i\in I$.
Then in any flavor of Burnside group we have
\begin{align*}
[&U\actsfromright G]:=
[X\actsfromright G]+\sum_{\emptyset\ne I\subseteq \cI}(-1)^{|I|}[\mathcal{N}_{D_I/X}\actsfromright G]^{\mathrm{naive}} \\
&\qquad=
[U\actsfromright G]^{\mathrm{naive}}+\sum_{\emptyset\ne I\subseteq \cI}(-1)^{|I|}[\mathcal{N}_{D_I/X}^\circ\actsfromright G]^{\mathrm{naive}}.
\end{align*}
The definition is \cite[Defn.\ 5.4]{BnG}, and the equality is \cite[Lemma 5.7]{BnG}.
This class is a birational invariant of $U\actsfromright G$, by \cite[Thm.\ 5.15]{BnG}.



\subsection{Comparisons}
\label{ss.comparisons}
The different versions of Burnside groups are related to each other, via forgetful homomorphisms. The homomorphism
$$
\Burn_n(G)\to \BC_n(G)
$$
forgets the birational type of the stratum $F$ in the symbol \cite[Prop.\ 8.2]{KT-struct}: 
$$
(H, Y\actsfromleft k(F), \beta) \mapsto [k':k](H, Y, \beta).
$$
For $G$ abelian, we 
have a further homomorphism \cite[Exa.\ 8.8]{KT-struct}
$$
\BC_n(G)\to \cB_n(G), 
$$
where we quotient by the subgroup generated by symbols with $H\subsetneq G$. 
A more refined formalism of intermediate quotients, with respect to natural filtrations on $\Burn_n(G)$ is described in Section~\ref{sect:filtration}.

\section{Equivariant Burnside groups -- properties}
\label{sect:first}
We continue to assume that $k$ has enough roots of unity, as in Sections \ref{sect:model} and \ref{sect:defi}.


\subsection{Vanishing in stable range}
\label{sect:vanishing}

In the analysis of relations in the various Burnside groups, we often use the following combinatorial consequence of the blow-up relation \cite[Prop.\ 4.7(i)]{BnG}: 
a symbol
$$
(H, Y\actsfromleft K, (b_1,\ldots, b_{n-d})) 
$$
vanishes in $\Burn_n(G)$, if $\sum_{i\in I}b_i=0$ for some $\emptyset\ne I\subseteq \{1,\dots,n-d\}$.
This implies the vanishing of any symbol of the form
\[ (H, Y\actsfromleft K(t_1,\dots,t_m), (b_1,\ldots, b_{n-d})) \]
where $m$ is at least one less than the minimum of the orders of the $b_i$ \cite[Prop.\ 4.1]{KT-toric}.
As a consequence, if $X$ is a smooth projective $n$-dimensional variety with generically free $G$-action, and we consider $X\times \bP^m$, with trivial $G$-action on $\bP^m$, then
$$
[X\times \bP^m\actsfromright G]=(1, G\actsfromleft k(X)(t_1,\ldots, t_m), ())\in \Burn_{n+m}(G),
$$
provided $m\ge -1+\max_{g\in G}|\langle g\rangle|$ \cite[Rmk.\ 4.2]{KT-toric}. 
(More generally, if $X$ is a smooth projective $G$-variety and $X_0$ an irreducible component of $X$, then this holds with $(1,\mathrm{D}(X_0)\actsfromleft k(X_0)(t_1,\dots,t_m),())$ on the right.)





\subsection{Filtrations}
\label{sect:filtration}
There is the possibility to restrict attention just to certain
stabilizer groups.
This can be done systematically using the notion of filter,
introduced in \cite[\S 3]{KT-struct}.
A \emph{$G$-filter} is a subset $\mathbf{H}$ of the set of pairs $(H,Y)$, consisting of an abelian subgroup $H\subseteq G$ and a subgroup $Y\subseteq Z_G(H)/H$,
that is closed under conjugation and satisfies the following additional property.
For $(H,Y)\in \mathbf{H}$, $H\ne 1$, and $g\in Z_G(H)$, with
$\bar g\in Y$ and $Y\subseteq Z_G(g)/H$, we require
$(\langle H,g\rangle,Y/\langle \bar g\rangle)\in \mathbf{H}$.
Then
$\Burn_n^{\mathbf{H}}(G)$
is defined to be the quotient of $\Burn_n(G)$ by the subgroup generated by all symbols $(H,Y\actsfromleft K,\beta)$ with $(H,Y)\notin \mathbf{H}$.
A similar definition is made for $\BC_n(G)$.
The properties of a filter guarantee \cite[Prop.\ 3.3, Prop.\ 8.7]{KT-struct} that
$\Burn_n^{\mathbf{H}}(G)$ is the abelian group,
generated by symbols $(H,Y\actsfromleft K,\beta)$ with $(H,Y)\in \mathbf{H}$, subject to relations
\textbf{(O)}, \textbf{(C)}, \textbf{(V)}, \textbf{(B)}, applied just to these symbols,
and the analogous fact holds for $\BC_n^{\mathbf{H}}(G)$.

There is also a comparison homomorphism between the Burnside and combinatorial Burnside groups (\S \ref{ss.comparisons}) in the presence of a filter $\mathbf{H}$, and this is compatible with the quotient maps $\Burn_n(G)\to \Burn_n^{\mathbf{H}}(G)$ and
$\BC_n(G)\to \BC_n^{\mathbf{H}}(G)$.



For instance, when $G$ is abelian we define (cf.\ \cite[\S 8]{BnG})
\[ \Burn_n^G(G):=\Burn_n^{\{(G,1)\}}(G),\qquad
\BC_n^G(G):=\BC_n^{\{(G,1)\}}(G).
\]
By \cite[Exa.\ 8.8]{KT-struct}, we have
\[ \BC_n^G(G)=\cB_n(G). \]




\subsection{Incompressibles}
\label{sect:incomp}
The analysis of the blow-up relation, and in particular, the determination of the vanishing of a symbol in $\Burn_n(G)$ can be tricky.
However, in many applications, $G$-actions can be distinguished upon projection to the subgroup, freely generated by {\em incompressible} symbols. These are divisor symbols, that is, symbols 
$$
(H,Y\actsfromleft K,(b))
$$ 
with $\mathrm{trdeg}_k(K)=n-1$ and nontrivial cyclic $H$, $H^\vee=\langle b\rangle$, that cannot be obtained, via the action construction, as rightmost term in a relation \eqref{eqn:v}, coming from relation \textbf{(B)}.
The notion appears in \cite[Defn.\ 3.3]{KT-vector}.

Among classical precursors of this notion are fixed curves of birational involutions on rational surfaces, leading to the classification of involutions as Bertini, Geisser, and
de Jonqui\`eres; cf.\ \cite{baylebeauville}. A more recent version of this is the {\em normalized fixed curve with action (NFCA)} invariant, introduced by Blanc \cite{blancsubgroups}.
This takes into account the stabilizer and the residual action, and gives a finer invariant for cyclic groups of nonprime order.      










\subsection{Specialization}
Let $\mathfrak{o}$ be a complete DVR with residue field $k$,
and let $K$ be the fraction field of $\mathfrak{o}$.
Fix a uniformizer $t\in \mathfrak{o}$;
so, $\mathfrak{o}\cong k[[t]]$.
We proceed to describe the  specialization map
\[ \rho_t^G\colon \Burn_{n,K}(G)\to \Burn_n(G). \]
The definition works with symbols in $\Burn_{n,K}(G)$, but the motivation is the specialization of $[X\actsfromright G]\in \Burn_{n,K}(G)$, where $X$ is a smooth projective $G$-variety over $K$ with generically free $G$-action.

We take $\cX$ to be a regular model, projective over $\mathfrak{o}$, such that the special fiber is a simple normal crossing divisor $D=D_1\cup\dots\cup D_\ell$, where the $G$-action extends to $\cX$, each $D_i$ is $G$-invariant, and there is $d_i\in \mathbb{N}$, common multiplicity of the components of $D_i$ in the special fiber.
The set-up reminds us of Section \ref{ss.opensubvariety}, and we have similarly the normal bundle $\mathcal{N}_{D_I/\cX}$ and punctured normal bundle $\mathcal{N}_{D_I/\cX}^\circ$ for $\emptyset\ne I\subseteq \cI:=\{1,\dots,\ell\}$.
We have a map, given by the canonical section of a twist of $\cO_{D_I}$, and isomorphism,
supplied by $t$:
\[
\bigotimes_{i\in I} \mathcal{N}_{D_i/\cX}^{\otimes d_i}|_{D_I}\to 
\biggl(\bigotimes_{i\in I} \mathcal{N}_{D_i/\cX}^{\otimes d_i}|_{D_I}\biggr)\otimes
\cO_{D_I}\Bigl(\sum_{j\in \cI\setminus I}d_jD_{I\cup\{j\}}\Bigr)\cong \cO_{D_I}.
\]
We define the regular function $\omega_I$ on $\mathcal{N}_{D_I/\cX}$ via the natural identification with $\bigoplus_{i\in I}\mathcal{N}_{D_i/X}|_{D_I}$,
to be the composite with the morphism
\[
\bigoplus_{i\in I}\mathcal{N}_{D_i/X}|_{D_I}\to
\bigotimes_{i\in I} \mathcal{N}_{D_i/\cX}^{\otimes d_i}|_{D_I},\qquad
(s_i)_{i\in I}\mapsto \prod_{i\in I}s_i^{d_i}.
\]
The locus of vanishing of $\omega_I$ is the union of the fibers over $D_{I\cup\{j\}}$ for $j\in \cI\setminus I$ and the zero-sections of the $\mathcal{N}_{D_i/X}|_{D_I}$.
Thus
\[
\omega_I^{-1}(\A^1\setminus\{0\})=\mathcal{N}_{D_I/\cX}^\circ.
\]

The specialization map is defined as follows.
Given a symbol
\[ (H,Y\actsfromleft K(W),\beta) \]
with $W$ a smooth projective $G$-variety over $K$, we take
$\mathcal{W}$ to be a regular model of $W$
as above.
So the $G$-action extends to
$\mathcal{W}$, projective over $\mathfrak{o}$,
with simple normal crossing special fiber
$D_1\cup\dots\cup D_\ell$, such that each $D_i$ is $G$-invariant with a multiplicity $d_i$.
There is the regular function $\omega_I$ on $\mathcal{N}_{D_I/\cW}$, as above.
Then
\[
\rho_\pi^G((H,Y\actsfromleft K(W),\beta)):=\sum_{\emptyset\ne I\subseteq \cI}(-1)^{|I|-1}(H,Y\actsfromleft k(\omega_I^{-1}(1)),\beta),
\]
which is well-defined, i.e.,
independent of the choice of $\cW$ \cite[Sect.\ 6]{BnG}.

We return to our motivation, to understand the image of $[X\actsfromright G]$ under $\rho_t^G$.
This was addressed in \cite[Thm.\ 6.6]{BnG}, with a formula as an explicit sum of symbols.
We state this in a conceptually simpler form.
\begin{theo}
\label{thm.revisitThm6.6}
Let $X$ be a smooth projective $G$-variety over $K$ with generically free $G$-action, and let $\cX$ be a regular model to which the action extends, such that the special fiber is a simple normal crossing divisor $D=D_1\cup \dots\cup D_\ell$, where each $D_i$ is $G$-invariant with a multiplicity $d_i$.
Then with $\omega_I$ as above we have
\[
\rho_t^G([X\actsfromright G])=
\sum_{\emptyset\ne I\subseteq \cI}
(-1)^{|I|-1}[\omega_I^{-1}(1)\actsfromright G]^{\mathrm{naive}}.
\]
\end{theo}

A technical point is that we wish to employ geometric arguments over $k$, but $X$ is defined over $K$.
N\'eron desingularization \cite[\S 3.6]{BLR} lets us replace $\mathfrak{o}$ by a suitable smooth $k[t]$-subalgebra $A$, and $K$ by the fraction field $K_A$ of $A$, thereby enabling direct geometric arguments.
First, $A$ may be chosen with projective $A$-scheme $X_A$, locally a complete intersection with $\cX\cong X_A\times_{\Spec(A)}\Spec(\mathfrak{o})$.
The simple normal crossing boundary translates into a local equation form for $\cX$, which we may assume valid for $X_A$.
Then $X_A$ is smooth over $k$ with simple normal crossing boundary $X_{A/tA}:=X_A\times_{\Spec(A)}\Spec(A/tA)$, and
the same holds when $A$ is replaced by a larger smooth $k[t]$-subalgebra of $\mathfrak{o}$.
By means of monoidal transform we may suppose that $X_{A/tA}$ is equivariantly isomorphic to $ D\times \Spec(A/tA)$ (where $G$ acts trivially on the second factor).

The normal crossing boundary of $X_A$ defines a morphism
\begin{equation}
\label{eqn.XAtoAlmodGml}
[X_A/G]\to [\A^\ell/\G_m^\ell],
\end{equation}
which by \cite[Rmk.\ 3.4]{berghrydh} is smooth.
The stack $[\A^\ell/\G_m^\ell]$ has a dense open substack $[(\A^1\setminus \{0\})^\ell/\G_m^\ell]\cong \Spec(k)$, with
pre-image $X_{A[t^{-1}]}=X_A\setminus X_{A/tA}$ by \eqref{eqn.XAtoAlmodGml}.
The pre-image of $X_{A[t^{-1}]}$ under
$\cX\to X_A$ is $X$.

We observe that
in Theorem \ref{thm.revisitThm6.6} there is no loss of generality in supposing $X$ to be in divisorial form.
For \eqref{eqn.XAtoAlmodGml} there is a relative notion of divisoriality and a functorial procedure to achieve this by iterated blow-up; see \cite[Thm.\ 6.6]{berghrydh}.
Each center of blow-up is smooth over $[\A^\ell/\G_m^\ell]$, i.e., smooth over $k$ and transverse to the normal crossing boundary.
We claim, also, transverse intersection on $\mathcal{N}_{D_I/\cX}^\circ$ with $\omega_I^{-1}(1)$, for every $I$.
This claim does not see the group action; with respect to a local trivialization of $\mathcal{N}_{D_I^\circ/\cX}$, we have $\omega_I$ given as unit times monomial,
so this is clear.
Thus, in Theorem \ref{thm.revisitThm6.6},
relative divisorialification preserves the hypothesis, as well as both sides of the equality
(by birational invariance of the class in the Burnside group).

The notion of divisoriality makes sense also for $\cX$, as explained in \cite[Sect.\ 6]{BnG}: $\cX$, to be divisorial, should have abelian stabilizers with $X$ divisorial, and for every $F\in \mathcal{S}(\cX,G)$ supported in the special fiber, with generic stabilizer group $H$, we should have surjective composite
\[
\Pic(\cX,G)\to 
\rH^1(\mathrm{D}(F), k(F)^\times)
\to \rH^1(H, k(F)^\times)^{\mathrm{D}(F)/H} \to H^\vee.
\]
But if $X$ is divisorial, then $\cX$ is divisorial.
The corresponding assertion for orbifolds is \cite[Lemma 9.1(ii)]{KT-stacks}; we give a proof in the present setting.
Say $X$ is divisorial with respect to linearized line bundles $L_1$, $\dots$, $L_r$.
We suppose $A$, as above, is taken so that the linearized line bundles are restrictions of linearized line bundles $M_1$, $\dots$, $M_r$ on $X_A$.
We apply relative divisorialification
to the map
\[
[X_A/G]\to B\G_m^r\times [\A^\ell/\G_m^\ell]
\]
that combines \eqref{eqn.toproductBGm} and \eqref{eqn.XAtoAlmodGml}.
A numerical quantity that measures the
relative stabilizers, if not identically zero, takes its maximal value on a locus $W$, smooth over $B\G_m^r\times [\A^\ell/\G_m^\ell]$, in particular, transverse to the normal crossing boundary.
By the equivariant isomorphism of
$X_{A/tA}$ with $D\times \Spec(A/tA)$ and the analysis of the numerical quantity in terms of representation theory (cf.\  proof of Proposition \ref{prop.divisorial}),
every component of $W$ is either contained in $X_{A[t^{-1}]}$ or meets every fiber of $X_{A/tA}$ nontrivially.
The image in $\Spec(A)$ of a component of $W$ cannot contain the generic point, but by transversality cannot be equal to $\Spec(A/tA)$.
Thus $W$ is contained in $X_{A[t^{-1}]}$, with $W\times_{X_{A[t^{-1}]}}X=\emptyset$.

The proof of Theorem \ref{thm.revisitThm6.6} will be helped by a more convenient description of $\omega_I$.
The monomial map
\[ [\A^\ell/\G_m^\ell]\to [\A^1/\G_m] \]
with exponents $d_1$, $\dots$, $d_\ell$,
when composed with \eqref{eqn.XAtoAlmodGml},
is the map given by the boundary divisor $X_{A/tA}\subset X_A$.
The chosen uniformizer $t$ supplies a lift of the composite map to $\A^1$ and a corresponding lift of \eqref{eqn.XAtoAlmodGml} to
\begin{equation}
\label{eqn.XAtoAlmodK}
[X_A/G]\to [\A^\ell/\cK],\qquad
\cK:=\ker(\G_m^\ell\to \G_m).
\end{equation}
Away from the boundary, \eqref{eqn.XAtoAlmodK} restricts to
\[ [X_{A[t^{-1}]}/G]\to [(\A^1\setminus \{0\})^\ell/\cK]\cong \A^1\setminus \{0\}, \]
given by $t$.
The morphism \eqref{eqn.XAtoAlmodK} is smooth.
Since formation of the normal bundle commutes with flat base change,
we may describe $\mathcal{N}_{D_I/\cX}$ by base change from $[\A^\ell/\cK]$,
where we have $D_i$ defined by $x_i=0$, and
$D_I$, by $x_i=0$ for all $i\in I$.
Then $\mathcal{N}_{D_I/\A^\ell}\cong \A^\ell$, but now $x_i$, for $i\in I$, is a normal bundle coordinate;
$\omega_I$ is given by $x_1^{d_1}\dots x_\ell^{d_\ell}$.
Summarizing, with $0_I\colon \A^\ell\to \A^\ell$ given by $(x_i)_{i\in \cI}\mapsto (\mathbbm{1}_{\cI\setminus I}(i)x_i)_{i\in \cI}$,
we have the diagram with fiber square
\begin{equation}
\begin{split}
\label{eqn.NDIX}
\xymatrix@C=36pt{
\mathcal{N}_{D_I/\cX} \ar[r]\ar[d] & [\A^\ell/\cK] \ar[r]^(0.58){x_1^{d_1}\cdots x_\ell^{d_\ell}} \ar[d]^{0_I} & \A^1 \\
\cX \ar[r] & [\A^\ell/\cK] \ar[r]^(0.58){x_1^{d_1}\cdots x_\ell^{d_\ell}} & \A^1
}
\end{split}
\end{equation}
where the composite top map is $\omega_I$, and the composite bottom map is $t$.

In the proof of Theorem \ref{thm.revisitThm6.6} we suppress the explicit mention of $X_A$, but its use is implicit in application of facts, valid for smooth morphisms, to a morphism such as $[\cX/G]\to [\A^\ell/\cK]$, where the justification will use the smooth morphism \eqref{eqn.XAtoAlmodK}.

\begin{proof}[Proof of Theorem \ref{thm.revisitThm6.6}]
We use the description of $\mathcal{N}_{D_I/\cX}$, and $\omega_I$, from the diagram \eqref{eqn.NDIX}.
Restriction to $\A^1\setminus \{0\}$ in the top row has the effect of replacing
$\mathcal{N}_{D_I/\cX}$ by $\mathcal{N}_{D_I/\cX}^\circ$ and making the right-hand map an isomorphism.
Stabilizer components in a smooth family are smooth, so
\begin{align}
\begin{split}
\label{eqn.EequalsF}
\sum_{E\in \cS(\mathcal{N}^\circ_{D_I/\cX},G)/G}(&H,Y\actsfromleft k(E\cap \omega_I^{-1}(1)),\beta_E(\mathcal{N}^\circ_{D_I/\cX})) \\
&=\sum_{F\in \cS(\omega_I^{-1}(1),G)/G}(H,Y\actsfromleft k(F),\beta_F(\omega_I^{-1}(1))).
\end{split}
\end{align}

We bridge the gap between $\rho_t^G([X\actsfromright G])$, a sum over
$\cS(X,G)/G$, and the sum over
$\cS(\mathcal{N}^\circ_{D_I/\cX},G)/G$ in \eqref{eqn.EequalsF}, by means of deformation to the normal bundle.
Classically, for a smooth scheme $X$ with a smooth closed subscheme $V$, the construction
\[ \cM_{V/X}^\circ:=B\ell_{V\times\{0\}}(X\times \A^1)\setminus B\ell_{V\times\{0\}}(X\times \{0\}) \]
yields a scheme with morphisms to $X$ and to $\A^1$, such that:
\begin{itemize}
\item The morphism to $\A^1$ is smooth.
\item The fiber over $0$ is $\mathcal{N}_{V/X}$, projecting to $V\subset X$.
\item The fiber over $\A^1\setminus \{0\}$ is $X\times (\A^1\setminus \{0\})$.
\end{itemize}
(Usually $V$ and $X$ are not assumed smooth;
the general case is called deformation to the normal cone.
In the first property ``smooth'' is replaced by ``flat'', and in the second, instead of the normal bundle we have the normal cone.)
We are interested in $D_I$ in $\cX$, but by flatness $\cM_{D_I/\cX}^\circ$ is obtained by base change from $D_I$ in $[\A^\ell/\cK]$.
In the latter case the deformation to the normal bundle yields $[\A^\ell/\cK]\times \A^1$.
Denoting by $u$ the coordinate on the factor $\A^1$, we have morphisms $((x_i)_{i\in \cI},u)\mapsto (u^{\mathbbm{1}_I(i)}x_i)_{i\in \cI}$ to $[\A^\ell/\cK]$ and projection to $\A^1$.

There is the additional morphism $((x_i)_{i\in \cI},u)\mapsto x_1^{d_1}\cdots x_\ell^{d_\ell}$ to $\A^1$.
The additional morphism restricts to $\omega_I$ on the fiber over $0$; we denote it by $\tilde\omega_I$.
The morphism $\tilde\omega_I$ is flat, and the restriction over $\A^1\setminus \{0\}$ is smooth.
Upon restriction we get smooth $(\tilde\omega_I,\mathrm{pr}_2)$ to $(\A^1\setminus \{0\})\times \A^1$, and
subject to the caveat mentioned just before the start of the proof, the same can be asserted
when we map the corresponding open subscheme
of $\cM_{D_I/\cX}^\circ$
to $(\A^1\setminus \{0\})\times \A^1$.
The fiber over $(\A^1\setminus \{0\})\times (\A^1\setminus \{0\})$ is isomorphic to $X\times (\A^1\setminus \{0\})$, and the fiber over $(\A^1\setminus \{0\})\times \{0\}$ is $\cN_{D_I/\cX}^\circ$.

Thus, we have a $G$-equivariant map $\cS(\mathcal{N}_{D_I/\cX}^\circ,G)\to \cS(X,G)$ 
and for $E\in \cS(\mathcal{N}_{D_I/\cX}^\circ,G)$ mapping to
$W\in \cS(X,G)$ we have
$\beta_E(\mathcal{N}_{D_I/\cX}^\circ)=\beta_W(X)$.
For $W\in \cS(X,G)$ the closure $\cW$ in $\cX$ is a regular model of $W$, with normal crossing special fiber $(D_1\cap \cW)\cup\dots\cup (D_\ell\cap \cW)$ and the same multiplicities $d_i$ as for the special fiber $D_1\cup\dots\cup D_\ell$ of $\cX$.
So
\begin{align*}
\rho^G_t&([X\actsfromright G])\\
&=\sum_{\emptyset\ne I\subseteq\cI}(-1)^{|I|-1}\sum_{W\in \cS(X,G)/G}
(H,Y\actsfromleft k(\cW\times_{\cX}\omega_I^{-1}(1)),\beta_W(X))\\
&=\sum_{\emptyset\ne I\subseteq\cI}(-1)^{|I|-1}\sum_{E\in\cS(\mathcal{N}_{D_I/\cX}^\circ,G)/G}(H,Y\actsfromleft k(E\cap \omega_I^{-1}(1)),\beta_E(\mathcal{N}_{D_I/\cX}^\circ)),
\end{align*}
and we conclude by \eqref{eqn.EequalsF}.
\end{proof}




\section{Applications}
\label{sect:apply}

We suppose that $k$ is algebraically closed. 
Our main interest is the study of finite subgroups of the Cremona group $\Cr_n$, i.e., automorphisms of rational varieties, see Section~\ref{sect:basic1}.
Representative examples of rational varieties include:
\begin{itemize}
\item quadric hypersurfaces and quadric bundles over $\bP^1$, 
\item del Pezzo surfaces and rational Fano threefolds,  
\item rational conic bundles over rational surfaces, 
\item singular cubic hypersurfaces (other than cones),
\item toric varieties,
\item equivariant compactifications of linear algebraic groups, 
\item generalized flag varieties. 
\end{itemize}
One of the main problems is the {\em linearization problem}, i.e., the determination whether or not a given action is equivariantly birational to the projectivization of a representation. This is settled in dimension 2 \cite{pinardin}, but is largely open in dimensions $\ge 3$. 
%For example, actions of cyclic groups on smooth quadrics are linearizable.
%but we do not know whether or not all actions 
%of the symmetric group $\fS_3$ on quadric threefolds are linearizable. 

Another important problem is to prove that 
a given $G$-action on a rational variety is 
{\em stably linearizable}. The main tool for this is Lemma~\ref{lemm:non-name} (No-name lemma), applied to various vector bundle constructions. This is particularly interesting when the action is not linearizable. 



In this section, we discuss several geometric applications 
of Burnside groups, quotient stacks, and cohomological invariants, with special regard to
\begin{itemize}
\item stabilizer stratification, Amitsur and Brauer groups,
\item filtrations, incompressibles,
\item specialization.
\end{itemize}




\subsection*{Fields of invariants}
An application that we have already seen, of Lemma ~\ref{lemm:non-name} (No-name lemma), is the stable $G$-birationality of projectivized representations from Corollary \ref{coro:no-name}.
By essentially the same argument, the quotients $\bP(V)/G$ and $\bP(W)/G$ are stably birational.
However, the corresponding questions for ($G$-)birationality are more subtle.

First examples of nonbirational linear $G$-actions appeared in \cite{RYinvariant}; these were based on 
Condition {\bf (Det)} in Section~\ref{sect:det}. Further examples, via the Burnside group formalism, can be found 
in \cite{CTZ-p2}   and \cite{tschinkelyangzhang}.

We are not aware of examples of 
nonbirational {\em quotients} $\bP(V)/G$, $\bP(W)/G$, for projectivized representations of the same dimension, where the $G$-actions are generically free.
The same questions can be asked for nonlinear actions, e.g., actions on quadrics. 

\subsection*{Translation actions}

Let $\mathsf G$ be a connected linear algebraic group over $k$; as a variety, $\mathsf G$ is rational. Let $G\subset \mathsf G$ be a finite subgroup. 
When is the natural translation action of $G$ on $\mathsf G$ (stably) linearizable?

The action is stably linearizable, when $\mathsf G$ is {\em special}, i.e., 
$$
\rH^1(K,\mathsf G)=1,\quad \forall\, K/k; 
$$
see, e.g., \cite[Prop. 4.1]{HT-quad}. 

Here we show how to address the stable linearizability problem for the {\em nonspecial} group $$
\mathsf G=\PGL_{n}, \quad n\ge 2.
$$
Let $V$ be a faithful representation of $G$,
and $K:=k(V)^G$ the field of invariants.
We have, by functoriality:
\begin{equation}
\label{eqn.toH1KsfG}
\rH^1(BG,\mathsf G) \to \rH^1([V/G], \mathsf G)\to \rH^1(K,\mathsf G).
\end{equation}
The inclusion of $G$ in $\mathsf G$ supplies an element $\alpha\in \rH^1(BG,\mathsf G)$.

\begin{lemm}
\label{lem.indepV}
If $\alpha$ has trivial image in $\rH^1(K,\mathsf G)$, then $\alpha$ also has trivial image in $\rH^1(k(W)^G,\mathsf G)$ for any other faithful representation $W$ of $G$.
\end{lemm}

\begin{proof}
Let $L:=k(W)^G$, and let $\beta$ denote the
image of $\alpha$ in $\rH^1(L,\mathsf G)$.
By the No-name lemma, we have on a suitable invariant open $W^\circ\subset W$ an equivariant identification $V\times W^\circ\cong \A^d\times W^\circ$, $d=\dim(V)$.
Since $\alpha$ is trivial in $\rH^1(K,\mathsf G)$ it is also trivial on a dense open subset of $[V/G]$, and this determines a dense open subset of $\A^d_L$ on which $\beta$ becomes trivial.
Since $L$-points in $\A^d_L$ are Zariski dense, we obtain the vanishing of $\beta$.
\end{proof}

We define
\[ \alpha^{\mathsf G}(G,V)\in \rH^1(k(V)^G,\mathsf G) \]
to be the image of $\alpha$ under \eqref{eqn.toH1KsfG}.
By Lemma \ref{lem.indepV},
the triviality of $\alpha^{\mathsf G}(G,V)$ is independent of the choice of faithful representation $V$.

The argument in the proof of \cite[Prop.\ 4.1]{HT-quad} implies: 

\begin{prop}
\label{prop.asinHTquad}
If $\alpha^{\mathsf G}(G,V)$ is trivial, then $\mathsf G$ is $G$-equivariantly stably linearizable.
\end{prop}

In fact, the triviality of $\alpha^{\mathsf G}(G,V)$ is both necessary and sufficient for the stable linearizability of the translation action, and indeed is encoded in the classical obstruction to lifting $G$ to a \emph{linear} representation
\[ G\to \GL_n. \]
There is a natural $\PGL_n$-equivariant compactification
$$
\PGL_n\subset \bP^{n^2-1},
$$
the projectivization of the coordinates of an $n\times n$ matrix. 
Those with a single nonzero column form an invariant linear subspace
$\bP^{n-1}\subset \bP^{n^2-1}$.
Since the inclusion induces an isomorphism on Picard groups, by \cite[Lemma 2.1]{KT-quotient}
we have $\Am^2(\bP^{n-1},G)=\Am^2(\bP^{n^2-1},G)$.
But $\Am^2(\bP^{n-1},G)$ is cyclic, generated by a class $\gamma\in \rH^2(G,k^\times)$ that represents the obstruction to lifting $G\to \PGL_n$
to a linear representation.
If $\gamma=0$,
then such a lift exists, and
by functoriality $\alpha^{\mathsf G}(G,V)$ lies in the image of
\[ \rH^1(K,\GL_n)\to \rH^1(K,\PGL_n), \]
hence is trivial, and
by Proposition \ref{prop.asinHTquad}, $\PGL_n$ is $G$-equivariantly stably linear.
If $\gamma\ne 0$, then we have nontrivial $\Am^2(\bP^{n^2-1},G)$, and this obstructs not only
the $G$-equivariant stable linearizability of $\PGL_n$ but also the $G$-equivariant unirationality.

\subsection*{Condition (A)}
Actions of abelian groups on del Pezzo surfaces have been classified in \cite{blanc-thesis}, in particular, there is a classification of all actions satisfying Condition {\bf (A)}. Such a classification for smooth Fano threefolds can be found in \cite{abban}. 


\subsection*{Amitsur and Brauer groups}

The Amitsur invariant $\Am^2(X,G)$ from \cite[Sect.\ 6]{blancfinite} has been computed for rational surfaces \cite[Prop. 6.7]{blancfinite} and can be extracted for smooth Fano threefolds
from the analysis in  \cite{sharma}. 
Computations of $\Am^3(X,G)$ for del Pezzo surfaces can be found in \cite{TZ-uni}.
The related Condition {\bf (T)} (Section \ref{sect:amitsur}) has been investigated for Fano threefolds and toric varieties in \cite{KT-uni}, \cite{TZ-toric}. It would be desirable to undertake a systematic study of $\Am^4(X,G)$. 

An algorithm for computing Brauer groups of stacks $[X/G]$ and smooth projective models of quotient spaces $X/G$ has been presented in \cite{KT-quotient}. 






\subsection*{Incompressibles}


The description of incompressible symbols is easy in $\Burn_2(G)$.
There, the incompressible symbols are the divisor symbols 
$(H,Y\actsfromleft k(C),(b))$,
where $C$ is either a curve of positive genus, or $C\cong \bP^1$ is acted on by a \emph{noncyclic} group $Y$.
Already the simplest example
$$
(C_2, \fS_3\actsfromleft k(\bP^1),(1)), 
$$
is useful, e.g., reproving a result of Iskovskikh \cite{isk-s3} about nonbirationality of a linear action and a toric action of $G=C_2\times \fS_3$, see \cite[\S 7.6]{HKTsmall}.


For $\Burn_3(G)$ the situation is more involved.
There are many examples of incompressible symbols, but a full classification is unavailable; it would require the detailed analysis of \cite{DI}. 
Proofs of nonlinerizability of $G$-actions on rational varieties of dimension $\ge 3$, based on incompressibles, can be found for 
\begin{itemize}
\item actions on singular cubic threefolds, e.g., \cite[Exa.\ 2.7]{CTZ-forum}, or
\item the Burkhardt quartic, \cite[Prop.\ 7.2]{CTZ-burk}.
\end{itemize}
Among further applications of incompressible symbols to proofs of nonbirationality of stably birational actions are:
\begin{itemize}
\item on $\bP^3$, see \cite[Thm.\ 11.2]{KT-vector}, or \cite[Exa.\ 8.1]{TYZ-survey}, 
\item 
cubic threefolds and degree 14 Fano threefolds \cite{TZ-14}.
\end{itemize}

 
\subsection*{Specialization}

An immediate application of Theorem~\ref{thm.revisitThm6.6} is the fact very general members of families of $G$-varieties equivariantly specializing to (mildly singular) nonlinearizable $G$-varieties are also not linearizable \cite[Cor.\ 6.8, Exa.\ 6.9]{BnG}.

This motivated a search for such families. For example, a systematic study of specialization patterns for singular cubic threefolds to cubics with cohomological obstructions to stable linearizability was offered in \cite{CTZ-forum}, \cite{CMTZ}. In particular, there is now a classification of actions on singular cubics $X\subset\bP^4$ failing Condition {\bf (H1)};  these come from specific $C_2$ or $C_3$-actions, see \cite[Sect.\ 2]{CTZ-real}. A representative example of such a specialization is given in \cite[Prop.~4.3]{CTZ-forum}.







\section{What if \dots}
\label{sect:scifi}


\subsection*{Nonabelian stabilizers}
\label{ss.nonabelian}
{\em A priori}, one could try to encode a given regular action of a finite group $G$ by the stabilizer stratification and the induced action in the normal 
bundles, which would decompose into irreducible representations of the stabilizer groups. There are at least two reasons for bypassing this, via passage to divisorial forms: 
\begin{itemize}
\item one would have to keep track of these representations, and it is easier to deal with sums of one-dimensional representations,
\item encoding the interaction of the action of the stabilizers with the residual action on the stratum might involve nonabelian second cohomology \cite{giraud}, which we wanted to avoid. 
\end{itemize}
This explains the importance of achieving abelian stabilizers.



\subsection*{Not in divisorial form}
\label{ss.notdivisorial}
The definition of divisorial form ensures that the formula for the class in the equivariant Burnside group yields a sum of symbols.

\begin{exam}
\label{exa.D8action}
Consider the action of $\mathfrak{D}_4$ on $\bP^3_{\C}$, the projective compactification of the action of Example \ref{exa.notcentral}.
The subvariety $F$ defined by $x=y=0$ has generic stabilizer $H$, cyclic of order $4$, and nontrivial residual action of $\mathfrak{D}_4/H$ on $\C(F)$.
As observed in \cite[Exa.\ 3.5]{BnG}, primitive characters of $H$ do not lift to $\rH^1(\mathfrak{D}_4,\C(F)^\times)$.
\end{exam}

Such lifts of characters are needed in the definition of relation \textbf{(B)}, which is crucial for the equivariant Burnside group.
The definition of symbols (\S \ref{ss.symbols}) restricts to actions $Y\actsfromleft K$, where $Y$ is quotient by $H$ of a subgroup $Y_0$ of $G$, \emph{in which $H$ is central}.
Only then is there the possibility for all characters of $H$ to lift to $\rH^1(Y_0,K^\times)$
(see Proposition \ref{prop.divisorialcentral} and \cite[Lemma 2.1]{KT-struct}).



\subsection*{Tori}
\label{ss.tori}
We suppose that $k$ is algebraically closed.
In principle, the theory developed above can be extended to the case of an algebraic torus $T=\bG_m^d$. 
We can imitate the construction of the equivariant Burnside group and define analogous groups, where the symbols involve subgroups, quotient groups, and sequences of characters. 

For example, there would be an analogous group 
$$
\cB_n(T),
$$
generated by unordered sequences of characters 
$$
\beta = (b_1,\ldots, b_n), \quad b_j \in M = \mathfrak X^*(T),
$$
and subject to the blow-up relation as in Section~\ref{ss.symbols}. 

The first issue is that 
$\cB_n(T)$ has infinitely many generators and relations, so that it is not obvious that one can effectively check whether or not two such symbols are equal in $\cB_n(T)$. 

However, there are geometric reasons for not pursuing the formalism for $\Burn_n(T)$. Indeed, $T$-torsors are Zariski locally trivial, and thus a $T$-action on $X$ is equivariantly birational to a $T$-action on $T\times Y$, with $Y=X/T$, standard action of $T$ on itself, 
and trivial action of $T$ on $Y$. Moreover, actions of $T$ on $X$ and $X'$ are equivariantly birational if and only if the quotients $Y$ and $Y'$ are (non-equivariantly) birational. 

\subsection*{Extensions of finite groups by tori}
Another direction would be to consider algebraic groups fitting into an exact sequence
$$
1\to T\to N\to G\to 1, 
$$
where $G$ is a finite group and $T=\bG_m^d$. 
%This can yield families of actions of infinite groups specializing to actions of finite groups.  
%We would like to develop a specialization theory in this context. 

\begin{exam}
\label{exam:cubic}
Consider the singular cubic threefold $X\subset \bP^4$ given by
$$
x_1x_2x_3+x_3^3+x_3x_4x_5+x_4^3+x_5^3=0, 
$$
with singularities at $[1:0:0:0:0]$ and $[0:1:0:0:0]$, and the action of 
$$
1\to \bG_m\to N\to C_2\to 1, 
$$
given by rescaling 
$$
\eta_a\colon x_1\mapsto a x_1, \quad x_2\mapsto a^{-1}x_2
$$ 
and $\tau\colon x_1 \leftrightarrow x_2$, switching the coordinates.  
The fixed locus  is the smooth curve given by $x_1=x_2=0$, of genus 1.  
%The generic stabilizer of this curve is nonabelian. 
We can consider $\langle \eta_a, \tau\rangle$, $a\in k^\times$, defining a family of actions of (infinite, for general $a$) dihedral groups on $X$. When $a$ is a primitive even root of unity, the action is not linearizable, by \cite[Prop.\ 5.17]{CMTZ}.
\end{exam}





\subsection*{Reductive groups}
\label{ss.reductive}
One might inquire, whether the theory developed above could be extended to linear algebraic groups, more general than tori.
Reichstein and Youssin have established the existence of a standard form, after suitable equivariant blow-up \cite{RYessential}.
However, several issues arise:
\begin{itemize}
\item The stabilizer groups are extensions by unipotent groups of diagonalizable groups, so their representation theory is not determined by characters.
\item The stabilizer groups can \emph{vary} in a stratum.
\end{itemize}

\begin{exam}
\label{exa.GL2}
Consider the action of $\GL_2$ by right multiplication on the affine space of $2\times 2$-matrices.
The stabilizer groups of
rank $1$ matrices consist of a whole conjugacy class of $\G_a$-extensions of $\G_m$.
As this stratum has codimension $1$,
there is no possibility to improve the situation by blowing up.
\end{exam}




\bibliographystyle{plain}
\bibliography{burnsurv}

\end{document} 





The class of an $n$-dimensional $G$-variety in this group 
is computed on an appropriate smooth $G$-birational model $X$, 
called {\em standard form}: 
after a sequence of $G$-equivariant blow-ups we may assume that 
\cite{RYessential}: 
\begin{itemize}
\item there exists a Zariski open $U\subset X$ such that the $G$-action on $U$ is free, 
\item the complement $X\setminus U$ is a $G$-invariant simple normal crossing divisor, 
\item for every $g\in G$ and every irreducible component $D$ of $X\setminus U$, either $g(D)=D$
or $g(D)\cap D = \emptyset$. 
\end{itemize}
The standard form is preserved under $G$-equivariant blow-ups with smooth centers which have normal crossings with respect to the components of $D$.
Moreover, 
the stabilizer of every $x$ on such $X$ is an {\em abelian} subgroup of $G$ \cite[Thm. 4.1]{reichsteinyoussinessential}.
On such a model, the class of $X\actsfromright G$ is defined by: 
\begin{equation}
\label{eqn:class}
[X\actsfromright G] :=\sum_{H\subseteq G} \sum_F \fs_F \,
%(H,N_G(H)/H\actsfromleft K,\beta)
 \in \Burn_n(G),
\end{equation}
with summation over conjugacy classes of {\em abelian} subgroups   $H\subseteq G$ and strata $F\subseteq X$ with 
{\em generic} stabilizer $H$; the symbol 
\begin{equation}
\label{eqn:symb}
\fs_F:=(H,N_G(H)/H\actsfromleft k(F),\beta_F(X))
\end{equation}
records the action of the normalizer $N_G(H)$ of $H$ 
on $k(F)$, the product of the function fields of the components of $F$, as well as the generic normal bundle representation $\beta_F(X)$ of $H$. 
The class in \eqref{eqn:class} 
takes values in a quotient of the free abelian group generated by such symbols, 
by certain {\em blow-up} relations, spelled out in \cite[Definition 4.2]{BnG}, 
and ensuring that this expression is a well-defined 
$G$-birational invariant \cite[Theorem 5.1]{BnG}.

In \cite{HKTsmall}, we presented first geometric applications of this invariant. 
Here, we continue to explore {\em functorial} and {\em combinatorial} properties of $\Burn_n(G)$. 
We introduce and study:
\begin{itemize}
\item filtrations on $\Burn_n(G)$,  
\item the restriction homomorphism 
$$
\Burn_n(G)\to \Burn_n(G'),
$$
where $G'\subset G$ is any subgroup,
\item products, 
\item 
a combinatorial analog $\BC_n(G)$ of $\Burn_n(G)$,  obtained by  
forgetting {\em field-theoretic} information, while keeping only discrete invariants encoded in a symbol \eqref{eqn:symb}. 
\end{itemize}
One of the motivating problems in this field is to distinguish equivariant birational types of (projectivizations of) linear actions (see, e.g., \cite{vinberg}, \cite{katsylo}). A sample question, raised in \cite[Section 8]{DolIsk}, is: {\it Are there 
isomorphic finite subgroups of $\mathrm{PGL}_3$ which are not conjugate in the plane Cremona group?} 

Examples of equivariantly nonbirational representations 
considered in \cite{reichsteinyoussininvariant}
required that $G$ contains an abelian $p$-subgroup of rank
equal to the dimension of the representation. 
Our formalism yields new examples 
%in all dimensions $\ge 2$, 
without the rank condition on the group, see Example~\ref{exa.Z5S3}.























\section{Generalities}
\label{sect:gen}

We adopt notational conventions from \cite{BnG}: 
\begin{itemize}
\item $G$ is a finite group,
\item $\mathrm{Ab}(G)$ is the set of abelian subgroups $H\subseteq G$, with character groups 
$$
H^\vee:=\Hom(H,k^\times),
$$
and 
$$
\mathrm{Ab}(G)^c = \mathrm{Ab}(G)/G,
$$
the set of conjugacy classes of abelian subgroups,
\item $\Bir_d(k)$ is the set of birational equivalence classes of $d$-dimensional algebraic varieties over $k$ , i.e., 
the set of finitely generated fields of transcendence degree $d$ over $k$; we identify a field with its isomorphism class in $\Bir_d(k)$, 
\item $\Alg_N(K_0)$ is the set of isomorphism classes of Galois algebras $K$ over $K_0\in \Bir_d(k)$
for the group
\[ N:=N_G(H)/H, \]
satisfying \\
\


{\bf Assumption 1}:  the composite homomorphism
\begin{equation}
\label{eqn.weak}
\rH^1(N_G(H),K^\times)\to \rH^1(H,K^\times)^N\to H^\vee
\end{equation}
is surjective (see \cite[Section 2]{BnG} for more details).
\item More generally, for a subgroup $M\subset N$ we denote by
$\Alg_M(K_0)$ the set of isomorphism classes of $M$-Galois algebras $K/K_0$
(i.e., Galois algebras $K$ over $K_0$ for the group $M$), such that 
$\Ind_M^N(K)$ 
satisfies
Assumption 1. Of particular interest is 
\[ 
Z:=Z_G(H)/H\subseteq N=N_G(H)/H.
\]
\end{itemize}



\begin{lemm}
\label{lem.Z}
Let $K\in \Alg_N(K_0)$. 
Then
$$
K\cong \Ind_{Z}^{N}(K')
$$ 
for some $K'\in \Alg_Z(K_0)$.
\end{lemm}

\begin{proof}
With notation of Assumption 1, we have
trivial $H$-action on 
$K^\times$, so 
$$
\mathrm{H}^1(H,K^\times)=\Hom(H,K^\times).
$$
Writing $K\cong K^1\times\dots\times  K^\ell$, where each $K^i$ is a field,
as rightmost map in \eqref{eqn.weak} we take
$$
\mathrm{H}^1(H,K^\times)\to\mathrm{H}^1(H,(K^1)^\times)\cong H^\vee.
$$
Projection $K^\times\to (K^1)^\times$ is equivariant for the
subgroup $Y\subseteq N$, defined by the
condition of sending $K^1$ to $K^1$, where the action on
$\mathrm{H}^1(H,(K^1)^\times)$ is given just by conjugation on $H$.
Assumption 1 implies that the conjugation action is trivial,
i.e., $Y\subseteq Z$.
So the result holds with 
$$
K'=\Ind_Y^Z(K^1).
$$
\end{proof}

\begin{rema}
\label{rem.groupH2}
Assumption 1, for an $N$-Galois algebra $K/K_0$, of the form
$\Ind_Z^N(K')$, where $K'/K_0$ is a
$Z$-Galois algebra, may be expressed as the
surjectivity of
\begin{equation}
\rH^1(Z_G(H),K'^\times)\to H^\vee
\label{eqn.groupH2}.
\end{equation}
Given this, the proof of
\cite[Prop.\ 2.2]{BnG} supplies an
equivalence of categories between
\begin{itemize}
\item 
$H$-Galois algebras over \'etale $K_0$-algebras and
\item 
$Z_G(H)$-Galois algebras with
equivariant homomorphism from $K'$;
\end{itemize}
in particular, there is then a $Z_G(H)$-Galois algebra $L/K_0$ with
homomorphism $K'\to L$,
compatible with the structure of Galois algebra for the
group $Z$, respectively $Z_G(H)$.
Assumption 1 is also implied by the existence
of such a Galois algebra $L$
and homomorphism $K'\to L$, as we see by using
the Hochschild-Serre spectral sequence and Hilbert's Theorem 90 to
identify $\rH^1(Z_G(H),K'^\times)$ with
$\rH^1(H\times Z_G(H),L^\times)$.
This allows us to view Assumption 1 as a lifting problem of Galois cohomology
\[ \rH^1(\Gal_{K_0},Z_G(H))\to \rH^1(\Gal_{K_0},Z) \]
and remark that the machinery of nonabelian cohomology
(cf.\ \cite[\S 1.3.2]{platonovrapinchuk})
supplies an obstruction to lifting in
$\rH^2(\Gal_{K_0},H)$.
\end{rema}




We now recall the definition of the {\em equivariant Burnside group}
$$
\Burn_n(G)=\Burn_{n,k}(G)
$$
following \cite[Section 4]{BnG}: it is a $\bZ$-module, generated by symbols
$$
\mathfrak s:=(H, N \actsfromleft K, \beta),
$$
where
\begin{itemize}
\item $H\subseteq G$ is an abelian subgroup,
\item $K\in \Alg_N(K_0)$, with $K_0\in \Bir_d(k)$, and $d\le n$,  
\item $\beta=(b_1,\dots,b_{n-d})$, a sequence of nonzero elements of the
character group $H^\vee$, that generate $H^\vee$.
\end{itemize}
The sequence of characters $\beta$ determines a faithful
$(n-d)$-dimensional representation of $H$ over $k$, with
trivial space of invariants.
As every $(n-d)$-dimensional representation of $H$ over $k$ splits as
a sum of one-dimensional representations, any
faithful $(n-d)$-dimensional representation of $H$ over $k$
determines a sequence of characters, generating $H^\vee$,
up to order.
The ambiguity of order gives us the first of several relations that we
impose on symbols:

\medskip
\noindent
\textbf{(O):}
$(H, N \actsfromleft K, \beta)=(H, N \actsfromleft K, \beta')$ if
$\beta'$ is a reordering of $\beta$.

\medskip
\noindent
The further relations are
\textbf{conjugation} and \textbf{blow-up} relations:

\medskip
\noindent
\textbf{(C):} 
$(H,N \actsfromleft K,\beta) = (H',N' \actsfromleft K,\beta')$,
when $H'=gHg^{-1}$ and $N'=N_G(H')/H'$,  with $g\in G$, and $\beta$ and $\beta'$ are related by conjugation by $g$.

\medskip
\noindent
\textbf{(B1):}
$(H,N \actsfromleft K, \beta)=0$ when $b_1+b_2=0$.

\medskip
\noindent
\textbf{(B2):}
$(H,N \actsfromleft K, \beta)=\Theta_1+\Theta_2$,
where
\[
\Theta_1= \begin{cases} 0, & \text{if $b_1=b_2$}, \\ 
(H,N \actsfromleft K, \beta_1)+(H,N \actsfromleft K, \beta_2),
& \text{otherwise},
\end{cases}
\]
with
\begin{equation}
\label{eqn:beta12}
\beta_1:=(b_1,b_2-b_1,b_3,\ldots, b_{n-d}), \quad \beta_2:=(b_2,b_1-b_2,b_3,\ldots, b_{n-d}),
\end{equation}
and
\[
\Theta_2= \begin{cases} 0, & \text{if $b_i\in \langle b_1-b_2\rangle$ for some $i$}, \\
(\overline{H},\overline{N} \actsfromleft \overline{K}, \bar{\beta}), & \text{otherwise}, 
\end{cases}
\]
with
\[
\overline{H}^\vee:=H^\vee/\langle b_1-b_2\rangle,
\quad 
\bar{\beta}:=(\bar b_2,\bar b_3,\dots,\bar b_{n-d}), \quad
\bar b_i\in \overline{H}^\vee,
\]
and $\overline{K}$ carries the action described in Construction \textbf{(A)}
in \cite[Section 2]{BnG}, applied to the character $b_1-b_2$.

\medskip

We permit ourselves to write a symbol in the form
\begin{equation}
\label{eqnwithM}
(H,M\actsfromleft K,\beta)
\end{equation}
with a subgroup $M\subset N$ and $K\in \Alg_M(K_0)$,
with
\[ (H,M\actsfromleft K,\beta):=(H,N\actsfromleft \Ind_M^N(K),\beta). \]
We further allow $K_0$ to be a \emph{product} of fields; 
then \eqref{eqnwithM} will denote the corresponding
sum of symbols, one for each factor.

By Lemma~\ref{lem.Z}, any symbol in $\Burn_n(G)$ is of the form
$$
(H,Z\actsfromleft K,\beta),
$$
with $K\in \Alg_Z(K_0)$.
In this notation, Construction \textbf{(A)} has a compact formulation. 
Applied to a single character $b\in H^\vee$,
this yields the subgroup
\[ \overline{H}:=\ker(b) \subset H\]
and the symbol
\[ (\overline{H}, Z_G(H)/\overline{H}\actsfromleft K(t),\bar\beta), \]
where a $Z_G(H)$-action on $K(t)$ arises by lifting $b$
via \eqref{eqn.groupH2} and is trivial on $\overline{H}$,
and $\bar\beta$ is obtained from
$\beta$ by applying the map $H^\vee\to \overline{H}^\vee$, as above.

\begin{rema}
\label{rema:iter}
Construction \textbf{(A)} may be applied to a collection of
characters, yielding the same outcome as when applied
iteratively, one character at a time.
\end{rema}

\medskip

A $G$-action on $X$ in standard form always satisfies

\medskip

{\bf Assumption 2}:  The stabilizers for the $G$-action on $X$ are abelian,
and for every $H$ and $F$ in \eqref{eqn:class} the
composite homomorphism
\[ \Pic^G(X)\to \rH^1(N_G(H),k(F)^\times)\to H^\vee \]
is surjective,
where the first map is given by restriction and the second is the
map from Assumption 1, with $K=k(F)$.

\

Note that Assumption 2 implies Assumption 1,
for every $H$ and every $N_G(H)/H\actsfromleft k(F)$
(see \cite[Rmk.\ 3.2(i)]{BnG}).

A variant, that will occur below, is the requirement of surjectivity, when we
restrict to a given subgroup of $\Pic^G(X)$.
Given this, we will say that Assumption 2 holds for the given subgroup
of $\Pic^G(X)$.


\section{Filtrations}
\label{sec.filtrations}
In this section, we explore additional combinatorial constructions on equivariant Burnside groups $\Burn_n(G)$, reflecting the geometry of the $G$-action on strata with given generic stabilizers. 




\begin{defi}
\label{defn:pre}
A $G$-prefilter 
is a collection $\mathbf{H}$ of pairs $(H,Y)$
consisting of an abelian subgroup
$H\subseteq G$ and a subgroup
$$
Y\subseteq Z=Z_G(H)/H,
$$ 
such that $\mathbf{H}$ is closed under conjugation,
i.e., for $(H,Y)\in \mathbf{H}$ we have
$$
(gHg^{-1},gYg^{-1})\in \mathbf{H},\quad \text{ for all } \, g\in G.
$$
\end{defi}

\begin{defi}
\label{defn.BurnH}
Given a $G$-prefilter $\mathbf{H}$, we let
$$
\Burn^{\mathbf{H}}_n(G)
$$ 
be the quotient of
$\Burn_n(G)$ by the subgroup generated by classes of the form
$$
(H, Y \actsfromleft K, \beta),
$$
where $K\in \Alg_Y(K_0)$ is a \emph{field}, and
\[ (H,Y)\notin \mathbf{H}. \]
\end{defi}

\begin{prop}
\label{prop.filtration}
Let $\mathbf{H}$  be a $G$-prefilter 
%be a collection of pairs $(H,Y)$ of subgroups,
%with $H$ abelian and $H\subseteq Y\subseteq Z_G(H)$, closed under
%conjugation, 
such that if $(H,Y)\in \mathbf{H}$, with $H$ nontrivial,
then $(\langle H,g\rangle,Y/\langle \bar g\rangle)\in \mathbf{H}$ for all
$g\in Z_G(H)$ satisfying
$$
\bar g\in Y\qquad\text{and}\qquad Y\subseteq Z_G(g)/H.
$$
Then $\Burn^{\mathbf{H}}_n(G)$ is generated by
triples
\[ (H, Y \actsfromleft K, \beta), \]
where $K\in \Alg_Y(K_0)$ is a field and
\[ (H,Y)\in \mathbf{H}, \]
subject to relations \emph{\textbf{(O)}}, \emph{\textbf{(C)}}, \emph{\textbf{(B1)}}, and \emph{\textbf{(B2)}} 
applied to these triples. 
%for
%$$
%(H, N \actsfromleft \Ind_{Y/H}^{N}(K'), \beta)
%$$ 
%with $K'\in \Alg_{Y/H}(K_0)$ a field and $(H,Y)\in \mathbf{H}$.
\end{prop}

\begin{proof}
For any
$$
(H, Y \actsfromleft K, \beta)
$$ 
with
$K\in \Alg_Y(K_0)$ a field,
the term $\Theta_2$ from \textbf{(B2)}, when nontrivial, consists of a subgroup
$\ker(b)$ of $H$ for some $b\in H^\vee$,
a field $K(t)$ with action of the
pre-image of $Y$ in $Z_G(H)/\ker(b)$,
and a sequence of characters.
If $(H,Y)\notin \mathbf{H}$, then by hypothesis the pair consisting of
$\ker(b)$ and the pre-image of $Y$
is not in 
$\mathbf{H}$.
This observation establishes the proposition,
since \textbf{(B1)} involves just one triple,
and in $\Theta_1$ from \textbf{(B2)} the group and algebra do not change.
\end{proof}

\begin{exam}
\label{exa.filtration}

\noindent
\begin{itemize}
\item For $G$ abelian, we have
$$
\Burn^G_n(G)=\Burn^{\{(G,\mathrm{triv})\}}_n(G),
$$
where the left side was introduced in \cite[\S 8]{BnG}:
this is the quotient of
$\Burn_n(G)$ by all triples whose first entry is a proper subgroup of $G$.
\item 
For $\mathbf{H}$ consisting of all $(H,Y)$ with $H$ nontrivial cyclic and
$Y$ noncyclic, and $k$ algebraically closed, 
$\Burn^{\mathbf{H}}_2(G)$ appeared in 
\cite[\S 7.4]{HKTsmall}. 
\end{itemize}
\end{exam}

\begin{rema}
\label{rema:alt}
One can additionally suppress the field information, which will lead to
{\em combinatorial} analogues of Burnside groups. We will explore this in Section~\ref{sec.combinatorial}.
\end{rema}



\section{Nontrivial generic stabilizers}
\label{sec.indexed}
In this section, we introduce a version of the equivariant Burnside group, relevant for considerations of actions with {\em nontrivial} generic stabilizer. 

Let $G$ be a finite group.
A variant of the equivariant Burnside group takes
the additional data of a finite index set
\[ I\subset \N. \]
The {\em equivariant indexed Burnside group}
$$
\Burn_{n,I}(G),
$$
is defined as a quotient of the $\bZ$-module generated by symbols
$$
(H\subseteq H', N' \actsfromleft K, \beta, \gamma),
$$
where
\begin{itemize}
\item $H\subseteq H'$ are abelian subgroups of $G$, 
\item $N':=N_{N_G(H)}(H')/H'$,
\item $K\in \Alg_{N'}(K_0)$, with $K_0\in \Bir_d(k)$, and $d\le n-|I|$,
\item $\beta=(b_1,\dots,b_{n-d-|I|})$, a sequence of nonzero
characters of $H'$,
trivial upon restriction to $H$,
that generate $(H'/H)^\vee$,
\item $\gamma=(c_i)_{i\in I}$ is a sequence of
elements of $H'^\vee$, such that
the images of $c_i$
in $H^\vee$ generate $H^\vee$.
\end{itemize}
As in Section \ref{sect:gen}, we permit ourselves to write a symbol in the form
$$
(H\subseteq H', M' \actsfromleft K, \beta, \gamma),
$$
where $M'\subset N'$ is a subgroup.
Every symbol may be expressed as
\[ (H\subseteq H',Z'\actsfromleft K,\beta,\gamma),\qquad
Z':=Z_G(H')/H'. \]
(Notice that $Z_G(H')=Z_{N_G(H)}(H')$.)

These symbols are subject to relations: 

\medskip
\noindent
\textbf{(O):}
$(H\subseteq H', N' \actsfromleft K, \beta,\gamma)=(H\subseteq H', N' \actsfromleft K, \beta',\gamma)$ if
$\beta'$ is a reordering of $\beta$. 

\medskip
\noindent
\textbf{(C):} 
$(H\subseteq H',N' \actsfromleft K,\beta,\gamma) = (gHg^{-1} \subseteq gH'g^{-1},gN'g^{-1} \actsfromleft K,\beta',\gamma')$
for $g\in G$, with $\beta$ and $\beta'$,
respectively $\gamma$ and $\gamma'$, related by conjugation by $g$.

\medskip
\noindent
\textbf{(B1):}
$(H\subseteq H',N' \actsfromleft K, \beta, \gamma)=0$ when $b_1+b_2=0$.

\medskip
\noindent
\textbf{(B2):}
$(H\subseteq H',N' \actsfromleft K, \beta, \gamma)=\Theta_1+\Theta_2$,
where $\Theta_1$ and $\Theta_2$ are as in Section \ref{sect:gen}, with $H$ prepended and
$\gamma$, respectively $\bar\gamma$,
appended to the corresponding symbols.

\begin{rema}
\label{rem.charactersgenerate}
By analogy with Remark \eqref{rem.groupH2},
we may express Assumption 1, for the
Galois algebra $K$, as the surjectivity of the
middle vertical map
\[
\xymatrix@C=8pt{
0 \ar[r] &
\rH^1(Z_G(H')/H,K^\times) \ar[r] \ar[d] &
\rH^1(Z_G(H'),K^\times) \ar[r] \ar[d] &
\rH^1(H,K^\times)^{Z_G(H')/H} \ar[d] \\
0 \ar[r] & (H'/H)^\vee \ar[r] & H'^\vee \ar[r] & H^\vee \ar[r] & 0
}
\]
Here, the top row comes from the Hochschild-Serre spectral sequence.
In a symbol, we have $\beta$ generating the left-hand group in the bottom row,
while $\gamma$ is a sequence of
characters of $H'$, whose images generate $H^\vee$.
Consequently, $\beta$ and $\gamma$
together generate $H'^\vee$.
Thus we have a homomorphism
\[
\psi_I\colon \Burn_{n,I}(G)\to \Burn_n(G),
\]
sending $(H\subseteq H',Z'\actsfromleft K,\beta,\gamma)$ to
\[
(H',Z'\actsfromleft K, \beta\cup\gamma)
\]
when $\gamma$ is a sequence of nontrivial characters,
otherwise to $0$.
\end{rema}

In order to explain the
relevance of this definition, we introduce a map which
converts some of the characters in $\gamma$ to a
transcendental extension of the Galois algebra.
Let
\[ J\subseteq I \]
be a subset.
Given a symbol $(H\subseteq H',Z'\actsfromleft K,\beta,\gamma)$,
we use $J$ to define subgroups
\begin{align*}
\overline{H}'&:=\bigcap_{i\in I\setminus J}\ker(c_i)\subseteq H', \\
\overline{H}&:=H\cap \overline{H}'\subseteq H.
\end{align*}
Then we define
\[ \omega_{I,J}\colon \Burn_{n,I}(G)\to \Burn_{n,J}(G), \]
by applying Construction \textbf{(A)} when possible:
\[
(H\subseteq H',Z'\actsfromleft K,\beta,\gamma)\mapsto
(\overline{H}\subseteq \overline{H}',
Z_G(H')/\overline{H}'\actsfromleft K((t_i)_{i\in I\setminus J}),\bar\beta,
\bar\gamma), 
\]
where $\bar\gamma=(\bar c_j)_{j\in J}$, when all of the characters of $\bar\beta$ are nonzero, and
$$
(H\subseteq H',Z'\actsfromleft K,\beta,\gamma)\mapsto 0, \quad \text{ otherwise.}
$$
This is compatible with relations:
the only one that is nontrivial to check is \textbf{(B2)},
where
$\Theta_1$ maps to
$\overline{\Theta}_1$, as we see by dividing into cases according to the vanishing of $\bar b_1$ or $\bar b_2$, or their equality, and
$\Theta_2$ maps to $\overline{\Theta}_2$, as we see using Remark~\ref{rema:iter}.

We recall the setting of
\cite[Defn.\ 5.4]{BnG}: Let 
$X$ be a smooth projective variety of dimension $n$, with a generically free action of $G$, satisfying Assumption 2. Let $D_1$, $\dots$, $D_{\ell}$ be $G$-stable divisors, 
with
$$
D_I:=\bigcap_{i\in I}D_i, \quad \text{
for } I\subseteq \cI:=\{1,\dots,\ell\}, \quad D_\emptyset=X.
$$
We suppose, for notational simplicity,
that for every $I$ the generic stabilizers of the
components of $D_I$
belong to a single conjugacy class of subgroups,
and take $H_I$ to be a representative.
Then
to $I\subseteq M\subseteq \cI$
we attach the following class in
$\Burn_{n,I}(G)$:
\[
\chi_{I,M}(X\actsfromright G,(D_i)_{i\in \cI}):=\sum_{H'\supseteq H_I}\sum_{\substack{W\subset D_I\\
\mathrm{generic\ stabilizer}\ H'\\
\{i\in \cI\,|\,W\subset D_i\}=M}}
\!\!\!(H_I\subseteq H',N'\actsfromleft k(W),\beta,\gamma),
\]
where 
\begin{itemize}
\item the first sum is over conjugacy class representatives $H'$ of abelian subgroups
of $N_G(H_I)$, containing $H_I$,
\item the second sum is over $N_{N_G(H_I)}(H')$-orbits of components
$W$ with generic stabilizer $H'$, contained in components of
$D_I$ with generic stabilizer $H_I$ and satisfying
$\{i\in \cI\,|\,W\subset D_i\}=M$,
\item $\beta=\beta_W(D_I)$ encodes the normal bundle to
$W$ in $D_I$, and
\item 
$\gamma=(c_i)_{i\in I}$, the characters coming from
$D_i$ with $i\in I$.
\end{itemize}

Then
\[
[\mathcal{N}_{D_I/X}\actsfromright G]^{\mathrm{naive}}=
\sum_{I\subseteq M\subseteq\cI}
\sum_{M\setminus I\subseteq J\subseteq M}
\!\psi_{I\cap J}(\omega_{I,I\cap J}(\chi_{I,M}(X\actsfromright G,(D_i)_{i\in \cI}))),
\]
where the terms with $J=\emptyset$ contribute
$[\mathcal{N}^{\circ}_{D_I/X}\actsfromright G]^{\mathrm{naive}}$.
This provides some insight to \cite[Lemma 5.7]{BnG}.



\section{Fibrations}
\label{sec.fibrations}
In this section, we define a projectivized version
of the equivariant indexed Burnside group and use it to give a formula for the class in $\Burn_n(G)$ of the projectivization of a sum of line bundles.   

Let $G$ be a finite group and $I\subset \N$ a \emph{nonempty} finite index set.
The \emph{equivariant projectively indexed Burnside group}
\[ \Burn_{n,\bP(I)}(G) \]
is defined with generators and relations as in
Section \ref{sec.indexed}, where
\begin{itemize}
\item $\beta$ consists of $n-d-|I|+1$ characters
(so $d\le n-|I|+1$),
\item 
the \emph{differences} of pairs of characters of 
$\gamma$ should generate $H^\vee$,
\item 
and there is an additional relation:
\end{itemize}

\medskip
\noindent
\textbf{(P):}
If
$\gamma'-\gamma$ is a constant sequence then
$$
(H\subseteq H', N' \actsfromleft K, \beta,\gamma)=(H\subseteq H', N' \actsfromleft K, \beta,\gamma').
$$


\

We define
\[ \omega_{\bP(I),J}\colon \Burn_{n,\bP(I)}(G)\to \Burn_{n,J}(G), \]
for a \emph{proper} subset
\[ J\subsetneq I, \]
by 
\begin{itemize}
\item 
choosing $i_0\in I\setminus J$,
\item 
applying \textbf{(P)} to get a representative symbol
$$
(H\subseteq H', N' \actsfromleft K, \beta,\gamma)
$$ 
with
$\gamma_{i_0}=0$, and
\item 
applying
$\omega_{I\setminus\{i_0\},J}$ to the class of
$(H\subseteq H', N' \actsfromleft K, \beta,(c_i)_{i\in I\setminus \{i_0\}})$ in $\Burn_{n,I\setminus \{i_0\}}(G)$.
\end{itemize}

Let $X$ be a smooth projective variety over $k$.
Assume that $X$ carries a $G$-action, and let
$L_0$, $\dots$, $L_r$ be $G$-linearized line bundles on $X$.
The next statement examines the condition, for
$G$ to act generically freely on $\bP(L_0\oplus\dots\oplus L_r)$, so
that Assumption 2 satisfied.

\begin{lemm}
\label{lem.assumptionbundles}
Let $X$ be a smooth projective variety over $k$ with a
$G$-action
and $G$-linearized line bundles $L_0$, $\dots$, $L_r$.
Let $H$ be the stabilizer at the generic point of a component of $X$,
and let us denote the $N_G(H)$-orbit of the component by $X'$.
The following are equivalent.
\begin{itemize}
\item[(i)] The $N$-action on $X'$ satisfies Assumption 2, and
$H$ is abelian with $H^\vee$ spanned by the differences of
characters determined by $L_0$, $\dots$, $L_r$.
\item[(ii)] The $G$-action on $\bP(L_0\oplus\dots\oplus L_r)$ is
generically free and satisfies Assumption 2.
\item[(iii)] The $G$-action on $\bP(L_0\oplus\dots\oplus L_r)$ is
generically free and satisfies Assumption 2 for
$L_0$, $\dots$, $L_r$, together with the
$G$-linearized line bundles on $X$ associated with
$N$-linearized line bundles on $X'$.
\end{itemize}
\end{lemm}

The statement is inspired by \cite[Lemma 7.3]{BnG}.

\begin{proof}
The action of $G$ on $\bP(L_0\oplus\dots\oplus L_r)$ is generically free
if and only if the
action of $N_G(H)$ on $\bP(L_0|_{X'}\oplus\dots\oplus L_r|_{X'})$ is generically free.
The latter has generic stabilizer $\bigcap_{i=1}^r\ker(b_i-b_0)$.
Thus the condition on $H$ in (i) is equivalent to the condition of
generically free action in (ii) and in (iii).
We assume this from now on.

An $N$-linearized line bundle on $X'$ determines an
$N_G(H)$-linearized line bundle on $X'$, with trivial $H$-action.
An $N_G(H)$-linearized line bundle on $X'$ determines, and is determined by,
$G$-linearized line bundle on $X$; this is the meaning of the
associated line bundles in (iii).
Conversely, a $G$-linearized line bundle on $X$ which restricts to an
$N_G(H)$-linearized line bundle on $X'$ with trivial $H$-action,
gives rise to an $N$-linearized line bundle on $X'$.

We start by showing (i) implies (iii), using the
interpretation of Assumption 2 in terms of the
representability of a certain morphism from the quotient stack to a
product of copies of $B\G_m$.
Given (i), we have such a representable morphism
\[ [X'/N]\to B\G_m\times\dots\times B\G_m. \]
Correspondingly, the fibers of the composite morphism
\[ [X/G]\to [X'/N]\to B\G_m\times\dots\times B\G_m \]
all have constant stabilizer group $H$.
The condition in (i) implies that the
$H$-representation given by $b_0$, $\dots$, $b_r$ is faithful.
With $r+1$ additional factors $B\G_m$ we get a
representable morphism from $[X/G]$, hence also from
$\bP(L_0\oplus\dots\oplus L_r)$.

Since trivially (iii) implies (ii), it remains only
to show (ii) implies (i).
Generally, a line bundle on a projective bundle is isomorphic to
the pullback of
a line bundle from the base, twisted by a power of the tautological line bundle.
Two vector bundles, one obtained from the other by tensoring
by a line bundle, have isomorphic projectivizations,
the tautological line bundle of one obtained from the other by
tensoring by the pullback of the line bundle from the base.
A sum of line bundles, after such tensoring, may be brought in a form with
trivial $i$th factor, for any $i$,
and this way see that any power of the tautological line bundle on
$\bP(L_0\oplus\dots\oplus L_r)$, restricted to
the open $U_i\subset \bP(L_0\oplus\dots\oplus L_r)$
defined by nonvanishing on the component $L_i$,
is identified with a line bundle pulled back from the base;
all of these assertions are valid in an equivariant setting.
With the notation of Assumption $2$ for $\bP(L_0\oplus\dots\oplus L_r)$,
we always have $\Spec(k(F))\subset U_i$ for some $i$.
So, (ii) implies that
the $G$-action on $\bP(L_0\oplus\dots\oplus L_r)$
satisfies Assumption 2 for $\Pic^G(X)$.
Since
\[ \bP(L_0\oplus\dots\oplus L_r)\to X \]
admits
equivariant sections, we deduce that some finite collection of
$G$-linearized line bundles on $X$ determines a representable morphism
\[ [X/G]\to B\G_m\times\dots\times B\G_m. \]
Replacing each by a tensor product with combinations of
$L_0$, $\dots$, $L_r$, we obtain a
$G$-linearized line bundle on $X$ that comes from an
$N$-linearized line bundle on $X'$.
The corresponding morphism
\[ [X'/N]\to B\G_m\times\dots\times B\G_m \]
is representable, and thus we have (i).
\end{proof}

\begin{prop}
\label{prop.fibration}
Let $X$ be a smooth projective variety of dimension $n-r$ over $k$ with a
$G$-action and $G$-linearized line bundles $L_0$, $\dots$, $L_r$.
We assume the conditions and adopt the notation of
Lemma \ref{lem.assumptionbundles}.
We define $I:=\{0,\dots,r\}$ and the following class in $\Burn_{n,\bP(I)}(G)$:
\[
\xi(X\actsfromright G,(L_i)_{i\in I}):=
\sum_{H'\supseteq H} \sum_{\substack{W\subset X'\\ \mathrm{generic\ stabilizer}\ H'}}
(H\subseteq H',N'\actsfromleft k(W),\beta,\gamma),
\]
where
\begin{itemize}
\item the first sum is over
abelian subgroups $H'$ of $G$ that contain $H$,
up to conjugacy in $N_G(H)$,
\item the second sum is over $N_{N_G(H)}(H')$-orbits of components $W\subset X'$ where the
generic stabilizer is $H'$,
\item $\beta=\beta_W(X')$ encodes the normal bundle to $W$ in $X'$, and
\item $\gamma=(c_i)_{i\in I}$, the characters coming from $L_i$ with $i\in I$.
\end{itemize}
Then
\[
[\bP(L_0\oplus\dots\oplus L_r) \actsfromright G]=\sum_{J\subsetneq I} \psi_J(\omega_{\bP(I),J}(\xi(X\actsfromright G, (L_i)_{i\in I})))
\]
in $\Burn_n(G)$.
\end{prop}

\begin{proof}
We identify each contribution to
$[\bP(L_0\oplus\dots\oplus L_r) \actsfromright G]$
as 
$$
V=\varphi_J^{-1}(W),
$$
for some $W$ in the
definition of $\xi(X\actsfromright G,(L_i)_{i\in I})$,
where $\varphi_J$ denotes the projection to $X$ from the projectivization of
$\bigoplus_{i\in I\setminus J}L_i$.
Then,
$$
(H\subseteq H',N'\actsfromleft k(W),\beta,\gamma)\in \Burn_{n,\bP(I)}(G)
$$
maps under $\psi_J\circ \omega_{\bP(I),J}$ to
$$
(\overline{H}',N_{N_G(H)}(\overline{H}')\actsfromleft k(V),\beta_V(X)).
$$
\end{proof}

\begin{exam}
\label{exa.Z5S3}
Let $G:=C_5\times \fS_3$, acting on
$X:=\bP^1$ via an irreducible $2$-dimensional representation of $\fS_3$.
We take $L_0$ to be trivial and
$L_1$ to be the twist of
$\cO_{\bP^1}(1)$ by a
nontrivial character $\chi$ of $C_5$.
Then we have the situation of
Lemma \ref{lem.assumptionbundles} with
$H=C_5$ and $N=\fS_3$, and the conditions
of the lemma are satisfied.
We have
\begin{align*}
\xi(X\actsfromright G,(L_0,L_1))&=
(C_5\subseteq C_5,\fS_3\actsfromleft k(\bP^1),\emptyset,(0,\chi))\\
&+
(C_5\subseteq C_5\times \langle (1,2)\rangle ,\mathrm{triv}\actsfromleft k,(0,1),(0,(\chi,0)))\\
&+
(C_5\subseteq C_5\times \langle (1,2)\rangle,\mathrm{triv}\actsfromleft k,(0,1),(0,(\chi,1)))\\
&+
(C_5\subseteq C_5\times \fA_3,\fS_3/\fA_3\actsfromleft k\times k,(0,1),(0,(\chi,1))).
\end{align*}
The outcome of Proposition \ref{prop.fibration} is
\begin{align*}
[&\bP(L_0\oplus L_1)\actsfromright G]=
(\mathrm{triv},G\actsfromleft k(\bP^1)(t),\emptyset)+
(\langle (1,2)\rangle,C_5\stackrel{\chi}{\actsfromleft} k(t),1)\\
&+{\color{red} (C_5,\fS_3\actsfromleft k(\bP^1),\chi)}+
(C_5\times\langle (1,2)\rangle,\mathrm{triv}\actsfromleft k,((0,1),(\chi,0)))\\
&+(C_5\times\langle (1,2)\rangle,\mathrm{triv}\actsfromleft k,((0,1),(\chi,1)))\\
&+(C_5\times \fA_3,\fS_3/\fA_3\actsfromleft k\times k,((0,1),(\chi,1)))\\
&+{\color{red} (C_5,\fS_3\actsfromleft k(\bP^1),-\chi)}+
(C_5\times\langle (1,2)\rangle,\mathrm{triv}\actsfromleft k,((0,1),(-\chi,0)))\\
&+(C_5\times\langle (1,2)\rangle,\mathrm{triv}\actsfromleft k,((0,1),(-\chi,1)))\\
&+(C_5\times \fA_3,\fS_3/\fA_3\actsfromleft k\times k,((0,1),(-\chi,1))).
\end{align*}
In the notation of Section~\ref{sec.filtrations} we 
observe that the $G$-prefilter
\[
\mathbf{H}:=\{(C_5,\fS_3)\}
\]
satisfies the condition of Proposition \ref{prop.filtration}.
Upon projection 
$$
\Burn_2(G)\to \Burn^{\mathbf{H}}_2(G)
$$
(see Definition~\ref{defn.BurnH}), we obtain the class  
$$
{\color{red}
(C_5,\fS_3\actsfromleft k(\bP^1),\chi)
} + 
{\color{red}
(C_5,\fS_3\actsfromleft k(\bP^1),-\chi)}
\in 
\Burn^{\mathbf{H}}_2(G).
$$
This class is nonzero. Moreover, it is different for $\chi\in \{\pm 1\}$ as compared to 
$\chi\in \{\pm 2\}$.

Geometrically, the situation above arises as follows: Consider the 3-dimensional 
representation $W_\chi = 1\oplus (V\otimes \chi)$ of $G$, sum of a trivial $1$-dimensional representation and twist by $\chi$ of the
standard $2$-dimensional representation
$V$ of $\fS_3$. This gives a generically free action of $G$ on $\bP^2=\bP(W_{\chi})$, 
with a $G$-fixed point $\mathfrak{p}$. 
To bring the $G$-action into a form where
Assumption 2 is satisfied, we need to blow up $\mathfrak{p}$, and
$$
[\bP(W_{\chi}) \actsfromright G] = [\bP(L_0\oplus L_1)\actsfromright G] \in \Burn_2(G).
$$
\end{exam}


\section{Products}
\label{sect:prod}

Let $G'$ and $G''$ be finite groups. 
Define a product map
\[ 
\Burn_{n'}(G')\times \Burn_{n''}(G'')\to \Burn_{n'+n''}(G'\times G''). 
\]
On symbols, it is given by
\begin{equation}\label{eqn:prod}
((H', Z'\actsfromleft K',\beta'), (H'', Z''\actsfromleft K'',\beta''))\mapsto
(H, Z\actsfromleft K,\beta), 
\end{equation}
where
\begin{itemize}
\item $H =H'\times H''$,
\item $Z=Z'\times Z''$,
\item $K=K'\otimes_{k} K''$, with the natural action of $Z$,
\item $\beta=\beta'\cup \beta''$.
\end{itemize}

\begin{prop}
The product map \eqref{eqn:prod} is well-defined, and satisfies
$$
([X'\actsfromright G'],[X''\actsfromright G''])\mapsto [X'\times X''\actsfromright G'\times G''].
$$
\end{prop}

\begin{proof}
The map clearly respects relations.
The only point to remark is that in
\textbf{(B2)}, the condition for
nontriviality of $\Theta_2$ holds for
$\beta'$ if and only if it holds for
$\beta=\beta'\cup \beta''$.
\end{proof}

\section{Restrictions}
\label{sec.restr}
Let $G$ be a finite group and $G'\subset G$ a subgroup. 
A $G$-action on a quasiprojective variety $X$ induces an action of $G'$,
and thus it is natural to propose the existence of a restriction homomorphism
from $\Burn_n(G)$ to $\Burn_n(G')$, acting by
\begin{equation}
\label{eqn:restr}
[X\actsfromright G]\mapsto [X\actsfromright G'].
\end{equation}
In this section we establish the existence and uniqueness of this homomorphism. 

\begin{exam}
Suppose $H$ is an abelian subgroup of $G$, contained in $G'$.
Symbols, identified in $\Burn_n(G)$ by relation \textbf{(C)},
might no longer be identified in $\Burn_n(G')$.
E.g., with $G=\mathfrak{D}_4$ and $G'=C_4$ the restriction of
$(G',G/G'\actsfromleft k\times k,1)\in \Burn_1(G)$
to $\Burn_1(G')$ has to be a sum of two symbols with distinct characters:
\[
(G',G/G'\actsfromleft k\times k,1)
\mapsto
(G',\mathrm{triv}\actsfromleft k,1)+(G',\mathrm{triv}\actsfromleft k,3).
\]
\end{exam}

\begin{theo}
\label{theo:res}
For all $n\ge 0$, 
there exists a unique homomorphism of abelian groups
$$
\mathrm{res}^G_{G'}: \Burn_n(G)\to \Burn_n(G').
$$
compatible with \eqref{eqn:restr}. 
\end{theo}

\begin{proof}
By Lemma \ref{lem.Z}, it suffices to consider
symbols of the form
$$
\mathfrak{s}=(H,Z\actsfromleft K,\beta).
$$ 
When we act by conjugation by some element of $G$, we obtain an
equivalent symbol, where $H$ is replaced by a conjugate,
the corresponding centralizer quotient replaces $Z$, and
conjugation is used to form from $\beta$ a sequence of
characters of the conjugate of $H$.
By conjugation we have a transitive action of $G$ on a set $\mathfrak{S}$ of
symbols, where $\mathfrak{s}\in \mathfrak{S}$ has stabilizer $Z_G(H)$.
The restriction of the action to $G'$ consists of finitely many orbits;
in the formula below the sum is over
orbit representatives
\[ \mathfrak{s}'=(H',Z'\actsfromleft K,\beta'), \]
such that the restriction $\beta'|_{H'\cap G'}$ of $\beta'$ to $H'\cap G'$ has trivial space of
invariants; here, $Z'$ denotes $Z_G(H')/H'$.
Then we define the restriction to $G'$ by
\[ \mathfrak{s} \mapsto \sum_{\mathfrak{s}'}
(H'\cap G',(Z_G(H')\cap G')/(H'\cap G')\actsfromleft K,
\beta'|_{H'\cap G'}); \]
this map respects relations.
Uniqueness follows from \cite[Rmk.\ 5.16]{BnG}.
\end{proof}



As an application of the restriction construction, we obtain a map 
$$
\Burn_{n'}(G)\times \Burn_{n''}(G)\to \Burn_{n'+n''}(G),
$$
using the product construction in 
Section~\ref{sect:prod} with $G'=G''=G$, followed restriction to the diagonal
$$
G\subset G\times G. 
$$
This map on Burnside groups sends
$$
([X'\actsfromright G],[X''\actsfromright G])\mapsto [X'\times X''\actsfromright G].
$$



\section{Combinatorial analogs}
\label{sec.combinatorial}

Here we define and study a quotient
$$
\Burn_n(G)\ra \BC_n(G)
$$
to a {\em combinatorial} version of the equivariant Burnside group, by forgetting the 
information about the Galois algebra. 

\begin{defi}
\label{defn:comb}
The \emph{combinatorial symbols group} 
$$
\BC_n(G)
$$
is the $\bZ$-module,
generated by symbols
\[ (H,Y,\beta) \]
with
$H$ abelian,
$Y\subseteq Z_G(H)/H$,
and $\beta$ a sequence of
nonzero elements generating $H^\vee$,
of length at most $n$,
modulo relations: 


\medskip
\noindent
\textbf{(O):}
$(H,Y,\beta)=(H,Y,\beta')$ if
$\beta'$ is a reordering of $\beta$.

\medskip
\noindent
\textbf{(C):} 
$(H,Y,\beta) = (gHg^{-1},gYg^{-1},\beta')$
for $g\in G$, with $\beta$ and $\beta'$
related by conjugation by $g$.

\medskip
\noindent
\textbf{(B1):}
$(H,Y, \beta)=0$ when $b_1+b_2=0$.

\medskip
\noindent
\textbf{(B2):}
$(H,Y, \beta)=\Theta_1+\Theta_2$,
where $\Theta_1$ and $\Theta_2$ are as in Section \ref{sect:gen}, 
i.e., 
\[
\Theta_1= \begin{cases} 0, & \text{if $b_1=b_2$}, \\ 
(H,Y,\beta_1)+(H,Y, \beta_2),
& \text{otherwise},
\end{cases}
\]
with $\beta_1$ and $\beta_2$ as in \eqref{eqn:beta12},
and
\[ \Theta_2= \begin{cases} 0, & \text{if $b_i\in \langle b_1-b_2\rangle$ for some $i$}, \\
(\overline{H},\overline{Y}, \bar{\beta}), & \text{otherwise}, 
\end{cases}
\]
where $\overline{H}=\ker(b_1-b_2)$,
$\overline{Y}$ is the pre-image of $Y$ in
$Z_G(H)/\overline{H}$, and $\bar \beta$ consists
of the restrictions to $\overline{H}$ of the characters of $\beta$.
\end{defi}




\begin{prop}
\label{prop:comb}
The map sending the class of a triple
$$
(H, Y\actsfromleft K, \beta) \in \Burn_n(G),
$$
for fields $K\in \Alg_Y(K_0)$, $K_0\in \Bir_d(k)$, with $d\le n$, 
to 
$$
[k':k](H,Y,\beta) \in \BC_n(G), 
$$
where $k'$ is the algebraic closure of $k$ in $K_0$,
gives a surjective homomorphism 
$$
\Burn_n(G) \ra  \BC_n(G).
$$
\end{prop}

\begin{proof}
This is clear from the description of the relations in
$\Burn_n(G)$ from Section \ref{sect:gen}.
\end{proof}

\begin{defi}
\label{def.BCH}
Given a $G$-prefilter $\mathbf{H}$, we let
\[
\BC_n^{\mathbf{H}}(G)
\]
be the quotient of $\BC_n(G)$ by the subgroup
generated by classes $(H,Y,\beta)$ with
$(H,Y)\notin \mathbf{H}$.
\end{defi}

Exactly as in Section \ref{sec.filtrations} we have

\begin{prop}
\label{prop.BCH}
Let $\mathbf{H}$ be a $G$-prefilter, satisfying the
hypothesis of Proposition \ref{prop.filtration}.
Then $\BC_n^{\mathbf{H}}(G)$ is generated by symbols
$(H,Y,\beta)$ for $(H,Y)\in \mathbf{H}$,
subject to relations \emph{\textbf{(O)}}, \emph{\textbf{(C)}}, \emph{\textbf{(B1)}}, and \emph{\textbf{(B2)}} 
applied to these symbols. 
\end{prop}


Additionally, upon passage to the combinatorial analogue we also have the other structures developed in this paper:
\begin{itemize}
\item equivariant (projectively) indexed combinatorial Burnside group;
\item product map;
\item restriction homomorphisms.
\end{itemize}



\begin{exam}
\label{exa.BCH}
Suppose that $G$ is abelian.
\begin{itemize}
\item We have (cf.\ \cite[\S 8]{BnG})
\[ \cB_n(G)=\BC_n^{(G,\mathrm{triv})}(G), \]
where $\cB_n(G)$ is the symbols group from \cite{kontsevichpestuntschinkel}.
\item There is a commutative diagram
\[
\xymatrix{
\Burn_n(G) \ar[d] \ar[r]  &   \BC_n(G) \ar[d] \\
   \Burn^G_n(G) \ar[r]    &    \cB_n(G)
}
\]
(The factor factor $[k':k]$ in Proposition \ref{prop:comb} matches the
similar factor in \cite[Prop.\ 8.1]{BnG}.)
\end{itemize}
\end{exam}

\section{Applications}
\label{sect:appl}



As a first application of the formalism in \cite{BnG} for nonabelian groups, we gave
in \cite{HKTsmall} an example of $G=C_2\times \fS_3$-actions on $\bP^2$ and a quadric surface $Q\subset \bP^3$,
which were distinguished by the respective classes in $\Burn_2(G)$. 
The actions are {\em stably} $G$-equivariantly rational \cite{lemire}. 


Here we give a further application, for $G=\fS_4$, acting on $\bP^2$ and a del Pezzo surface of degree 6. This example
was treated in  \cite{BanTok}, via birational rigidity techniques (Noether inequality).  



\

\noindent
We recall basic facts about the subgroup lattice of $G=\fS_4$:
\begin{itemize}
\item Conjugacy classes of  nonabelian subgroups:  $\fS_3$, $\sD_4$, $\fA_4$, 
\item Conjugacy classes of abelian subgroups: 
trivial, even $\bZ/2$, odd $\bZ/2$, $\bZ/3$, $\bZ/4$, even $\mathfrak K_4\simeq \bZ/2\oplus \bZ/2$, odd $\mathfrak K_4$.
\end{itemize}


\

\noindent
\textbf{First action:} On $\bP^2$, we consider the projectivization of the standard 3-dimensional representation  $V_3$,
with respect to basis 
$$
(-1,1,1,-1), \quad (1,-1,1,-1), \quad (1,1,-1,-1),
$$
given by the 4 matrices below. 
$$
\sigma:=\begin{pmatrix} 
0 & 1 & 0 \\ 
1 & 0 & 0\\
0 & 0 & 1
\end{pmatrix}
\quad 
\tau:=
\begin{pmatrix}
0 & 0 & 1\\
1 & 0 & 0\\
0 & 1 & 0
\end{pmatrix} 
$$
$$
\lambda_1:=
\begin{pmatrix}
-1& 0& 0\\
0 & 1 & 0\\
 0& 0 & -1 
\end{pmatrix} 
\lambda_2:=
\begin{pmatrix}
-1& 0 & 0\\
0& -1& 0\\
0 & 0 & 1
\end{pmatrix} 
$$
Here $\sigma$ and $\tau$ generate $\fS_3$ and $\lambda_1,\lambda_2$ the even
$\mathfrak K_4$.

The restriction of $V_3$ to $\sD_4$ decomposes as 
$1$-dimensional plus irreducible $2$-dimensional representation. 
Each $\sD_4\subset G$, gives a distinguished line and a point; together these form
a triangle.
We blow up the $3$ points to get a hexagon of lines, which form two
triples of disjoint lines, each line has faithful Klein $4$-group action and
generic stabilizer even $\bZ/2$.
The intersection points of the lines have stabilizer even $\mathfrak K_4$.
We blow up these intersection points.
The result is a wheel of $12$ rational curves:
\[
\xymatrix{   D_1 \ar@{-}[r] \ar@{-}[d] & R_1 \ar@{-}[r]  & D_1' \ar@{-}[r]  & R_2 \ar@{-}[r]  & D_2  \ar@{-}[r] &  R_3 \ar@{-}[d]  \\ 
R_6 \ar@{-}[r]& D_3'  \ar@{-}[r]   & R_5  \ar@{-}[r] & D_3 \ar@{-}[r] & R_4  \ar@{-}[r] &   D_2' 
}
\]
Each rational curve has generic stabilizer even $\bZ/2$, and their
intersection points have stabilizer even $\mathfrak K_4$.
The $12$ curves form $3$ orbits:
$\{ D_1,D_2,D_3\}$ and $\{ D_1',D_2',D_3'\}$ consist of
lines with generic stabilizer even $\bZ/2$ and faithful $\mathfrak K_4$-action, and 
the lines in the $G$-orbit $\{ R_1, \dots, R_6\}$ 
have generic stabilizer even $\bZ/2$ and a nontrivial $\bZ/2$-action. 



The restriction of $V_3$ to $\fS_3$ also decomposes into a
$1$-dimensional and an irreducible $2$-dimensional representation. 
Looking at $G$-orbits, we find a $G$-orbit of $4$ distinguished  $\fS_3$-lines, which intersect in $6$ points with
odd $\mathfrak K_4$ stabilizer. We also have a $G$-orbit of
$4$ distinguished points with $\fS_3$-stabilizer.
We blow up the points with odd $\mathfrak K_4$-stabilizer and also those with $\fS_3$-stabilizer.
We get a $G$-orbit of $6$ lines, each with $\bZ/2$-action and
generic stabilizer odd $\bZ/2$, as well as an orbit of 4 exceptional curves with $\fS_3$-action. 

\

\noindent
\textbf{Second action:} 
Let $X$ be a del Pezzo surface of degree 6 given by 
$$
x_0y_0z_0=x_1y_1z_1 \subset \bP^1\times \bP^1\times \bP^1.
$$ 
We write the action of $G$ in coordinates 
$$
x:=x_0/x_1, \quad  y:=y_0/y_1, \quad z:=z_0/z_1.
$$ 
Then,
$\fS_3=\langle \sigma, \tau\rangle$ acts by permuting $x, y, z$;  $\lambda_1$ changes 
signs on $x$ and $z$, and $\lambda_2$ changes signs on  $x$ and $y$. 

There are three orbits of points, of length 4, with stabilizer $\fS_3$ (see \cite[Lemma 1.3]{BanTok}). 
Blowing these up, we obtain 3 $G$-orbits of $\fS_3$-lines. These 
do not contribute to $[X\actsfromright G]$. 
There are also two $G$-orbits of points with $\sD_4$-stabilizers, these points are precisely the intersection points of  
the 6 lines at infinity, i.e., in the locus 
$$
\{x_0=0\} \cup \{ y_0=0\} \cup \{z_0=0\}.
$$
These lines have generic stabilizer even $\bZ/2$ and a nontrivial $\bZ/2$-action, and they form a single $G$-orbit. 
After we blow up the two orbits of 3 points points, we obtain precisely the wheel configuration we described above. 





\


To summarize, the difference 
\[
[X\actsfromright G] - [\bP^2\actsfromright G]
\]
is a symbol 
\begin{equation}
\label{eqn.difference}
(\text{odd }\bZ/2, \bZ/2 \actsfromleft k(t),(1))
\end{equation}
corresponding to a $G$-orbit of 6 lines with generic stabilizer odd $\bZ/2$ and nontrivial $\bZ/2$-action.
By a computation, analogous to the determination of $\Burn_2(\bZ/2\oplus \bZ/2)$ in
\cite[\S 5.4]{HKTsmall},
the class \eqref{eqn.difference} is nontrivial in $\Burn_2(G)$.





\end{document}




\begin{itemize}
\item $\mathfrak C_n$: We have
$\bP^1=\bP(\mathrm I \oplus V_{\chi})$,
where $V_{\chi}$ is a faithful representation of $\mathfrak C_n$, with (primitive) character $\chi$. Then
$$
[\bP^1\actsfromright \mathfrak C_n] =  (\mathfrak C_n, 1\actsfromleft k, (\chi)) +(\mathfrak C_n, 1\actsfromleft k, (-\chi)) \in \Burn_1(\mathfrak C_n). 
$$
\item 


\item  $\fA_5$:
\begin{multline*}
[\bP^1\actsfromright \fA_5] = 
(\mathfrak C_2, \mathfrak K_4/\mathfrak C_2\actsfromleft k^2 , (1)) +  
(\mathfrak C_3, 1 \actsfromleft k, (1)) + \\
(\mathfrak K_4, 1 \actsfromleft k, (\chi_4) +
(\mathfrak C_5, 1 \actsfromleft k, (\chi_5))
\end{multline*}

\end{itemize}


%%%%%%%%

\subsection{$p$-adic integration}
\label{sect:p-adic}

To our knowledge, the first connections of $p$-adic integrals with birational properties appeared in \cite{batyrev}, which in turn was inspired by \cite{BT-ani}, \cite{BT-toric}. 

To explain these connections, we start with a detour:
Consider 
$$
H_p(x):=\max(1,|x|_p), \quad   x\in \mathbb Q_p,
$$
a {\em local height} on $p$-adic points of the affine line $\bA^1$. 
We proceed to compute the integral 
$$
I_p(s):=\int_{\bA^1(\mathbb Q_p)} H_p(x_p)^{-s} \mathrm dx_p, \quad \Re(s)\gg 0,
$$
following \cite[Section 3]{CLT1}.
Here $\mathrm dx_p$ is the standard Haar measure on $\bQ_p$ normalized by the condition that the volume of $\bZ_p$ equals 1. 
To this end, we introduce the domains
 $$
U(0):=\{ x\,|\, |x|_p\le 1\}, \quad \mathrm{vol}(U(0))=1,
$$
$$
U(j):=\{ x\,|\, |x|_p = p^j\}, \quad \mathrm{vol}(U(j))=p^j(1-\frac{1}{p}),
$$
and 
find that 
$$
I_p(s) =  \int_{U(0)} H_p(x_p)^{-s} \mathrm dx_p +
\sum_{j\ge 1} \int_{U(j)}H_p(x_p)^{-s} \mathrm dx_p  
 =  1+\sum_{j\ge 1} p^{-js}{\rm vol}(U(j)).
$$
It follows that 
$$
I_p(s) = \frac{1-p^{-s}}{1-p^{-(s-1)}}.
$$
Evaluating at $s=2$, which corresponds to the anticanonical class $\cO(2)$ on $\bP^1$, we obtain 
\begin{equation}
\label{eqn:p1}
I(2)= (1+\frac{1}{p})  =\frac{ \# \mathbb P^1(\mathbb F_p)}{p}.
\end{equation}
We observe that the $p$-adic integral over the Zariski open $\bA^1\subset \bP^1$ recovers the number of $\bF_p$-points on $\bP^1$. 

More generally, consider $X=X_{\Sigma}$, a smooth projective equivariant compactification of the algebraic torus $T=\mathbb G_m^d$. Very briefly, $X$ is described by combinatorial data of  
\begin{itemize}
\item $N\simeq \bZ^d$, the lattice of cocharacters, and $M= \mathrm{Hom}(N,\bZ)$, the lattice of characters of $T$,
\item $\Sigma=\{ \sigma\}$, a  {\em fan}, spanned by vectors
$e_1,\ldots, e_n \in N$, subject to certain conditions ensuring smoothness and projectivity of $X$, 
\end{itemize}
see \cite{fulton} for basic notions concerning toric varieties. 

There is a fundamental exact sequence describing the Picard group:
 $$
0\to M\ra \mathrm{PL}(\Sigma) \to \mathrm{Pic}(X) \to 0,
$$
where $\mathrm{PL}(\Sigma)$ is the group of piece-wise linear functions on $\Sigma$; a function 
$$
\varphi_{\mathbf{s}} \in  \mathrm{PL}(\Sigma), \quad {\bf s}=(s_1, \ldots, s_n),
$$ 
is determined by its values on the generators $e_j$: $\varphi_{\bf{s}}(e_j)=s_j$.

Using the fact that 
$$
T(\mathbb Q_p)/T(\mathbb Z_p) = N
$$
we can introduce {\em local heights}
$$
H_p(\varphi_{\bf s}, t_p):= p^{\varphi_{\bf{s}}(\bar{t}_p)}, $$
where $\bar{t}_p$ is the image of $t_p$ in $N$,
and consider the {\em height integrals}
$$
I_p({\bf s}):=\int_{T(\bQ_p)} H_p (\varphi_{\bf s}, t_p)^{-1} \, \mathrm dt_p, 
$$
where $\mathrm dt_p$ is a Haar measure on $T(\bQ_p)$, normalized by the compact subgroup $T(\bZ_p)$, and $s_j\in \bC$ have $\Re(s_j)>1$, so that the integral converges. 


Following \cite[Section 2]{BT-ani}, we proceed to compute $I_p({\bf s})$ by
reorganizing the sum over all lattice points in $N$ as the alternating sum of contributions from $\Sigma(r)$, the set of $r$-dimensional subcones of $\Sigma$:
$$
I_p({\bf s})=
 \sum_{r=1}^d \sum_{\sigma\in \Sigma(r)} (-1)^{r} \left( \sum_{n\in \sigma\cap N} p^{-\varphi_{\bf s} (n)} \right).
 $$
We obtain
$$
I_p({\bf s})= \sum_{r=1}^d \sum_{\sigma\in \Sigma(r)} (-1)^r \prod_{e_j\in \sigma}  \frac{1}{1-p^{-s_j }}. 
$$
Note that this sum is unchanged, if we subdivide a cone $\sigma \in \Sigma$, and replace $\varphi_{\bf s}$ by the induced piece-wise linear function $\tilde{\varphi}_{\bf s}$ on the new fan $\tilde{\Sigma}$. Such a subdivision corresponds to a toric blow-up of $X$ and the straighforward observation is that $I_p({\bf s})$ remains invariant under such blow-ups. 

Evaluating at the distinguished element 
${\bf s}=(1,\ldots, 1)$ (which corresponds to the anticanonical class $-K_X\in \Pic(X)$), we obtain 
\begin{equation}
\label{eqn:toric}
I_p((1,\ldots, 1)) =\frac{ \# X(\bF_p)}{p^d}, 
\end{equation}
which, as in the case of \eqref{eqn:p1}, is 
interpreted  as a {\em volume} with respect to a natural measure $\tau_p$, the {\em Tamagawa measure} introduced in \cite{Peyre}. Again, the $p$-adic integral over $T(\bQ_p)$ recovers the number of $\bF_p$-points on the compactification $X$. 

The Tamagawa measure is defined for all smooth projective $X$ equipped with a {\em metrization} of the canonical line bundle $K_X$. Metrizations of line bundles naturally arise for projective varieties over number fields $k$, it suffices to choose models over the integers. 


In detail, for $x\in X(k_v)$, where $k_v$ is the completion of $k$ at a nonarchimedian place $v$, choose local analytic coordinates $x_{1}, \ldots, x_{d}$, in a
neighborhood $U_x$. In $U_x$, a section of the canonical line bundle has  
the form $\mathsf s :=\mathrm dx_1\wedge \ldots \wedge \mathrm dx_d$. Put
$$
\tau_v=\tau_{X,v} :=\|\mathsf s\|_{v} \mathrm dx_1\cdots \mathrm dx_d,
$$
where $\mathrm dx_1\cdots \mathrm dx_d$ is the standard normalized Haar measure on $k_v^d$; it globalizes to $X(k_v)$. 
For almost all $v$, and Zariski open $U\subset X$, one has
$$
\int_{U(k_v)}\!\! \tau_{v} = 
\int_{X(k_v)}\!\! \tau_{v} = \int_{X(\mathfrak o_v)}\!\! \tau_{v}  =\sum_{\tilde{x}\in X(\mathbb F_q)} \int_{\pi^{-1}(\tilde{x})} \tau_{v} = \frac{\# X(\F_q)}{q^{d}},
$$
where $\pi:X(k_v)\to X(\bF_q)$ is the reduction map. 
The first equality is due to the fact that the boundary $X\setminus U$ has measure zero. 
Again, the integral over the $v$-adic points of the Zariski open $U$ yields information about $\bF_q$ points on the compactification $X$. 

This leads to a surprising application \cite{batyrev}: Let 
$X$, $Y$ be {\em birational} Calabi-Yau varieties of dimension $d$. The canonical bundle of a Calabi-Yau variety has a canonical metrization, and the corresponding Tamagawa measure integrates to 
\begin{equation}
\label{eqn:tamagawa}
\int_{X(k_v)}\tau_v = \frac{\# X(\mathbb F_q)}{q^{d}},  \quad \forall' v.
\end{equation}
Birationality of $X$ and $Y$ implies that there exists a commond Zariski open subset $U$, 
$$
X\supset U\subset Y.
$$
But then 
$$
\frac{\# X(\mathbb F_q)}{q^{d}} = \int_{X(F_v)} \tau_{v} = \int_{U(F_v)} \tau_{v}  = \int_{Y(F_v)} \tau_v = 
\frac{\# Y(\mathbb F_q)}{q^{d}}, \quad \forall' v. 
$$
Combining with the Weil conjectures one obtains information about cohomology of $X$ and $Y$: they have equal {\em Betti numbers}. 


More generally, following \cite{CLT-I}, let 
$X$ be a smooth projective variety of dimension $d$ over a local field $k_v$. Consider a normal crossings divisor $D\subset X$,
with
$$
D=\cup_{\alpha\in \mathcal A} D_{\alpha},    \quad -K_X = \sum \rho_{\alpha} D_{\alpha},
$$ 
where $D_{\alpha}$ are geometrically irreducible, smooth, and intersecting transversally. For $A\subset  \mathcal A$ let 
$$
D_A:=\cap_{\alpha\in A} D_{\alpha} , \quad D_A^{\circ} = D_A\setminus \cup_{A'\supset A} D_{A'},
$$
note that $D_A\subset X$ is smooth, of codimension~$\#A$ (or empty).

Let $U:=X\setminus D$ and consider the {\em local height}
$$
H_{\alpha}: U(k_v) \ra \bR_{\ge 0}, 
$$
that should be thought of as the $v$-adic distance to the boundary component $D_{\alpha}$.
The corresponding height integral is 
$$
I_v(\mathbf s):=\int_{U(F_v)} \,\, \prod_{\alpha \in \mathcal A} \,\,  H_{\alpha} (x)^{-s_{\alpha}}  \mathrm  d \tau_v,
$$ 
where $\mathrm  d \tau_v$ is the Tamagawa measure from \cite{Peyre}. 

As in \cite[Section 4]{CLT-I}, we can perform the computation
in {\em charts}, via partition of unity:
in a neighborhood of $x\in D_A^\circ(\bF_q)$ it takes the form
$$
 \int \prod_{\alpha\in A}      |x_{\alpha}|_v^{s_\alpha -\rho_\alpha} \,\mathrm d\tau_v.
$$
Essentially, this is a product of integrals of the form
$$
\int_{|x|_v\le 1} |x|_v^{s-1} \mathrm dx_v.
$$
{\em Denef's formula} yields (for almost all $v$):
\begin{equation}
\label{eqn:den}
I_v(\mathbf s) = \sum_A  \frac{\# D_A^\circ(\F_q)}{q^{d}} 
\prod_{\alpha\in A} \frac{q-1}{q^{s_\alpha-\rho_\alpha+1}-1} .
\end{equation}
Evaluating at $s_\alpha=\rho_{\alpha}$, for all $\alpha\in \mathcal A$, we obtain, as in \eqref{eqn:p1} and \eqref{eqn:toric}:
\begin{equation}
\label{eqn:igusa}
 I_v(\rho) = \sum_A  \frac{\# D_A^\circ(\F_q)}{q^{d}}  = \frac{\#X(\mathbb F_q)}{q^{d}}.
\end{equation}
One should view this formula as a refinement of \eqref{eqn:tamagawa}. Thus, the integral \eqref{eqn:den} (which plays a central role in analytic/spectral approches to Manin's conjectures, volume asymptotics, etc., see \cite{CLT-I}), is a {\em birational invariant}. 
In particular, it can be defined for singular $X$, in which case it will encode information about the singularities of $X$. 








%\centerline{\includegraphics[width=8cm]{p2}}



%%%%%%%%



\section{Computing the class}



\subsection{Curves}

For smooth projective curves, (equivariant) birationality equals (equivariant isomorphism). Correspondingly, there are no relations between symbols in $\Burn_n(G)$ for $n=1$. On the other hand, these symbols enter in the analysis of relations for $n\ge 2$. 

Recall that the only groups acting generically freely on $\bP^1$ are 
$$
\mathfrak C_n, \mathfrak D_{2n}, \fA_4, \fS_4, \fA_5.
$$
Cyclic actions are linear, the dihedral group acts linearly iff $n$ is even, and the actions of the tetrahedral, octahedral, and icosahedral groups are not even stably linearizable: the Amitsur invariant $\mathbf{(Ami)}$ does not vanish. We list the corresponding classes:


We compute classes in $\Burn_1(G)$ of $G$-actions on $\bP^1$ (using the centralizer notation):
$$
\begin{array}{r|c} 
G & [\bP^1 \actsfromright G]\\
\hline
C_n &  (C_n, 1\actsfromleft k, \epsilon_n)  + (C_n, 1\actsfromleft k, -\epsilon_n) \\
\mathfrak D_{2n} &  \\
\fA_4  & (C_2, \bZ/2\actsfromleft k^2,\epsilon_2) + (C_3, 1\actsfromleft k,\epsilon_3)\\
\fS_4 &  (C_2, \bZ/2\actsfromleft k^2,\epsilon_2) + (C_3, 1 \actsfromleft k,\epsilon_3) + (C_4, 1\actsfromleft k,\epsilon_4) \\
\fA_5 & (C_2, \bZ/2\actsfromleft k^2, \epsilon_2) + 
(C_3,  1\actsfromleft k, \epsilon_3) + (C_5, 1\actsfromleft k , \epsilon_5)
\end{array}
$$

\subsection{Surfaces: linear actions}

In this section, we illustrate the machinery of the previous sections, in its application 
to linear and projective actions of the 
icosahedral group $G=\fA_5$. 


Recall some basic facts about
(conjugacy classes of) subgroups of $\fA_5$:
\begin{itemize}
\item Nonabelian subgroups: $\mathfrak S_3$ (length 10),  $\sD_5$ (of length 6), and $\fA_4$ (length 5), 
\item Abelian subgroups: 
$\bZ/2$ (length 15), $\bZ/3$ (length 10), $\bZ/5$ (length 6), and the Klein 4-group $\fK_4 = \bZ/2\oplus \bZ/2$ (length 5). The 
corresponding normalizers are $\fK_4$, $\fS_3$, $\sD_5$, and $\fA_4$, respectively. 
\end{itemize}
Here are the irreducible representations of $\fA_5$ (notation of \cite{ChS}):
\begin{itemize}
\item $I$ - the trivial representation,
\item $W_5, W_4, W_3, W_3'$ -- 5, 4, 3-dimensional (faithful), 
\end{itemize}
and of $2.\fA_5$ (the binary icosahedral group):
\begin{itemize}
\item $U_4, U_2, U_2'$ -- 4 and 2-dimensional.
\end{itemize}

This gives actions of $G=\fA_5$ on $\bP^3$, via
$$
\bP(W_4), \bP(1 \oplus W_3),  \bP(U_2\oplus U_2), \bP(U_2\oplus U_2'), \bP(U_4)
$$
By \cite[Remark 1.2.1]{ChS}, neither of the first two $G$-actions is $G$-equivariant to the third or the fourth action -- the Klein subgroup $\fK_4=\bZ/2\oplus \bZ/2\subset G$
has fixed points for the first two actions, but not for the third and the fourth. 

The action on $\bP^2=\bP(W_3)$ is via the projectivization of one of the 3-dimensional representations $W_3$ of $\fA_5$. 
We analyze this action in two ways: first geometrically, and then in terms of the combinatorial structures introduced in Section~\ref{sect:models}. 



The representation-theoretic facts about restrictions of $W_3$ to subgroups are summarized in, e.g., \cite[Corollary 5.2.2]{ChS}: 
\begin{itemize}
\item
to $\fS_3$ and $\sD_5$ decompose into nontrivial 1-dimensional and a 2-dimensional representation, 
\item to $\bZ/5$ is a sum
of a trivial and two different 1-dimensional, 
\item to $\fK_4=\bZ/2\oplus\bZ/2$ is a sum of three different nontrivial characters, 
\item 
to $\bZ/3$ is a sum three of different characters,
\item to $\bZ/2$ gives one trivial and two nontrivial characters. 
\end{itemize}
Note  that $\sD_5$ has two conjugacy classes of elements of order 5: one acts with weights (2,3) and the other with weights (1,4) in 
the 2-dimensional representation. In $\fA_5$ there is only one conjugacy class of elements of order 5. 




This translates into the following orbit description, see \cite[Proposition 6.1.2]{ChS}; we write $\Sigma_m$ for an orbit of points and $L_m$ for an orbit of lines, with nontrivial stabilizers and of length $m\mid 60$, respectively:

\

\begin{itemize}
\item $\Sigma_6$ / stabilizer $\sD_5$ -- irreducible representation on the tangent space, after blowing up 
get an orbit of $\sD_5$-lines, no generic stabilizer, 
fixed points with stabilizers $\bZ/2$ and $\bZ/5$.
\item $\Sigma_{10}$ / stabilizer $\fS_3$ --   
irreducible representation on the tangent space, blow-up gives $\fS_3$-lines, no generic stabilizer,
fixed points with stabilizers $\bZ/2$ and $\bZ/3$.
\item $\Sigma_{12}$ / stabilizer $\bZ/5$ -- note that
any subgroup $\bZ/5\subset \fA_5$ fixes three points,
one of them with the larger stabilizer $\sD_5$,
the other two swapped by the action of an element of
$\sD_5$, not in $\bZ/5$.
\item $\Sigma_{15}$ /  stabilizer $\fK_4$ -- three points stabilized by a given
$\fK_4\subset \fA_5$ lie on a triangle of lines, each having as generic stabilizer
one of the order $2$ subgroups of $\fK_4$.
\item $\Sigma_{20}$  / stabilizer $\bZ/3$ -- note that
any subgroup $\bZ/3\subset \fA_5$ fixes three points, one of them with
stabilizer $\fS_3$, the other two swapped by the action of an odd element of
$\fS_3$.
\item $L_{15}$ / stabilizer $\bZ/2$ -- arising as generic stabilizer of a line in the
triangle associated with a subgroup $\fK_4\subset \fA_5$.
\end{itemize}

Besides the orbits of lines $L_{15}$ with generic stabilizer $\bZ/2$ there are
orbits of lines $L_6$ and $L_{10}$, each with trivial generic stabilizer,
which appear as projectivization of the $2$-dimensional representation
in the decomposition of the
restriction to $\sD_5$, respectively $\fS_3$.
Table \ref{table61015} summarizes the intersections among the
orbits of lines $L_6$, $L_{10}$, $L_{15}$.
\begin{table}[h]
\[
\begin{array}{c|c|c|c}
&L_6&L_{10}&L_{15} \\
\hline
L_6&\text{5 $\fK_4$-points}&\text{5 $\fK_4$-points($\times 2$)}&
\begin{array}{c}\text{5 $\fK_4$-points($\times 2$)}\\ \text{5 $\bZ/2$-points}\end{array}
%\medskip
\\
\hline
L_{10}&\text{3 $\fK_4$-points($\times 2$)}&
\begin{array}{c}\text{3 $\fK_4$-points}\\ \text{6 $\bZ/2$-points}\end{array} &
\begin{array}{c}\text{3 $\fK_4$-points($\times 2$)}\\ \text{9 $\bZ/2$-points}\end{array}
%\medskip
\\
\hline
L_{15}&
\begin{array}{c}\text{2 $\fK_4$-points($\times 2$)}\\ \text{2 $\bZ/2$-points}\end{array}&
\begin{array}{c}\text{2 $\fK_4$-points($\times 2$)}\\ \text{6 $\bZ/2$-points}\end{array}&
\begin{array}{c}
\text{2 $\fK_4$-points}\\
\text{2 $\fS_3$-points($\times 2$)}\\
\text{2 $\sD_5$-points($\times 4$)}
\end{array}
\end{array}
\]
\caption{Intersections of lines in $\bP(W_3)$.
Left-hand column refers to \emph{each} line,
entries in the row accounts for intersections with \emph{all} of the
lines, respectively in the case of diagonal entries, all of the
other lines, in the indicated orbits of lines.
Multiplicities are indicated in case of points of intersection
with more than one line.}
\label{table61015}
\end{table}

In order to achieve Assumption 2 and compute the class in the
Burnside group, we certainly
need to blow up all points with nonabelian stabilizer.
At the $\fK_4$-points, we would need two independent characters to arise
from $\fA_5$-linearized line bundles, and this is immediate if we blow these
up as well.
So we blow up $\Sigma_6\cup \Sigma_{10}\cup \Sigma_{15}$ to obtain
$\widetilde{\bP}^2$ with orbits of exceptional curves
$E_6$, $E_{10}$, and $E_{15}$.
We have
\begin{itemize}
\item
Each curve in $E_6$ has $2$ points with stabilizer $\bZ/5$, and
$5$ points on the proper transform of
$L_{15}$ with stabilizer $\bZ/2$, the generic stabilizer of $L_{15}$.
\item 
Each curve in $E_{10}$ has $2$ points with stabilizer $\bZ/3$, and
$3$ points on the proper transform of $L_{15}$.
\item
Each curve in $E_{15}$ has generic stabilizer $\bZ/2$ and has
two $K_4$-points, the intersection points with the proper transform of
$L_{15}$.
\end{itemize}

We compute $[\bP^2\actsfromright \fA_5]\in \Burn_2(\fA_5)$
directly from the model $\widetilde{\bP}^2$:
\begin{itemize}
\item $\bZ/2$: $L_{15}$ contributes
$(\bZ/2,\fK_4/(\bZ/2)\actsfromleft \bC(t),1)$,
$E_{15}$ contributes
$(\bZ/2,\fK_4/(\bZ/2)\actsfromleft \bC(t),1)$.
\item $\bZ/3$: $\Sigma_{20}$ contributes
$(\bZ/3,\fS_3/(\bZ/3)\actsfromleft \bC^2,(1,2))$,
points on $E_{10}$ contribute
$(\bZ/3,\fS_3/(\bZ/3)\actsfromleft \bC^2,(1,1))$.
\item $\fK_4$: two orbits, contribution
$2(\fK_4,\fA_4/\fK_4\actsfromleft \bC^3,((1,0),(0,1)))$.
\item $\bZ/5$: $\Sigma_{12}$ contributes
$(\bZ/5,\sD_5/(\bZ/5)\actsfromleft \bC^2,(1,2))$,
points on $E_6$ contribute
$(\bZ/5,\sD_5/(\bZ/5)\actsfromleft \bC^2,(1,3))$.
\end{itemize}

%{\bf rewrite in terms of centralizers?}

Using relations in the Burnside group, the contributions from
$\bZ/2$ and $\fK_4$, as well as those from $\Sigma_{20}$, cancel. 
Direct computation reveals contribution from $\bZ/5$ and hence the entire  class 
$[\bP^2\actsfromright \fA_5]\in \Burn_2(\fA_5)$
as
\begin{equation}
\label{eqn:p2a5}
(\bZ/5,\sD_5/(\bZ/5)\actsfromleft \bC^2,(1,1)) +  (\bZ/3,\fS_3/(\bZ/3)\actsfromleft \bC^2,(1,1))
\end{equation}
The classes of
$(\bZ/5,\sD_5/(\bZ/5)\actsfromleft \bC^2,(a,b))$ for nonzero weights
$a$, $b$ span a subgroup of $\Burn_2(\fA_5)$ isomorphic to $\bZ/2\oplus \bZ/2$.
In this subgroup, $(\bZ/5,\sD_5/(\bZ/5)\actsfromleft \bC^2,(1,1))$ is nonzero.

Notice, $\bP(W'_3)$ differs from $\bP(W_3)$ only by $5$th roots of unity.
The computation in that case would yield
$(\bZ/5,\sD_5/(\bZ/5)\actsfromleft \bC^2,(2,2))$.
But this has the same class 
in the subgroup of $\Burn_2(\fA_5)$, mentioned above.
This is consistent with the known equivariant birational equivalence of
$\bP(W_3)$ and $\bP(W'_3)$ \cite[Remark 6.3.9]{ChS}.


\subsection{Surfaces: nonlinear actions}





%%%%%%%%%%


\section{From the Stacky paper}




Here is the road map of this paper: 
\begin{itemize}
\item We recall basic definitions and examples of stacks. In particular, we discuss notions of birationality in the category of stacks and of toric stacks. 
\item We explain basic (stable) birational invariants of stacks, such as torsion in Picard group, Brauer group, class in the Burnside group, and present examples of 
computations. 
\item Toric stacks and new obstructions for toric stacks in Levchenko.
\item We construct a specialization homomorphism
\item We adopt the universal torsor formalism of Colliot-Th\'el\`ene and Sansuc to stacks. 
\item Examples (hopefully, for toric surfaces -- non-birational but stably birational stacks). 

\end{itemize}






For the rest of the section we suppose that $k$ contains all roots of unity.
We are also interested in the question, to what extent we can obtain such an $\cX$, whose stabilizer groups are in a certain sense no worse than those of $\cU$.
Formal definition would be that for every finite group $H$ every component of the (nonsingular locally closed) substack of $\cX$ with geometric stabilizer group isomorphic to $H$ has nontrivial intersection with $\cU$.

Construction.
We apply the divisorialification algorithm (Proposition \ref{prop.orbifolddivisorialification}, and its proof) to $\cX$, with snc divisor $\cD$.
This leads to a sequence of blow-ups $\cY\to\dots\to \cX$, where each is a blow-up in a smooth center that meets the snc divisor transversely (Remark \ref{rem.orbifolddivisorialification}), let us say
\[ \cE_1,\dots,\cE_r; \]
and with each blow-up the snc divisor is pulled back and the new exceptional divisor is appended.
Finally $\cY$ is divisorial with respect to its snc divisor.
The copy of $\cU$ in $\cX$ has also been iteratively blown up to obtain
$\cY\setminus \cD$.
We relabel $\cY$ as $\cX$, and $\cY\setminus \cD$ as $\cU$.
Then we have $\cX$, divisorial with respect to snc divisor
\[ \cD\cup \cE, \qquad \cD=\cD_1\cup\dots\cup \cD_\ell,\qquad \cE=\cE_1,\dots,\cE_r, \]
so $\cU=\cX\setminus \cD$ is divisorial with respect to $\cE\cap \cU$.

We suppose that $\cX$ is \emph{not} a minimal compactification of $\cU$.
That is, there exists $x\in \cX$, with stabilizer $H_x$ and component
$\Xi_x$ of the locus with stabilizer $H_x$, such that
\[
\Xi_x\subset \cD.
\eqno{(*)}
\]

Since $\cX$ has $\cD\cup \cE$ as snc divisor, $(\cX,\cD)$ is a \emph{standard pair} \cite[Defn.\ 3.5]{berghrydh} over $[\A^r/\G_m^r]$, i.e., the associated morphism $\cX\to [\A^r/\G_m^r]$ is smooth (cf.\ Remark \ref{rem.orbifolddivisorialification}), as are the restrictions to the $\cD_i$ and their intersections.
We may therefore apply functorial destackification to $(\cX,\cD)$ over $[\A^r/\G_m^r]$ to obtain a sequence of stacky blow-ups, whose centers are smooth over $[\A^r/\G_m^r]$ and have normal crossing with $\cD$, to obtain an orbifold, which is obtained by an iterated root construction from an orbifold $\cX'$ with representable smooth morphism to $[\A^r/\G_m^r]$.

Since $\cU\to [\A^r/\G_m^r]$ is already representable, the destackification procedure leaves $\cU$ untouched.
Since the destackification is carried out over $[\A^r/\G_m^r]$,
only blow-ups whose center meets $\cE$ transversally take place
(where $\cE$ is pulled back after each blow-up).
Each root operation is along a divisor in the complement of $\cU$, and has also the impact of a root operation on $\cE$.
As a consequence, the snc divisor $\cE'$ determined by $\cX'\to [\A^r/\G_m^r]$ is the closure of $\cE\cap \cU$.

\begin{exam}
\label{exa.Xprimenogood}
There is no reason to expect $\cX'$ to be a minimal compactification of $\cU$.
Indeed it is possible for $\cX$ (divisorial compactification with respect to snc divisor $\cD\cup \cE$ with $\cU=\cX\setminus \cD$) to be representable over $[\A^r/\G_m^r]$ -- in which case destackification over $[\A^r/\G_m^r]$ does nothing -- yet not a minimal compactification.
We use $\PP(\ell,m,n)$, $\G_m$-quotient of $\A^3\setminus \{0\}$ analogous to $\PP(m,n)$, in particular $\PP(1,1,2)$, which we observe has three coordinate substacks: $\PP(1,1)\cong \PP^1$, a copy of $\PP(1,2)$, and another copy which we denote by $\PP(1,2)'$.
Let $\cU=\PP(1,1,2)\times \A^1$.
The natural candidate for $\cX$ is $\PP(1,1,2)\times \PP^1$, with $\cD=\PP(1,1,2)\times \{\infty\}$ and $\cE=\PP(1,2)\times \PP^1$.
However, we rather take
\[ \cX=B\ell_{\PP(1,2)'\times \{\infty\}}(\PP(1,1,2)\times \PP^1), \]
with $\cD=\cD_1\cup \cD_2$ and $\cE$ the pre-image of $\PP(1,2)\times \PP^1$.
Then $\cX$ is representable over $[\A^1/\G_m]$, but there is an isolated point with stabilizer $\Z/2\Z$, so $\cX$ is not a minimal compactification of $\cU$.
Thus it is possible to start with a divisorial compactification $(\cX_0,\cD_0\cup \cE_0)$, $\cU=\cX_0\setminus \cD_0$, with $\cX_0$ a minimal compactification, and blow it up along a smooth substack meeting $\cE_0$ transversely and having normal crossing with $\cD_0$ to obtain $(\cX,\cD\cup \cE)$, such that $\cX$ is not a minimal compactification of $\cU$.
\end{exam}

Local analysis near a point $x\in \cX$, satisfying $(*)$.
Let $H=H_x$, and let $d$ denote the
dimension of $\Xi_x$.
After renumbering, we may suppose that $x$ lies on $\cD_1$, $\dots$, $\cD_a$, $\cE_1$, $\dots$, $\cE_b$, and no others.
We take $c$, so that
\[ a+b+c+d=n. \]
Then on a local chart of the form $[V/H]$ at a fixed point of affine $V$ over $x\in \cX$ there are the corresponding characters
\[ \chi_1,\dots,\chi_a,
\chi'_1,\dots,\chi'_b,\chi''_1,\dots,\chi''_c,\overbrace{0,\dots,0}^d. \]
There is not much point to local analysis.
We had guessed that the condition
\[ \langle \chi_1,\dots,\chi_a\rangle \cap \langle \chi'_1,\dots,\chi'_b\rangle=0, \]
which translates into a locus in $\cU$ with stabilizer $H_0$ running into $\Xi_x$ with stabilizer $H$, collapsing to stabilizer $H_0$ when mapped to the relative coarse moduli space over $[\A^r/\G_m^r]$, would be preserved under stacky blow-up with center meeting $\cE$ transversely and having normal crossing with $\cD$.
But the condition already fails to be preserved under ordinary blow-up, see Example \ref{exa.Xprimenogood}.

In any event let us record a local example in toric language \cite{BCS}.
Given finite-index $N'\subset N$ we set $H=N/N'$.
Then, for any $\Z$-basis $v_1$, $\dots$, $v_n$ of $N'$ we obtain
\[ H\hookrightarrow \G_m^n\actsfromleft \A^n, \]
given by ($c_1$, $\dots$, $c_n\in \Q$ such that $\sum c_jv_j\in N$)
\[ \sum c_jv_j\mapsto (e^{2\pi ic_1},\dots,e^{2\pi ic_n}) \]
and thus $n$ characters $\chi_1$, $\dots$, $\chi_n\in H^\vee$.
(Detail: $\Hom(H,\G_m)=\ker(\Hom(N,\mu_\infty)\mapsto \Hom(N',\mu_\infty))$,
identify $\mu_\infty\cong\Q/\Z$, $\Hom(N,\Q)\cong\Hom(N',\Q)$, get $\Hom(H,\G_m)\cong M'/M$, thus $H\cong\Spec k[M'/M]$.)
Three worked examples ($n=2$, $a=1$, $b=1$, $N=\Z^2$):
\begin{itemize}
\item $N'=4\Z\oplus \Z$, $v_1=(4,1)$, $v_2=(4,2)$, $\chi_1(\bar e_1)=-1$, $\chi_2(\bar e_1)=-i$, $H_0=\Z/2\Z$, $H=\Z/4\Z$, $|\langle \chi_1\rangle\cap \langle \chi_2\rangle|=2$, solution root and blow up to obtain new basis $(12,4)$, $(4,2)$, new group $\Z/4\Z\oplus \Z/2\Z$,
new characters $(i,-1)$, $(-1,-1)$ spanning disjoint subgroups.
\item $N'=4\Z\oplus \Z$, $v_1=(4,2)$, $v_2=(4,1)$, $\chi_1(\bar e_1)=-i$, $\chi_2(\bar e_1)=-1$, $H_0=0$, $H=\Z/4\Z$, $|\langle \chi_1\rangle\cap \langle \chi_2\rangle|=2$, solution blow up twice new basis $(12,4)$, $(4,1)$, new characters $-i$, $1$ spanning disjoint subgroups.
\item $N'=2\Z\oplus \Z$, $v_1=(2,1)$, $v_2=(0,1)$, $\chi_1(\bar e_1)=-1$, $\chi_2(\bar e_1)=-1$, $H_0=0$, $H=\Z/2\Z$, $|\langle \chi_1\rangle\cap \langle \chi_2\rangle|=2$, solution blow up new basis $(2,2)$, $(0,1)$, new characters $-1$, $1$ spanning disjoint subgroups.
\end{itemize}
The strategy, to root boundary divisors and perform blow-ups where they meet components of $\cE$, to eliminate $\langle \chi_1,\dots,\chi_a\rangle\cap \langle\chi'_1,\dots,\chi'_b\rangle$, is
not viable due to the lack of preservation of the vanishing of $\langle \chi_1,\dots,\chi_a\rangle\cap \langle\chi'_1,\dots,\chi'_b\rangle$, mentioned above.

How to proceed without minimal compactifications?
There is \cite[Lemma 7.5]{BnG}, which makes use of $[X]$, $X$ a coarse moduli space,
If we compactify $\cU$ to $\cX$ without attention to minimality, could we define $[\cU]^{\mathrm{naive}}$ in $\oBurn_n$ as a sum over (terminology from Section \ref{sect:g}) stabilizer components of $\cU$ (just copy the definition with $\cU$ in place of $\cX$)?
In the sum each $\cK$ is the closure in $\cU$ of a locus $\cK^\circ$ where the stabilizer is constant,
$\cK^\circ=\cK\setminus \cK^{>}$ (where $\cK^{>}$ is the locus in $\cK$ where the stabilizer jumps).
We form closures $\overline{\cK}$, $\overline{\cK^{>}}$ in $\cX$ and associate the class of
\[ \overline{\cK}\setminus \overline{\cK^{>}}, \]
in $\oBurn_n$, specifically, the pair consisting of the coarse moduli space and the normal bundle representation at the generic point.





\begin{exam}
Let $G:=C_2^2\times \fS_3$. 
The symbol 
$$
(C_2,C_2\times \fS_3\actsfromleft k(Q),(1)),
$$
where $Q\subset \bP^3$ is the diagonal smooth quadric, with action of
$C_2$ by sign on the first coordinate and
of $\fS_3$ by permutation on the remaining coordinates, is incompressible in $\Burn_3(G)$.
We have to distinguish this from symbols
$$
(\bar H,\bar{Y} \actsfromleft k(S), (b))
$$
arising as $\Theta_2$ in a blow-up relation.
A symbol $(H,\{1\}\actsfromleft k,(b_1,b_2,b_3))$ can only give rise
to symbol with an action of an \emph{abelian} group $H/\bar H$, which
$\fS_3\times C_2$ is not.
What remains is a symbol
$(H,Y\actsfromleft k(C), (b_1,b_2))$.
For $k(C)(t)$ to be rational, we must have $C$ rational, $k(C)\cong k(\bP^1)$.
Now we must have $Y=G/H$ with $H$ \emph{central} in $G$, containing the first factor $C_2$, the possibilities are
$H=C_2$, $H=C_2^2$.
From $H=C_2$ and $b_1=b_2=1$ we obtain
\[ (C_2,C_2\times \fS_3\actsfromleft k(\bP^1)(t),(1)), \]
the action coming from the unique faithful action of $C_2\times \fS_3$ on
$k(\bP^1)$ and trivial action on $t$.
From $H=C_2^2$, the condition $\ker(b_1-b_2)=C_2$ (first factor) forces
(up to order) $b_1=e_1$, $b_2=e_1+e_2$, and we obtain
\[ (C_2,C_2\times \fS_3\actsfromleft k(\bP^1)(t),(1)), \]
where $\fS_3$ acts on $k(\bP^1)$ and $C_2$ acts on $t$.
The existence of fixed points for the subgroup $C_2\times \fS_2$
(on a projective model with regular action), respectively, the multiplicity
of the incompressible symbol
\[ (C_2,\fS_3\actsfromleft k(\bP^1),(1))\in \Burn_2(C_2\times \fS_3) \]
distinguishes the birational equivalence classes of the
$C_2\times \fS_3$-actions on $k(\bP^1)(t)$.
\end{exam}


%%%%%%%%%%%%%%





\section{Limitations}
\label{sect:limit}


\subsection*{Fixed points}

The Burnside group formalism does not capture the obstruction {\bf (Fix)} from Section~\ref{sect:fix} -- existence of fixed points upon restriction to abelian subgroups. This is easy to see from the fact that $\Burn_n(G)$ is a {\rm group}: the blow-up relation {\bf (B)} produces two fixed points out of one. 

\

\subsection*{Determinant}

The determinant condition {\bf (Det)} of Section~\ref{sect:det} is also not preserved, in general. A fixed point for the action of an abelian subgroup $H\subseteq G$ does produce a symbol with weights
$$
\beta=(b_1,\ldots, b_n), \quad b_j\in H^\vee, 
$$
and it is the case that every equivariant blow-up will have a $H$-fixed points with  
weights
$$
\tilde{\beta}=(\tilde{b}_1,\ldots, \tilde{b}_n)
$$
with 
$$
b_1\wedge \cdots \wedge b_n = 
\tilde{b}_1\wedge \cdots \wedge \tilde{b}_n,
$$
but as elements in $\Burn_n(G)$, these may vanish, and thus be invisible.  

\subsection*{Cohomology}
\label{sect:coho}

The Burnside formalism fails to capture Amitsur invariants and cohomology. Indeed, by the observation in Section~\ref{sect:vanishing}, the nontrivial terms in the class of $X\times \bP^m$ vanish, for $m > |G|$. On the other hand, cohomological obstructions such as the nonvanishing of the Amitsur groups, or of $\rH^1(G,\Pic(X))$ are stable birational invariants. 



%%%%%%%%%%




%\subsection{Specialization sends birational types to birational types}
%\label{sec.specialization}
%We keep the notation of the previous section.
%We wish to verify that $\rho^G_t([X\actsfromright G])$ lies in the subgroup of $\Burn_n(G)$ spanned by the classes of projective $G$-varieties.

%POSSIBLY AN EXAMPLE HERE.
%Group of order $2$ acting freely on a curve, %specializing to
%linear combination of classes of curves with %group action,
%where there are fixed points.



%(Look at formula in Theorem 6.6 of \cite{BnG}, find a suitable
%replacement for the punctured normal bundle.
%For instance, $I=\{1,\dots,j\}$, we might consider a compactification
%of $\mathcal{N}_{D_1/\cX}|_{D_I}\oplus\dots\oplus \mathcal{N}_{D_{j-1}/\cX}|_{D_I}$, where an interesting possibility is the outcome of the
%iterated blow-up codimension $j-1$ then codimension $j-2$ and so on
%of the projectivization of
%\[
%\mathcal{N}_{D_1/\cX}|_{D_I}\oplus 
%(\mathcal{N}_{D_1/\cX}\otimes \mathcal{N}_{D_2/\cX})|_{D_I}\oplus
%\dots\oplus 
%(\mathcal{N}_{D_1/\cX}\otimes \dots\otimes\mathcal{N}_{D_{j-1}/\cX})|_{D_I}\oplus
%\cO_{D_I}
%\]
%which has interesting symmetries and might at the cost of
%combinatorial sophistication lead to an analysis that is similar to
%the proof of Lemma 5.7 of \cite{BnG}.) This is nonsense.

%By \cite[Rem.\ 6.3]{BnG}, the specialization of $[X\actsfromright G]$ remains unchanged after
%replacing $t$ by $t^{1/N}$, for any $N$. 
%Choosing a suitable $N$, we may assume that $\cX$ is a \emph{semistable} model, with special fiber $D_1\cup \dots\cup D_\ell$ (and $d_1=\dots=d_\ell=1$).



%Statement 1. There is an equality
%\begin{align*}
%&\sum_i [D_i\actsfromright G]\\
%&-\sum_{i<i'} [\bP(\cO_{\cX}(D_i)|_{D_{ii'}}\oplus 1)\actsfromright G]\\
%&+\sum_{i<i'<i''} [\bP(\cO_{\cX}(D_i)|_{D_{ii'i''}}\oplus \cO_{\cX}(D_i+D_{i'})|_{D_{ii'i''}}\oplus 1)\actsfromright G]\\
%&-\dots=
%\sum_i [D_i^\circ\actsfromright G]^{\mathrm{naive}}\\
%&-\sum_{i<i'} [\bP(\cO_{\cX}(D_i)|_{D_{ii'}^\circ}\oplus 1)^\circ\actsfromright G]^{\mathrm{naive}}+\dots.
%\end{align*}
%Notation: $D_{ii'}=D_i\cap D_{i'}$, etc.,
%punctured loci $D_i^\circ$ is the complement in $D_i$ of its intersection with any $D_{i'}$ ($i'\ne i$),
%$\bP()^\circ$ is punctured projectivized bundle over punctured locus, the locus of vanishing of each summand is removed and the union of deeper strata in the base is removed.

%Statement 2. The right-hand side of the equality in Statement 1
%is equal to the specialization of $[X\actsfromright G]$.

%Note of explanation.
%Geometrically, 
%$$
%\bP(\cO_{\cX}(D_i)|_{D_{ii'i''}^\circ}\oplus \cO_{\cX}(D_i+D_{i'})|_{D_{ii'i''}^\circ}\oplus 1)
%$$ 
%is sensitive to the order of $D_1$, $\dots$, $D_\ell$.
%But the class 
%$$
%[\bP(\cO_{\cX}(D_i)|_{D_{ii'i''}^\circ}\oplus \cO_{\cX}(D_i+D_{i'})|_{D_{ii'i''}^\circ}\oplus 1)\actsfromright G]
%$$ is invariant.
%Indeed, $\bP(\cO_{\cX}(D_i)|_{D_{ii'i''}^\circ}\oplus \cO_{\cX}(D_i+D_{i'})|_{D_{ii'i''}^\circ}\oplus 1)$ is an equivariant compactification of
%$\bP(\cO_{\cX}(D_i)|_{D_{ii'i''}^\circ}\oplus \cO_{\cX}(D_i+D_{i'})|_{D_{ii'i''}^\circ}\oplus 1)^\circ$.
%The latter is the $\G_m^2$-torsor associated with
%the restrictions to $D_{ii'i''}^\circ$ of
%$\cO_{\cX}(D_i)\oplus \cO_{\cX}(D_i+D_{i'})$.
%Generally, the $\G_m^n$-torsor associated with %$L_1\oplus\dots\oplus L_n$ may be identified with the
%$\G_m^n$-torsor associated with
%$L'_1\oplus\dots\oplus L'_n$, for
%$L'_i=L_1^{\otimes a_{i1}}\otimes\dots\otimes L_n^{\otimes a_{in}}$,
%$(a_{ij})\in \GL_n(\bZ)$.
%Combining with the canonical trivialization of
%$\cO_{\cX}(D_1+\dots+D_\ell)$,
%which when restricted to $D_{ii'i''}^\circ$ gives a
%trivialization of $\cO_{\cX}(D_i+D_{i'}+D_{i''})$,
%we see the invariance.
%(And not just for $i$, $i'$, $i''$, but for any number of indices.)


